\chapter{The Monster acts on the moonshine module: lifts, triality, the commutant and the converse}\label{ch:22}

\providecommand{\chtwentytwoM}{\mathbb{M}}
\providecommand{\chtwentytwoB}{\mathcal{B}}
\providecommand{\chtwentytwoA}{\mathcal{A}}
\providecommand{\chtwentytwoX}{\mathcal{X}}
\providecommand{\chtwentytwoW}{\mathcal{W}}
\providecommand{\chtwentytwoGvis}{G_{\mathrm{vis}}}
\providecommand{\chtwentytwoGmat}{G_{\mathrm{mat}}}
\providecommand{\chtwentytwoAutO}{\mathrm{AutO}}
\providecommand{\chtwentytwoCo}{\mathrm{Co}}
\providecommand{\chtwentytwoiota}{\iota_{\natural}}
\providecommand{\chtwentytwolean}[1]{\texttt{\directlua{ local s = "\luaescapestring{\detokenize{#1}}" for piece, sep in string.gmatch(s, "([^_.]*)([_.]?)") do if not (piece == "" and sep == "") then tex.sprint(-2, piece .. sep) if not (sep == "") then tex.sprint(string.char(92) .. "allowbreak{}") end end end}}}
\chaptermark{The Monster acts on the moonshine module}
\begingroup\setlength{\emergencystretch}{3em}\sloppy

This chapter documents the release of the library named in the front matter; every Lean quotation below is checked by
\texttt{scripts/textbook/sync.py} to occur verbatim (up to whitespace and comments) in the source files of that release.  Within a
\texttt{leancode} block, declarations appear in file order; proof bodies are omitted, and a comment line
\texttt{-- ⋯} marks a junction between two separately quoted declarations, where intervening declarations of the file may be
skipped (or the next declaration comes from another file named in the citation).  It is the
last chapter of the part ``Monstrous moonshine'', and it contains the part's headline.

Chapter~\ref{ch:20} proved the Conway--Norton theorem for every automorphism of $\Vnat$: the Lean theorem
\lean{conwayNorton : ConwayNortonConjecture} says that for every $g\in\Aut(\Vnat)$ the McKay--Thompson series
$T_g$ is the normalised Hauptmodul of a genus-zero group containing some $\Gamma_0(N)$.  It is a statement about
the group $\Aut(\Vnat)$, \emph{whatever that group is}.  Chapter~\ref{ch:21} imported, from the AxiomMath CFSG
library, an independent model of the Monster (``the Monster'' is this book's name for the Lean group
$\chtwentytwoAutO$; that it is a simple group of order $\approx8\cdot10^{53}$ is classical and is proved in AxiomMath's
library on Lean~4.35, but is \emph{not} a statement of LeanVOA): the rational Griess algebra $\chtwentytwoB$ (the Lean type
\lean{SpaceB}, a $\Q$-vector space of dimension $196884$), its automorphism group $\chtwentytwoAutO$, the visible
subgroup $\chtwentytwoGvis$ (classically isomorphic to $2^{1+24}.\chtwentytwoCo_1$; this isomorphism type is
\emph{not} formalised), the triality element $\tau_B$, and the generated group
$\chtwentytwoGmat=\langle\chtwentytwoGvis,\tau_B\rangle$; it also built an explicit isomorphism of algebras with forms
$\C\otimes_\Q\chtwentytwoB\cong\Vnat_2$ (\lean{griessIso}), and the language in which an element of
$\chtwentytwoAutO$ ``acts on $\Vnat$'' (\lean{IsVnatLiftable}, \lean{monsterLift}).

The present chapter connects the two groups.  It proves, in Lean:
\begin{enumerate}[label=(\Alph*)]
\item \textbf{Every element of $\chtwentytwoGmat$ acts on $\Vnat$} (\S\S\ref{ch22:sec:natural}--\ref{ch22:sec:tauB}):
  each element of the visible group and the triality element $\tau_B$ is the restriction to the Griess algebra of an
  automorphism of $\Vnat$.  Consequently Conway--Norton holds for every element of $\chtwentytwoGmat$, with
  no hypothesis beyond the membership $g\in\chtwentytwoGmat$ (\lean{conwayNorton_gmat}).
\item \textbf{The converse} (\S\S\ref{ch22:sec:traces}--\ref{ch22:sec:converse}): every automorphism of $\Vnat$ is
  rational on the Griess algebra, so $\Aut(\Vnat)$ is isomorphic to the subgroup of $\chtwentytwoAutO$ consisting of
  the elements that act on $\Vnat$ --- again with \emph{no hypotheses}
  (\lean{nonempty_autVnat_mulEquiv_liftable}).  The proof does not use Tits' theorem; it uses the moonshine theorem
  itself (rationality of all graded traces, \lean{rationalTraces}) and an absolute-irreducibility statement for
  $\chtwentytwoGmat$ (the \emph{commutant lemma}, \lean{gmatCommutantScalar}), proved by explicit Clifford theory and
  kernel-checked $M_{24}$ certificates.
\item \textbf{The headline} (\S\ref{ch22:sec:headline}): given one named hypothesis
  \lean{GmatEqAutO}, which is AxiomMath's statement $\chtwentytwoGmat=\chtwentytwoAutO$ after renaming (with one
  definition, \lean{SpinLattice.phi0}, restated in a definitionally equal form), there is a group isomorphism $\chtwentytwoAutO\cong\Aut(\Vnat)$ compatible with the Griess algebras, and monstrous
  moonshine holds for every element of the Monster (\lean{monstrousMoonshine}).
\end{enumerate}
The last statement is \textbf{conditional}: the hypothesis \lean{GmatEqAutO} is an explicit argument of the theorem.
Its \texttt{\#print axioms} output is clean, but, as Chapter~\ref{ch:01} (\S\ref{ch01:sec:uncond}) explains,
\texttt{\#print axioms} does not report explicit hypotheses.  This book never calls a theorem with the argument
\lean{(hM : GmatEqAutO)} unconditional.  Everything in (A) and (B) is unconditional.

\paragraph{The status of the hypothesis, up front.}  \lean{GmatEqAutO} is \emph{proved in AxiomMath's library}
(\lean{CFSG.Blueprint.NativeSporadic.MonsterModel.gmat_eq_griess}, Lean~4.35, commit \texttt{69cb972f}).  Its proof passes
through literature ``milestones'' whose names begin \lean{cl_lit_}; the ones on its path that were inspected
(\lean{CFSG.Milestones.cl_lit_stroth_leon_sims_baby_monster_recognition} and its \lean{_models} variant, and
\lean{CFSG.Milestones.cl_lit_goldschmidt_strongly_closed_abelian} from which they are derived) are \emph{theorems} in
that library, not axioms.  Its full axiom footprint has \emph{not} been re-checked: that needs \texttt{\#print axioms} on a
Lean~4.35 build of AxiomMath's library, which nobody in this project has run.  Our proposition agrees with AxiomMath's
statement after renaming, with \lean{phi0} restated definitionally equal (source-text audit, not a cross-toolchain
proof; \S\ref{ch22:sec:headline}).

As everywhere in this book, our concern is the \emph{meaning} of the Lean statements.  Proofs are explained
mathematically; we do not describe how the Lean proof scripts are organised beyond, at most, naming the file.

\section{Overview: the Monster as the symmetry group of the moonshine module}\label{ch22:sec:overview}

\subsection{The classical picture}

Frenkel, Lepowsky and Meurman constructed $\Vnat$ in 1988 as a $\Z/2$-orbifold of the Leech lattice vertex operator
algebra,
\begin{equation*}
  \Vnat \;=\; V_\Leech^+\;\oplus\;(V_\Leech^T)^+,
\end{equation*}
(Chapters~\ref{ch:02}--\ref{ch:14}) and proved that the Monster acts on it.  Their argument has two halves.
\begin{itemize}
\item \emph{The centralizer of an involution acts.}  The operator $z$ which is $+1$ on $V_\Leech^+$ and $-1$ on
  $(V_\Leech^T)^+$ is an automorphism of $\Vnat$ (it is the ``quantum'' or dual $\Z/2$ symmetry of the orbifold).  The
  group of automorphisms of $\Vnat$ that are visible in the orbifold description --- lattice automorphisms of
  $\Leech$, lifted, and the extraspecial group $2^{1+24}$ acting on the twisted sector --- has the shape
  $C=2^{1+24}.\chtwentytwoCo_1$, and it is the centralizer of $z$.
\item \emph{Triality.}  There is an extra automorphism $\sigma$ of $\Vnat$, not in $C$, constructed from a second
  description of $\Vnat$ as an orbifold (Chapters~\ref{ch:11}--\ref{ch:14}).  The group generated by $C$ and $\sigma$
  acts on the weight-two space $\Vnat_2$ (the Griess algebra) by algebra automorphisms; by Griess' and Conway's work
  it is the Monster $\chtwentytwoM$.  The converse, that $\Aut(\Vnat)$ is no larger, follows from Tits' theorem that
  $\chtwentytwoM$ is the full automorphism group of the Griess algebra, together with the fact that $\Vnat$ is
  generated by $\Vnat_2$.
\end{itemize}

\subsection{The route taken in Lean}

The Lean development reproduces the first half literally, but in AxiomMath's coordinates, and replaces the
second half by a new argument.  The chain is the following; each arrow is a Lean theorem.
\begin{enumerate}[label=(\arabic*)]
\item \textbf{Lifts of $C$.}  A pair (automorphism $h$ of $V_\Leech^+$, compatible isometry $\varphi$ of
  $(V_\Leech^T)^+$) is an automorphism of $\Vnat$ (\lean{vnatAutOfPair}).  Every isometry $\iota$ of $\Leech$ gives
  such a pair (\lean{liftTwVnatAut}); so does every $\mu\in\Leech$ (\lean{extraVnatAut}).
\item \textbf{Block forms and the rational criterion.}  On $\Vnat_2\cong\C\otimes\chtwentytwoB$ these lifts act by
  explicit block matrices; the lift of $\iota$ is \emph{liftable} (comes from $\chtwentytwoAutO$) exactly when one
  $4096\times4096$ ``spin'' block is rational, and that reduces to a rational intertwining equation.
\item \textbf{The visible group, uniformly.}  For every compatible pair $(s,g)$ defining an element of
  $\chtwentytwoGvis$, the matrix $s^{\mathsf T}$ solves the equation up to an extraspecial correction; hence every element
  of $\chtwentytwoGvis$ is liftable (\lean{liftVisRAllVisible}).  This proof does \emph{not} go through generators.
  Separately, the equation is solved explicitly for the Golay sign changes, the code-preserving coordinate permutations
  ($2^{12}{:}M_{24}$) and Conway's non-monomial element $\xi$; these explicit lifts, with their explicit block forms, are
  what the commutant lemma (6) needs, because there one must know \emph{how} specific elements of $\chtwentytwoGvis$ act
  block by block, not merely that they act.
\item \textbf{Triality.}  $\sigma$ is rational on $\Vnat_2$, so it is the lift of an element $\sigma_B\in\chtwentytwoAutO$
  (\lean{sigmaBOdd_liftable}); conjugation by $\sigma_B$ carries the parity involution $x_B$ to $y_B$ or $z_B$; the
  parity-centralizer theorem $\chtwentytwoAutO\cap C(x_B)=\chtwentytwoGvis$ of Chapter~\ref{ch:21} then shows that
  $\tau_B$ is liftable (\lean{isVnatLiftable_tauB}).  Hence every element of $\chtwentytwoGmat$ is liftable.
\item \textbf{Rational traces.}  Galois invariance $T_{g^u}=T_g$ (a consequence of the moonshine theorem) and a
  Ramanujan-sum argument show that every graded trace of every automorphism of $\Vnat$ is rational
  (\lean{rationalTraces}).
\item \textbf{The commutant lemma.}  Every $\Q$-linear endomorphism of $\chtwentytwoB$ commuting with
  $\chtwentytwoGmat$ is $a\cdot1+b\cdot P_0$, $P_0$ the projection onto the line of the identity element
  (\lean{gmatCommutantScalar}).
\item \textbf{The converse.}  (5), (6) and Schur's orthogonality relations show that every automorphism of $\Vnat$
  maps the rational lattice $\mathrm{griessNat}(\chtwentytwoB)\subset\Vnat_2$ into itself (\lean{autVnat_rational});
  therefore $\Aut(\Vnat)\cong$ the liftable subgroup of $\chtwentytwoAutO$ (\lean{nonempty_autVnat_mulEquiv_liftable}).
\item \textbf{The headline.}  If $\chtwentytwoGmat=\chtwentytwoAutO$ (\lean{GmatEqAutO}), the liftable subgroup is all of
  $\chtwentytwoAutO$, and (3), (4), (7) give $\chtwentytwoAutO\cong\Aut(\Vnat)$ and monstrous moonshine for the Monster
  (\lean{monsterEquivAutVnat}, \lean{monstrousMoonshine}).
\end{enumerate}

\begin{physics}
\textbf{The Monster as the full symmetry group of the moonshine CFT.}  $\Vnat$ is the chiral algebra of a
holomorphic two-dimensional conformal field theory with central charge $24$; it has a single irreducible module
(itself), so its torus partition function is the modular function $J(q)=q^{-1}+196884q+\cdots$.  Its weight-one
space vanishes, so the theory has no conserved currents and hence no continuous symmetries: every global symmetry
is a finite-order automorphism of the chiral algebra, and the group of all of them is $\Aut(\Vnat)$.  The
statement that this group is the Monster is the statement that the largest sporadic simple group is \emph{exactly}
the global symmetry group of this CFT --- not merely a group of symmetries.  For a symmetry $g$, the twisted
partition function $T_g(q)=\tr_{\Vnat}(g\,q^{L_0-1})$ (the torus amplitude with a $g$-twist along the time
circle) is the McKay--Thompson series, and moonshine says that each of these is the partition function of a
genus-zero modular curve.  The present chapter is the part of the story that identifies the symmetry group.
\end{physics}

\begin{physics}
\textbf{The $\Z/2$-orbifold and its centralizer $2^{1+24}.\chtwentytwoCo_1$.}  The Leech lattice CFT $V_\Leech$
(24 chiral bosons on the torus $\mathbb{R}^{24}/\Leech$) has the reflection symmetry $X\mapsto-X$.  Gauging it (the orbifold)
keeps the invariant states $V_\Leech^+$ and adds the invariant part of the twisted sector, in which the bosons are
antiperiodic and the ground states form the $2^{12}$-dimensional representation $T$ of a Clifford-like extension
of $\Leech/2\Leech$.  The orbifold theory has a new $\Z/2$ symmetry $z$ (``quantum symmetry''), $+1$ on the
untwisted, $-1$ on the twisted sector.  Symmetries of the parent theory that commute with the reflection descend:
the automorphisms of the lattice ($\chtwentytwoCo_0=2.\chtwentytwoCo_1$, acting on the bosons by rotations) and
the momentum phases $e^{i\pi\langle\mu,p\rangle}$, $\mu\in\Leech$, which act on the twisted ground states through the
extraspecial group.  Together they form $2^{1+24}.\chtwentytwoCo_1$, the centralizer of $z$ in the Monster.  These
are the elements lifted in \S\S\ref{ch22:sec:natural}--\ref{ch22:sec:visible}.
\end{physics}

\subsection{What is recalled from Chapter~\ref{ch:21}}\label{ch22:sec:recall}

The following objects are defined and fidelity-checked in Chapter~\ref{ch:21}; we recall each in the form our
statements need, and quote the definitions a referee must see to read the statements of this chapter.

\paragraph{The rational Griess algebra.}  $\chtwentytwoB=$ \lean{SpaceB} $=\chtwentytwoA\times\chtwentytwoX\times
\chtwentytwoW$ with $\chtwentytwoA=$ \lean{symMat} the symmetric $24\times24$ rational matrices (dimension $300$),
$\chtwentytwoX=$ \lean{SpaceX} $=\Q^{\mathrm{ShortIdx}}$ with one coordinate for each pair $\pm\alpha$ of minimal
vectors of the Leech lattice (there are $196560/2=98280$ such pairs; $\mathrm{ShortIdx}$ is AxiomMath's index set for
them, Chapter~\ref{ch:21}), and $\chtwentytwoW=$ \lean{SpaceW}
the $4096\times24$ rational matrices; $300+98280+98304=196884$.  It carries AxiomMath's commutative product
\lean{gmul} and symmetric bilinear form \lean{bform}.  An element is written $u=(A,x,N)$; in Lean the components
are \lean{u.1}, \lean{u.2.1}, \lean{u.2.2}.

\paragraph{Matrix conventions.}  AxiomMath's matrices act on \emph{row} vectors: for a rational matrix $g$ of
$\chtwentytwoCo_0$, AxiomMath's visible element acts on $\chtwentytwoA$ by $A\mapsto g^{\mathsf T}Ag$.  The library's own
matrix of a Leech isometry $\iota$ (\lean{liftBlockMatrix ι}, written $g_\iota$ below) acts on \emph{column} vectors, and
the corresponding lift acts on $\chtwentytwoA$ by $A\mapsto g_\iota Ag_\iota^{\mathsf T}$.  For the same isometry the two
are transposes of each other, $g_\iota=g^{\mathsf T}$ (\lean{liftVisBlockMatrix}, \S\ref{ch22:sec:visible}), so the two
formulas agree.

\paragraph{The dictionary in symbols.}  Chapter~\ref{ch:21} builds the rational dictionary
$\mathrm{griessNat}:\chtwentytwoB\to\Vnat_2$, $\mathrm{griessNat}(A,x,N)=(v_{\mathrm{untw}},v_{\mathrm{tw}})\in
V_\Leech^+\oplus T^+$, with
\begin{equation*}
  v_{\mathrm{untw}}=q_C(A^\pi)+\sum_{c}x_c\,\zeta'_c\,x_{\alpha_c},\qquad
  v_{\mathrm{tw}}=\lambda\sum_{i=0}^{23}h_{f_i}(-\tfrac12)\otimes\Psi\bigl(N_{\cdot,\pi(i)}\bigr),
\end{equation*}
where: Conway's coordinates realise $\Leech$ inside $\Z^{24}$ as the vectors $u$ with $u_i\equiv m+2d_i\pmod4$ for a
common $m$ and a Golay codeword $d$, and $\sum_iu_i\equiv4m\pmod8$, with inner product $u\cdot v/8$ (so minimal vectors
have norm $4$); $f_i$ is the $i$-th Conway axis vector $8e_i$, used as an oscillator direction; $\pi$ is the fixed
permutation of $\{0,\dots,23\}$ matching AxiomMath's numbering with the library's oscillator frame, and $A^\pi$ the
correspondingly permuted matrix; $q_C(M)$ is the weight-two quadratic oscillator state
$\sum_{i,j}M_{ij}h_{f_i}(-1)h_{f_j}(-1)\mathbb 1$ (with the normalisation of Chapter~\ref{ch:21}); $\alpha_c$ is a chosen
minimal vector in the pair $c$ and $x_\alpha=e^\alpha+e^{-\alpha}$; $\zeta'_c$ (\lean{fullShortSign}) is a fixed complex
phase chosen so that the dictionary is compatible with AxiomMath's product; $\Psi$ (\lean{spinIso}) identifies AxiomMath's
spin module $\Q^{4096}$ (tensored with $\C$) with the $4096$-dimensional module $T$ of the twisted sector; $\lambda$
(\lean{spinDictScale}) is a normalising scalar.  The twisted sector starts at weight $\tfrac32$ (the ground states of
$24$ antiperiodic bosons, $24/16$), so its weight-two part is $h(-\tfrac12)\otimes T$, of dimension $24\cdot4096=98304$:
this is why $\chtwentytwoW$ consists of $4096\times24$ matrices.  \lean{griessIso} is the $\C$-linear extension of
\lean{griessNat}; we write $\mathcal G$ for it.

\paragraph{The Monster.}  $\chtwentytwoAutO$ is the subgroup of $\mathrm{GL}_\Q(\chtwentytwoB)$ of maps preserving
the product and the form:
\begin{leancode}
def IsAlgAut (f : SpaceB ≃ₗ[ℚ] SpaceB) : Prop := ∀ u v, f (gmul u v) = gmul (f u) (f v)
\end{leancode}
\begin{leancode}
def AutO : Subgroup (SpaceB ≃ₗ[ℚ] SpaceB) where
  carrier := {g | IsAlgAut g ∧ ∀ u v, bform (g u) (g v) = bform u v}
\end{leancode}
(\leanpath{VOA/Moonshine/MonsterGriess/AutO.lean}).  The visible group is the image of AxiomMath's map
\lean{thetaEquiv} on compatible pairs $(s,g)$, $s$ a rational $4096\times4096$ matrix normalising the extraspecial
group and $g$ a rational matrix of $\chtwentytwoCo_0$; $\mathrm{thetaEquiv}(s,g)$ acts by
$(A,x,N)\mapsto(g^{\mathsf T}Ag,\ \text{a signed permutation of }x,\ s^{\mathsf T}Ng)$ (Chapter~\ref{ch:21}):
\begin{leancode}
def Gvis : Subgroup (SpaceB ≃ₗ[ℚ] SpaceB) where
  carrier := Set.range thetaEquiv
\end{leancode}
(\leanpath{VOA/Moonshine/MonsterGriess/VisiblePair.lean}).  The parity involution and its two triality conjugates,
and the generated group:
\begin{leancode}
noncomputable def xB : SpaceB ≃ₗ[ℚ] SpaceB :=
  LinearEquiv.ofInvolutive parityMap fun u => by simp [parityMap_apply]
\end{leancode}
(\leanpath{VOA/Moonshine/MonsterGriess/GriessTrilinear.lean}; \lean{parityMap} is $(A,x,N)\mapsto(A,x,-N)$),
\begin{leancode}
noncomputable def yB : SpaceB ≃ₗ[ℚ] SpaceB := tauB * xB * tauB⁻¹
-- ⋯
noncomputable def zB : SpaceB ≃ₗ[ℚ] SpaceB := tauB ^ 2 * xB * (tauB ^ 2)⁻¹
-- ⋯
noncomputable def Gmat : Subgroup (SpaceB ≃ₗ[ℚ] SpaceB) :=
  Subgroup.closure ((Gvis : Set (SpaceB ≃ₗ[ℚ] SpaceB)) ∪ {tauB})
\end{leancode}
(\leanpath{VOA/Moonshine/MonsterGriess/TrialityFourGroup.lean}).  Here \lean{tauB} is AxiomMath's explicit rational
matrix of order $3$ (Chapter~\ref{ch:21}; \lean{tauB_pow_three}, \lean{trace_tauB : ... = 783}); it fixes every
diagonal matrix in $\chtwentytwoA$ and permutes cyclically three ``triality basis'' vectors in each of the pair, odd and
octad blocks (\lean{tauB_diag}, \lean{tauB_pair}, \lean{tauB_odd}, \lean{tauB_oct}).  Chapter~\ref{ch:21} also proves
$\tau_B\in\chtwentytwoAutO$ (\lean{tauB_mem_AutO}), $\chtwentytwoGmat\le\chtwentytwoAutO$ (\lean{Generated.Gmat_le_AutO}) and the
\emph{parity-centralizer theorem}
\begin{equation}\label{ch22:eq:parity}
  \chtwentytwoAutO\cap C(x_B)\;=\;\chtwentytwoGvis
  \qquad(\text{\lean{Parity.AutO_inf_centralizer_xB_eq_Gvis}}),
\end{equation}
ported from AxiomMath; it is used twice below.

\paragraph{The Griess isomorphism and the lift.}  \lean{griessNat : SpaceB →+ VnatVOA.grading 2} is the explicit
rational dictionary, and \lean{griessIso : ℂ ⊗[ℚ] SpaceB ≃ₗ[ℂ] VnatVOA.grading 2} its $\C$-linear extension
($1\otimes u\mapsto\mathrm{griessNat}(u)$); it carries \lean{gmul} to the Griess product of $\Vnat$ and
$\mathrm{natForm}=2\cdot$\lean{bform} (\lean{griessIso_mul}, \lean{griessIso_form}).  For
$g\in\mathrm{GL}_\Q(\chtwentytwoB)$ the \emph{Griess conjugate} is the transported map on $\Vnat_2$,
\begin{leancode}
noncomputable def griessConj (g : SpaceB ≃ₗ[ℚ] SpaceB) :
    VnatVOA.grading 2 ≃ₗ[ℂ] VnatVOA.grading 2 :=
  griessIso.symm ≪≫ₗ griessBaseChange g ≪≫ₗ griessIso
\end{leancode}
and \lean{vnatWeightTwoRestrict} restricts an automorphism of $\Vnat$ to $\Vnat_2$:
\begin{leancode}
noncomputable def vnatWeightTwoRestrict : AutVnat →* (VnatVOA.grading 2 ≃ₗ[ℂ] VnatVOA.grading 2) :=
  GriessAutVnat.subtype.comp vnatGriessRestrict
\end{leancode}
(\leanpath{VOA/Moonshine/MoonshineCNMonsterActionConj.lean}).  An element $g\in\chtwentytwoAutO$ \emph{acts on}
$\Vnat$, or \emph{is liftable}, if some automorphism of $\Vnat$ restricts to its Griess conjugate:
\begin{leancode}
def IsVnatLiftable (g : AutO) : Prop :=
  ∃ h : AutVnat, vnatWeightTwoRestrict h = griessConj (g : SpaceB ≃ₗ[ℚ] SpaceB)
-- ⋯
noncomputable def liftableSubgroup : Subgroup AutO :=
  vnatGriessRestrict.range.comap griessConjAut
-- ⋯
noncomputable def monsterLift : liftableSubgroup →* AutVnat :=
  (MonoidHom.ofInjective vnatGriessRestrict_injective).symm.toMonoidHom.comp
    (griessConjAut.subgroupComap vnatGriessRestrict.range)
-- ⋯
noncomputable def monsterLiftOfAll (hlift : ∀ g : AutO, IsVnatLiftable g) : AutO →* AutVnat :=
  monsterLift.comp (MonoidHom.codRestrict (MonoidHom.id AutO) liftableSubgroup
    fun g => mem_liftableSubgroup.2 (hlift g))
\end{leancode}
(\leanpath{VOA/Moonshine/MoonshineCNMonsterActionLift.lean}).  The liftable elements form a subgroup
(\lean{mem_liftableSubgroup}); the lift of a liftable $g$ is \emph{unique}, because an automorphism of $\Vnat$ is
determined by its restriction to $\Vnat_2$ (\lean{vnatWeightTwoRestrict_injective}: an automorphism fixing
$\Vnat_2$ pointwise is the identity, Chapter~\ref{ch:20}); \lean{monsterLift} is the resulting injective homomorphism, with
$\mathrm{monsterLift}(g)(\mathrm{griessNat}\,u)=\mathrm{griessNat}(g u)$ (\lean{monsterLift_apply_griessNat}).  Two
criteria from Chapter~\ref{ch:21} are used repeatedly:
\begin{leancode}
theorem isVnatLiftable_of_rational (h : AutVnat)
    (hrat : ∀ u : SpaceB, ∃ u' : SpaceB, h (griessNat u : Vnat) = griessNat u') :
    ∃ g : liftableSubgroup, monsterLift g = h
\end{leancode}
(an automorphism of $\Vnat$ that maps the rational lattice $\mathrm{griessNat}(\chtwentytwoB)$ into itself \emph{is} the
lift of an element of $\chtwentytwoAutO$; \leanpath{VOA/Moonshine/MoonshineCNMonsterActionLift.lean}), and
\lean{conwayNorton_autO_of_generators} (\leanpath{VOA/Moonshine/MoonshineCNMonsterActionMoonshine.lean}): if a
generating set lifts, every element lifts, and then the moonshine theorem applies to the lifts.

\paragraph{Why liftability is a rationality question.}  This is the idea on which the whole chapter rests.  $\Vnat_2$ is a
complex vector space carrying a rational structure, the $\Q$-span $\mathrm{griessNat}(\chtwentytwoB)$ of the dictionary
vectors.  $\chtwentytwoAutO$ consists, by definition, of the $\Q$-linear automorphisms of that rational structure which
preserve the product and the form.  An automorphism $h$ of $\Vnat$ is only $\C$-linear, but it always preserves the
Griess product and the form on $\Vnat_2$.  Hence:
\begin{quote}
\emph{$h$ comes from an element of $\chtwentytwoAutO$ if and only if $h|_{\Vnat_2}$ maps the rational lattice
$\mathrm{griessNat}(\chtwentytwoB)$ into itself,}
\end{quote}
and then the element is unique, since $h$ is determined by $h|_{\Vnat_2}$.  This is \lean{isVnatLiftable_of_rational} with
\lean{mem_range_monsterLift_iff} (Chapter~\ref{ch:21}).  Every later section is an application: to show that a given
element of $\chtwentytwoAutO$ acts on $\Vnat$, we construct \emph{some} automorphism of $\Vnat$ that acts on $\Vnat_2$ by a
rational matrix equal to it (\S\S\ref{ch22:sec:natural}--\ref{ch22:sec:tauB}); to show that every automorphism of $\Vnat$ comes
from $\chtwentytwoAutO$, we prove that its matrix on $\Vnat_2$ is rational
(\S\S\ref{ch22:sec:traces}--\ref{ch22:sec:converse}).

\paragraph{The moonshine statement.}  From Chapter~\ref{ch:16}: \lean{AutVnat} $=\Aut(\Vnat)$ is the group of
vertex-operator-algebra automorphisms of \lean{VnatVOA}; \lean{mckayThompsonFun g} is
$T_g(\tau)=q^{-1}+\sum_{n\ge1}\tr(g|\Vnat_{n+1})q^n$ on the upper half-plane;
\begin{leancode}
def ConwayNortonFor (g : AutVnat) : Prop :=
  ∃ (Γ : Subgroup (GL (Fin 2) ℝ)) (N : ℕ), IsMoonshineGroup Γ N ∧
    IsNormalizedHauptmodul Γ (mckayThompsonFun g)
-- ⋯
def ConwayNortonVia {𝕄 : Type*} [Group 𝕄] (ι : 𝕄 →* AutVnat) : Prop :=
  ∀ g : 𝕄, ConwayNortonFor (ι g)
-- ⋯
def RationalTraces : Prop :=
  ∀ (g : AutVnat) (n : ℕ), ∃ r : ℚ, PowerSeries.coeff n (mckayThompson g) = r
\end{leancode}
(\leanpath{VOA/Moonshine/MoonshineCNStatement.lean}), where \lean{IsMoonshineGroup Γ N} says $N\ge1$,
$\Gamma\le\mathrm{SL}_2(\mathbb{R})$, $\Gamma_0(N)\le\Gamma$ and $\Gamma$ is commensurable with $\mathrm{SL}_2(\Z)$, and
\lean{IsNormalizedHauptmodul Γ f} is the eight-clause definition of a normalised Hauptmodul of
Chapter~\ref{ch:16}.  The theorem \lean{ConwayNorton.conwayNorton : ConwayNortonConjecture} (Chapter~\ref{ch:20})
has no hypotheses.

\begin{physics}
The dictionary $\C\otimes\chtwentytwoB\cong\Vnat_2$ says that the $196884$ weight-two states of the moonshine CFT
--- the stress tensor $\omega$ and $196883$ further primary fields of dimension $2$ --- close under the
$(1)$-product of their operator product expansion into AxiomMath's rational algebra.  ``$g$ is liftable'' means
that the symmetry $g$ of this finite set of OPE coefficients extends to a symmetry of the \emph{whole} chiral
algebra.  Since the weight-two fields generate the theory, such an extension is unique when it exists.
\end{physics}

\subsection{Notation used in this chapter}\label{ch22:sec:notation}

Several letters are used by the sources for different objects.  This chapter fixes the following meanings.
\begin{center}\small
\begin{longtable}{p{0.2\textwidth}p{0.74\textwidth}}
\toprule
symbol & meaning\\
\midrule
\endhead
$\chtwentytwoB=\chtwentytwoA\oplus\chtwentytwoX\oplus\chtwentytwoW$ & AxiomMath's rational Griess algebra and its three blocks
  (symmetric matrices, short coordinates, $4096\times24$ spin matrices)\\
$\mathcal D,\ \mathcal D_0,\ \mathcal O$ & diagonal, traceless diagonal, off-diagonal symmetric matrices in $\chtwentytwoA$\\
$\mathbf{e}_1$; $P_0$ & the vector $(I,0,0)\in\chtwentytwoB$; the $b$-orthogonal projection onto $\Q\mathbf{e}_1$\\
$b$ & AxiomMath's form \lean{bform}\\
$\mathcal G$ & \lean{griessIso}$:\C\otimes\chtwentytwoB\to\Vnat_2$; $\mathrm{griessNat}(u)=\mathcal G(1\otimes u)$\\
$T$, $T^+$ & the $4096$-dimensional extraspecial module; the twisted sector $(V_\Leech^T)^+$\\
$T_g$ & the McKay--Thompson function of $g\in\Aut(\Vnat)$\\
$c$ & an element of \lean{ShortIdx} (a pair $\pm\alpha_c$ of minimal vectors), \emph{except} in $c(\alpha)$\\
$c(\alpha)$ & the coboundary \lean{latCob} of a lattice isometry\\
$\iota$; $O_\iota$; $g_\iota$ & a Leech isometry; its orthogonal $24\times24$ complex matrix on the bosons; its rational matrix
  in AxiomMath's coordinates\\
$\rho$; $\lambda_\iota$ & a charge intertwiner on $T$; the normalising scalar of the lift of $\iota$\\
$d$ & a Golay codeword (in sign changes), or a scalar in the commutant proof (stated there)\\
$\pi$ & a permutation of the $24$ coordinates (code-preserving in \S\ref{ch22:sec:generators}); also the dictionary
  permutation in $A^\pi$\\
$\pi_A,\pi_B,\pi_C$ & the three explicit code-preserving permutations of the commutant proof\\
$Q$; $Q_\alpha$ & the extraspecial group $2^{1+24}$ acting on $\chtwentytwoB$; the rational unphased spin matrix of the class of
  $\alpha\in\Leech$ \\
$q_C$ & the quadratic oscillator state of the dictionary\\
$Z(\alpha)$ & the sign refinement (spin phase) of $\alpha$; not to be confused with the involution $z$\\
$\sigma$; $\sigma_0$; $\sigma_B$ & FLM's triality automorphism of $\Vnat$; its restriction to $V_0^+$; the element of
  $\chtwentytwoAutO$ it induces\\
$\Phi$, $\Theta$; $\mathfrak D$ & the twisted halves of triality; the triality data (Lean \lean{SigmaData})\\
$z$, $z'$ & the orbifold involution ($+1$ on $V_\Leech^+$, $-1$ on $T^+$); the frame-sector sign\\
$x_B,y_B,z_B$ & AxiomMath's parity involution and its two $\tau_B$-conjugates; through the dictionary $x_B\leftrightarrow z$,
  $\{y_B,z_B\}\leftrightarrow\{z',zz'\}$ (\S\ref{ch22:sec:tauB})\\
$U_{jk}$ & the off-diagonal symmetric matrix unit (\lean{sigmaBlockRH j k}, \lean{basisA s(j,k)})\\
$\mathcal K$; $\mathcal R$ & an endomorphism commuting with $\chtwentytwoGmat$ (commutant lemma); its residual $\mathcal K-d\cdot1$\\
$H$ & the complex matrix of $h|_{\Vnat_2}$ (\S\ref{ch22:sec:rationalgriess})\\
$\mathrm{gmatLift}$ & the injective homomorphism $\chtwentytwoGmat\to\Aut(\Vnat)$, $g\mapsto\mathrm{monsterLift}(g)$
  (\lean{MonsterAction.gmatLift}, Theorem~\ref{ch22:thm:omega})\\
\bottomrule
\end{longtable}
\end{center}

\paragraph{Lean combinators used in the quotations.}  \lean{e₁ ≪≫ₗ e₂} is the composite of linear equivalences, first
\lean{e₁}; \lean{MonoidHom.ofInjective} turns an injective homomorphism into an isomorphism onto its range;
\lean{subgroupComap} restricts a homomorphism to the preimage of a subgroup; \lean{MonoidHom.codRestrict} restricts the
codomain to a subgroup containing the range; \lean{(h).choose} is the witness picked by \lean{Classical.choice} from an
existence proof \lean{h}.  \lean{LinearMap.smulRight φ v} is $u\mapsto\varphi(u)v$.

%======================================================================
\section{Lifting \texorpdfstring{$C=2^{1+24}.\mathrm{Co}_1$}{C = 2\^{}(1+24).Co1} to automorphisms of \texorpdfstring{$V^\natural$}{V-natural}}\label{ch22:sec:natural}
%======================================================================

All declarations of \S\S\ref{ch22:sec:natural}--\ref{ch22:sec:tauB} whose names begin with \lean{lift}, \lean{sigmaB},
\lean{tauBLift}, \lean{extra} live in the namespace \chtwentytwolean{VOA.Moonshine.ConwayNorton.M8} (the suffix \lean{M8} is a historical node label).  The files are
\leanpath{VOA/Moonshine/MoonshineLRLift*.lean}.

\subsection{Naturality: compatible pairs are automorphisms}

Recall from Chapters~\ref{ch:09}--\ref{ch:10} that the vertex operator of $\Vnat=V_\Leech^+\oplus T^+$ (we write
$T^+$ for $(V_\Leech^T)^+$, the Lean submodule \lean{TP}, and \lean{SkewVP} for $V_\Leech^+$) is the block operator
\begin{equation}\label{ch22:eq:block}
  Y^\natural\bigl((u,t),z\bigr)=\begin{pmatrix}Y(u,z)&\overline W(t,z)\\ W(t,z)&Y^T(u,z)\end{pmatrix},
  \qquad W(t,z)u=e^{zL_{-1}}Y^T(u,-z)t,
\end{equation}
where $Y^T$ is the module action of $V_\Leech^+$ on $T^+$ (\lean{plusModule.YM}) and $\overline W$ is the adjoint of
$W$ with respect to the invariant forms $B_V$ (\lean{leechBV}) on $V_\Leech^+$ and \lean{TPForm} on $T^+$.

\begin{theorem}[Naturality of $\Vnat$]\label{ch22:thm:natural}
Let $h$ be an automorphism of the vertex operator algebra $V_\Leech^+$ and $\varphi:T^+\to T^+$ a $\C$-linear
bijection such that
\begin{enumerate}[label=(\roman*)]
\item $\varphi(u_{(n)}t)=(hu)_{(n)}\varphi(t)$ for all $u\in V_\Leech^+$, $n\in\Z$, $t\in T^+$, and
\item $\mathrm{TPForm}(\varphi t,\varphi t')=\mathrm{TPForm}(t,t')$ for all $t,t'$.
\end{enumerate}
Then $h\oplus\varphi:(u,t)\mapsto(hu,\varphi t)$ is an automorphism of the vertex operator algebra $\Vnat$.
\end{theorem}
\begin{leancode}
noncomputable def vnatAutOfPair (h : Aut leechVOAPlus) (φ : TP ≃ₗ[ℂ] TP)
    (hφ : ∀ (u : SkewVP) (n : ℤ) (t : TP),
      φ (mode (plusModule.YM u) n t) = mode (plusModule.YM (h u)) n (φ t))
    (hT : ∀ t t' : TP, TPForm (φ t) (φ t') = TPForm t t') : AutVnat where
-- ⋯
@[simp] theorem vnatAutOfPair_apply (u : SkewVP) (t : TP) :
    vnatAutOfPair h φ hφ hT ((u, t) : Vnat) = (h u, φ t) := rfl
\end{leancode}
\begin{leancode}
theorem lrLift_leechBV (h : Aut leechVOAPlus) (u v : SkewVP) :
    leechBV (h u) (h v) = leechBV u v
\end{leancode}
(\leanpath{VOA/Moonshine/MoonshineLRLiftNatural.lean}; full names \chtwentytwolean{VOA.Moonshine.ConwayNorton.M8.vnatAutOfPair},
\lean{...M8.vnatAutOfPair_apply}, \lean{...M8.lrLift_leechBV}.)
\begin{dictionary}
\textbf{Decoding.}  \lean{Aut leechVOAPlus} is the type of vertex-operator-algebra automorphisms of $V_\Leech^+$
(bijective linear maps commuting with all modes and fixing the vacuum and the conformal vector; Chapter~\ref{ch:01}).
\lean{TP ≃ₗ[ℂ] TP} is a $\C$-linear bijection of $T^+$.  \lean{mode (plusModule.YM u) n t} is $u_{(n)}t$ (the
convention \lean{mode f n = f (-n-1)} of Chapter~\ref{ch:01}).  The output type \lean{AutVnat} is $\Aut(\Vnat)$; the
\lean{where} block (not reproduced) builds the automorphism, and \lean{vnatAutOfPair_apply} says, by \lean{rfl},
that it is the block-diagonal map $(u,t)\mapsto(hu,\varphi t)$ on \lean{Vnat = SkewVP × TP}.

\textbf{Hypotheses.}  Exactly (i) (\lean{hφ}) and (ii) (\lean{hT}).  Two conditions one might expect are \emph{not}
hypotheses: that $h$ preserves $B_V$, and that $\varphi$ commutes with $L_{-1}$.  The first is the separate lemma
\lean{lrLift_leechBV}, valid for \emph{every} automorphism of $V_\Leech^+$; the second follows from (i) at
$u=\omega$ because $h\omega=\omega$.

\textbf{Same statement.}  Yes; the theorem is a construction, and the \lean{_apply} lemma pins the underlying map.
\textbf{Could it be vacuous?}  No: $h=1$, $\varphi=1$ satisfies the hypotheses, and the extraspecial and lattice lifts
below provide many nontrivial instances.
\end{dictionary}
\begin{proof}[Proof on paper]
We must check $g\,Y^\natural(a,z)\,g^{-1}=Y^\natural(ga,z)$ for $g=h\oplus\varphi$, block by block in
\eqref{ch22:eq:block}.  The top-left block is $hY(u,z)h^{-1}=Y(hu,z)$ because $h$ is an automorphism.  The
bottom-right block is hypothesis (i).  Applying (i) to the conformal vector $u=\omega$ (fixed by $h$) gives
$\varphi L_n=L_n\varphi$ for all $n$, in particular for $L_{-1}$ and $L_0$, so $\varphi$ preserves conformal weights.
For the bottom-left block, the skew-symmetry definition of $W$ gives
\begin{equation*}
  \varphi\,W(t,z)u=\varphi\,e^{zL_{-1}}Y^T(u,-z)t=e^{zL_{-1}}Y^T(hu,-z)\varphi t=W(\varphi t,z)\,hu.
\end{equation*}
For the top-right block, $\overline W$ is defined in Chapter~\ref{ch:09} as the adjoint of $W$ for the two
invariant forms, $B_V(\overline W(t,z)a,u)=\mathrm{TPForm}(a,W(t^*,z^{-1})u)$, where $t\mapsto t^*$ is the
contragredient (BPZ-type) operation of Chapter~\ref{ch:09}, built from $L_0$, $L_1$ and fixed scalars; $\varphi$ commutes
with all $L_n$, hence with $t\mapsto t^*$.  Since $h$ preserves $B_V$ and $\varphi$ preserves
$\mathrm{TPForm}$ (hypothesis (ii)), the identity just proved for $W$ transposes to
$h\,\overline W(t,z)a=\overline W(\varphi t,z)\varphi a$.  That $h$ preserves $B_V$: the pulled-back form $B_V(h\,\cdot,h\,\cdot)$ is again a fully invariant
bilinear form on $V_\Leech^+$ taking the value $1$ on the vacuum, and such a form is unique (the library's canonical
invariant form, \lean{canonicalInvariantForm}, used in the proof of \lean{lrLift_leechBV}).  Vacuum and conformal vector of $\Vnat$ lie in $V_\Leech^+$ and are fixed by
$h$.  The file is \leanpath{VOA/Moonshine/MoonshineLRLiftNatural.lean}.
\end{proof}
\begin{physics}
The orbifold theory is assembled from the untwisted sector and the twisted sector glued by the intertwining
fields $W$, $\overline W$.  A symmetry of the untwisted chiral algebra extends to the orbifold as soon as one specifies
how it acts on the twisted sector compatibly with the module structure and the two-point function: the
twisted-untwisted OPE $W$ is \emph{determined} by the module action through locality (skew-symmetry), and
$\overline W$ by the two-point function.  This is the familiar statement that a symmetry of the parent theory that
commutes with the orbifold group lifts to the orbifold once its action on twisted ground states is fixed.
\end{physics}
\begin{status}
Unconditional.  \lean{vnatAutOfPair}, \lean{vnatAutOfPair_apply}, \lean{lrLift_leechBV}: \texttt{\#print axioms} in
\leanpath{VOA/Moonshine/MoonshineLRLiftAudit.lean} gives \lean{[propext, Classical.choice, Quot.sound]}.
\end{status}

\subsection{\texorpdfstring{$\mathrm{Co}_0$}{Co0} on \texorpdfstring{$V_\Lambda^+$}{V\_Lambda+} and on \texorpdfstring{$T^+$}{T+}}

Let $\iota:\Leech\to\Leech$ be a surjective isometry (an element of $\chtwentytwoCo_0=\Aut(\Leech)$).  In the lattice
vertex algebra, $\iota$ acts by the morphism $\Phi_\iota$ of Chapter~\ref{ch:03} (\lean{latMorph}): it rotates the
bosons by $\iota$ and sends $e^\alpha\mapsto c(\alpha)e^{\iota\alpha}$, where the \emph{coboundary}
$c(\alpha)$ (\lean{latCob}) compares the cocycle $\varepsilon$ with $\varepsilon\circ(\iota\times\iota)$.

\begin{definition}[$\chtwentytwoCo_0$ on $V_\Leech^+$]\label{ch22:def:latplus}
For a surjective isometry $\iota$ of $\Leech$, $\mathrm{leechPlusLatMorphAut}(\iota)\in\Aut(V_\Leech^+)$ is the restriction of
$\Phi_\iota$ to $V_\Leech^+$.  (This makes sense because $\Phi_\iota$ is bijective and commutes with the involution
$\theta$.)
\end{definition}
\begin{leancode}
variable (ι : leech.carrier →ₗ[ℤ] leech.carrier)
  (hι : ∀ α β, leech.form (ι α) (ι β) = leech.form α β) (hsurj : Function.Surjective ι)
-- ⋯
def leechPlusLatMorphAut : Aut leechVOAPlus :=
  leechPlusAutOfTheta (latMorphAut ι hι hsurj) fun x => latMorph_thetaLat leech leech ι x
\end{leancode}
(\leanpath{VOA/Moonshine/MoonshineLRLiftLattice.lean}.)
\begin{dictionary}
\lean{leech.carrier →ₗ[ℤ] leech.carrier} is a $\Z$-linear endomorphism of the Leech lattice; \lean{hι} says it
preserves the form, \lean{hsurj} that it is onto (an isometry of a positive-definite lattice is injective, so it is an
automorphism).  \lean{latMorphAut} packages $\Phi_\iota$ as an element of \lean{Aut leechVOA} (bijectivity from
\lean{latMorph_injective} and the range criterion \lean{mem_range_latMorph_iff} of Chapter~\ref{ch:03});
\lean{leechPlusAutOfTheta} restricts a $\theta$-commuting automorphism to $V_\Leech^+$.  The definition is exactly the
textbook one.  It is a construction for \emph{one} $\iota$: $\iota\mapsto\Phi_\iota$ is not a homomorphism on the nose
($\Phi_{\iota\iota'}$ and $\Phi_\iota\Phi_{\iota'}$ differ by a coboundary character), and nothing below needs it to
be.
\end{dictionary}
\begin{physics}
A rotation of the 24 bosons by a lattice symmetry, dressed by the phases that keep the vertex operators $e^\alpha$
mutually local; it commutes with $X\mapsto-X$, so it is a symmetry of the untwisted sector of the orbifold.
\end{physics}
\begin{status}Unconditional; audited in \leanpath{VOA/Moonshine/MoonshineLRLiftAudit.lean}, standard axioms.\end{status}

The twisted sector is $T^+\subset S[\theta]\otimes T$, where $S[\theta]$ is the twisted Fock space of the
half-integrally moded bosons $h(-\tfrac12),h(-\tfrac32),\dots$ and $T$ is the $4096$-dimensional irreducible module of the
twisted group algebra of $\Leech/2\Leech$ (Chapters~\ref{ch:06}--\ref{ch:08}), on which $e^\alpha$ acts by
$e_\alpha$ (\lean{eT}).  The natural candidate for $\varphi_\iota$ is the tensor product of
\begin{itemize}
\item the orthogonal change of oscillator frame on $S[\theta]$ by the matrix $O_\iota$ of $\iota$
  (\lean{liftTwFockMap}), which preserves the Fock form, and
\item a linear bijection $\rho$ of $T$ with $\rho\,e_\alpha=c(\alpha)\,e_{\iota\alpha}\,\rho$ for all $\alpha$.
\end{itemize}

\begin{theorem}[The charge intertwiner exists]\label{ch22:thm:rho}
For every isometry $\iota$ of an even unimodular lattice $L$, there is a $\C$-linear bijection $\rho$
of the extraspecial module with $\rho(e_\alpha t)=c(\alpha)\,e_{\iota\alpha}(\rho t)$ for all $\alpha\in L$, $t$.
\end{theorem}
\begin{leancode}
theorem liftTwChargeIntertwiner_exists :
    ∃ ρ : ExtraT L hEven hU ≃ₗ[ℂ] ExtraT L hEven hU,
      ∀ α t, ρ (eT L hEven hU α t) =
        latCob L L ι α • eT L hEven hU (ι α) (ρ t)
\end{leancode}
(\leanpath{VOA/Moonshine/MoonshineLRLiftTwisted2.lean}; the section variables are a lattice \lean{L}, its evenness
\lean{hEven}, unimodularity \lean{hU : IsUnit (gramMat L).det}, an endomorphism \lean{ι} and the isometry hypothesis
\lean{hι}.)
\begin{proof}[Proof on paper]
The operators $\alpha\mapsto c(\alpha)e_{\iota\alpha}$ satisfy the same multiplication law
$x_\alpha x_\beta=\varepsilon(\alpha,\beta)x_{\alpha+\beta}$ and the same relation modulo $2L$ as the $e_\alpha$
(\lean{liftTwChargeOp_mul}, \lean{liftTwChargeOp_add_two}; this is what the coboundary $c$ is for).  So they define a
second representation of the twisted group algebra of $L/2L$ on $T$.  That algebra is a full matrix algebra (its action
on $T$ is surjective onto $\End T$, \lean{liftTwChargeAction_surjective}), and a full matrix algebra has a unique
irreducible module up to isomorphism; the isomorphism is $\rho$.
\end{proof}
\begin{physics}
The twisted ground states carry a projective representation of the momentum lattice modulo $2$; a rotation of the
lattice maps it to an equivalent projective representation, and $\rho$ is the unitary (up to normalisation) that
implements the equivalence --- a ``Clifford'' intertwiner, as for spinors under rotations.
\end{physics}
\begin{status}Unconditional; audited (\leanpath{VOA/Moonshine/MoonshineLRLiftAudit.lean}), standard axioms.  $\rho$
is unique up to a scalar; the Lean definition \lean{liftTwChargeEquiv} picks one by choice.\end{status}

\begin{theorem}[$\chtwentytwoCo_0$ on $T^+$, covariantly]\label{ch22:thm:twcov}
For a surjective isometry $\iota$ of $\Leech$, the map $\varphi_\iota=O_\iota\otimes\rho$ restricted to $T^+$ satisfies
$\varphi_\iota(u_{(n)}t)=(\Phi_\iota u)_{(n)}\varphi_\iota(t)$ for all $u\in V_\Leech^+$, $n$, $t$; and some nonzero
multiple $\lambda_\iota\varphi_\iota$ preserves \lean{TPForm}.
\end{theorem}
\begin{leancode}
theorem liftTwPlusEquiv_mode (hsurj : Function.Surjective ι)
    (u : (orbifold (Subgroup.zpowers leechTheta)).carrier) (n : ℤ) (t : TP) :
    liftTwPlusEquiv ι hι (mode (plusModule.YM u) n t) =
      mode (plusModule.YM (leechPlusLatMorphAut ι hι hsurj u)) n (liftTwPlusEquiv ι hι t)
\end{leancode}
\begin{leancode}
theorem liftTwExists_formScalar : ∃ c : ℂˣ, ∀ t t' : TP,
    TPForm (liftTwNormalizedPlusEquiv ι hι c t)
      (liftTwNormalizedPlusEquiv ι hι c t') = TPForm t t'
\end{leancode}
(\leanpath{VOA/Moonshine/MoonshineLRLiftTwisted2.lean}, \leanpath{VOA/Moonshine/MoonshineLRLiftTwisted3.lean}.)
\begin{dictionary}
\lean{liftTwPlusEquiv ι hι : TP ≃ₗ[ℂ] TP} is $\varphi_\iota$ on $T^+$ (built from the chosen $\rho$);
\lean{orbifold (Subgroup.zpowers leechTheta)} is the fixed-point subalgebra $V_\Leech^+$, the same type as
\lean{SkewVP}; \lean{liftTwNormalizedPlusEquiv ι hι c} is $c\cdot\varphi_\iota$ for a unit $c\in\C^\times$.  The second
statement is an existence of the scalar, as the textbook statement says.
\end{dictionary}
\begin{proof}[Proof on paper]
\emph{Covariance.}  On the generating fields the claim is a computation: the twisted Heisenberg fields transform by
$O_\iota$ because $O_\iota$ is orthogonal (\lean{liftTwFockMap_modeT}); the twisted vertex operator of $e^\alpha$ is
(dressing by oscillator exponentials)$\otimes e_\alpha$, the dressing transforms covariantly under the frame change
(\lean{liftTwFockMap_DD}), and $\rho e_\alpha=c(\alpha)e_{\iota\alpha}\rho$ supplies exactly the coboundary that
$\Phi_\iota$ puts on $e^\alpha$.  Every state of $V_\Leech$ is a ``word'' (oscillators applied to an $e^\alpha$), and its
twisted field is obtained from the fields of the letters by iterated twisted normal-ordered products; conjugation by
a linear bijection preserves those products (\lean{liftTwConjField_star}), so covariance propagates to all states by
induction on words (\lean{liftTwLatticeEquiv_word}, \lean{liftTwLatticeEquiv_YML}).  Both $O_\iota$ and $\rho$ commute
with $\theta$, so $\varphi_\iota$ restricts to $T^+$.

\emph{The form.}  \lean{TPForm} is the restriction of (Fock form)$\otimes$($T$-form).  The Fock form is invariant
under the orthogonal frame change (\lean{liftTwFockMap_form}).  The $T$-form pulled back by $\rho$ has the same
adjoint operation $A\mapsto A^*$ on $\End T$ as the original (both make $e_\alpha$ adjoint to $\pm e_\alpha$), and two
nondegenerate forms with the same adjoint operation on all of $\End T$ are proportional (a rank-one-operator argument,
\lean{liftTwForm_unique}).  So $\varphi_\iota$ scales \lean{TPForm} by some $z\ne0$, and $\lambda_\iota=z^{-1/2}$ (any square root)
normalises it.
\end{proof}
\begin{physics}
Covariance of every twisted vertex operator is the statement that the rotated twisted sector is again a twisted
sector of the rotated theory --- automatic in a path-integral picture, but here it must be checked mode by mode.  The
form normalisation fixes the overall scale of the twisted ground-state wave functions so that the two-point function
is preserved.
\end{physics}
\begin{status}Unconditional; audited, standard axioms.\end{status}

\begin{definition}[The lift of a Leech isometry]\label{ch22:def:liftTw}
$\mathrm{liftTwVnatAut}(\iota)=\Phi_\iota|_{V_\Leech^+}\oplus \lambda_\iota\varphi_\iota\in\Aut(\Vnat)$, with the scalar $\lambda_\iota$ of
Theorem~\ref{ch22:thm:twcov} chosen once and for all.
\end{definition}
\begin{leancode}
def liftTwVnatAut (hsurj : Function.Surjective ι) : AutMonster.AutVnat :=
  liftTwVnatAutOfForm ι hι hsurj (liftTwExists_formScalar ι hι).choose
    (liftTwExists_formScalar ι hι).choose_spec
-- ⋯
@[simp] theorem liftTwVnatAut_apply (hsurj : Function.Surjective ι)
    (u : SkewVP) (t : TP) :
    liftTwVnatAut ι hι hsurj ((u, t) : Vnat) =
      (leechPlusLatMorphAut ι hι hsurj u,
        liftTwNormalizedPlusEquiv ι hι (liftTwExists_formScalar ι hι).choose t) := rfl

end Normalization
\end{leancode}
(\leanpath{VOA/Moonshine/MoonshineLRLiftTwisted3.lean}.)
\begin{dictionary}
\lean{liftTwVnatAutOfForm} is Theorem~\ref{ch22:thm:natural} applied to the pair; \lean{(…).choose} is the scalar
chosen by \lean{Classical.choice}.  The section variables are \lean{ι} and \lean{hι} as above.  Two features must be
kept in mind.  (a) The lift is a \emph{chosen} element: $\lambda_\iota$ is determined only up to sign and $\rho$ only up to a
scalar, so different choices differ by the extraspecial elements of the next subsection (and by the central element
$z$).  Nothing below depends on which choice was made: every later statement is of the form ``\emph{the chosen} lift
is liftable / has these blocks'', and the visible group is matched by multiplying with extraspecial lifts.  (b)
$\iota\mapsto\mathrm{liftTwVnatAut}(\iota)$ is not claimed to be a homomorphism.
\end{dictionary}
\begin{physics}
The element of $C=C_{\chtwentytwoM}(z)$ above the lattice rotation $\iota\in\chtwentytwoCo_0$; it is defined only up to the
kernel $2^{24}$ of $C\to\chtwentytwoCo_1$ and the centre, exactly as a rotation lifts to the spin group only up to sign.
\end{physics}
\begin{status}Unconditional definition; \lean{liftTwVnatAut} audited, standard axioms.\end{status}


\subsection{The extraspecial group}

For $\mu\in\Leech$, the character automorphism $e^\alpha\mapsto(-1)^{\langle\mu,\alpha\rangle}e^\alpha$ of $V_\Leech$
(\lean{extraLeechAut}) commutes with $\theta$ and restricts to $V_\Leech^+$ (\lean{extraLeechPlusAut}); on the twisted
sector, $\mathrm{id}\otimes e_\mu$ intertwines it (\lean{extraPhi_YML}: $e_\mu e_\alpha=(-1)^{\langle\mu,\alpha\rangle}
e_\alpha e_\mu$ up to the cocycle).  The raw operator scales \lean{TPForm} by the sign
$(-1)^{\langle\mu,\mu\rangle/2}\varepsilon(\mu,\mu)$:
\begin{leancode}
theorem extraPlusPhi_form (x y : TP) :
    TPForm (extraPlusPhi μ x) (extraPlusPhi μ y)
      = ((-1 : ℂ) ^ (leech.form μ μ / 2) * latticeCocycle leech μ μ) * TPForm x y
\end{leancode}
(\leanpath{VOA/Moonshine/MoonshineLRLiftExtra.lean}), so a square root of that sign normalises it
(\lean{extraPlusIsometry}).

\begin{definition}[Extraspecial lifts]\label{ch22:def:extra}
$\mathrm{extraVnatAut}(\mu)=\mathrm{extraLeechPlusAut}(\mu)\oplus\mathrm{extraPlusIsometry}(\mu)\in\Aut(\Vnat)$.
\end{definition}
\begin{leancode}
def extraVnatAut (μ : leech.carrier) : AutVnat :=
  vnatAutOfPair (extraLeechPlusAut μ) (extraPlusIsometryEquiv μ)
    (fun u n t => extraPlusIsometry_mode μ u n t) (fun x y => extraPlusIsometry_form μ x y)
-- ⋯
theorem extraVnatAut_apply (μ : leech.carrier) (u : SkewVP) (t : TP) :
    extraVnatAut μ ((u, t) : Vnat) = (extraLeechPlusAut μ u, extraPlusIsometry μ t)
\end{leancode}
(\leanpath{VOA/Moonshine/MoonshineLRLiftExtraAut.lean}, \leanpath{VOA/Moonshine/MoonshineLRLiftExtra.lean}.)
\begin{dictionary}
Exactly the textbook pair, assembled by Theorem~\ref{ch22:thm:natural}.  As for the lattice lifts, each
$\mathrm{extraVnatAut}(\mu)$ is an individual element: the map $\mu\mapsto\mathrm{extraVnatAut}(\mu)$ is a
homomorphism only up to central signs, which is not claimed (the module docstring says so).
Chapter~\ref{ch:21} proves these elements liftable (\lean{exists_monsterLift_eq_extraVnatAut}).
\end{dictionary}
\begin{physics}
The momentum-phase symmetries $e^{i\pi\langle\mu,p\rangle}$ of the Leech CFT; on twisted ground states they act by the
Clifford generators $e_\mu$, and together they generate the extraspecial group $2^{1+24}$, the normal subgroup of
$C_{\chtwentytwoM}(z)$.
\end{physics}
\begin{status}Unconditional; audited, standard axioms.\end{status}

%======================================================================
\section{Block forms in AxiomMath's coordinates and the rational criterion}\label{ch22:sec:blocks}
%======================================================================

\subsection{The untwisted blocks}

The rational dictionary of Chapter~\ref{ch:21} sends $u=(A,x,N)\in\chtwentytwoB$ to the vector of $\Vnat_2$ whose
$V_\Leech^+$-component is $Q_C(A^\pi)+\sum_c x_c\zeta'_c\,x_{\alpha_c}$ (a quadratic oscillator state in Conway's
coordinates plus a combination of the states $x_\alpha=e^\alpha+e^{-\alpha}$ for the short vectors $\alpha_c$, with phases
$\zeta'$) and whose $T^+$-component is $\lambda\sum_i h_{f_i}(-\tfrac12)\otimes\Psi(N_{\cdot,\pi i})$.  The ``untwisted
dictionary'' \lean{untwistedDictNat ζ} uses the first component only, with a choice of phases $\zeta$; the phases
\lean{griessShortSign} of Chapter~\ref{ch:21} make it rational, and the full dictionary uses
$\zeta'=$ \lean{fullShortSign}, related by the diagonal sign $c\mapsto\varepsilon(\alpha_c,\alpha_c)$ (\lean{griessTrX},
\leanpath{VOA/Moonshine/MoonshineLRGriessTransfer.lean}, documented in Chapter~\ref{ch:21}).

\begin{theorem}[Untwisted block of a lattice lift]\label{ch22:thm:untwistedblock}
For every surjective isometry $\iota$ of $\Leech$ there are a rational $24\times24$ matrix $g_\iota$ (the matrix of $\iota$ in
AxiomMath's coordinates) and a signed permutation $s_\iota$ of the $98280$ short coordinates such that, for all
$A\in\chtwentytwoA$ and $x\in\chtwentytwoX$,
\begin{equation*}
  \mathrm{liftTwVnatAut}(\iota)\bigl(\mathrm{untwistedDictNat}(A,x,0)\bigr)=\mathrm{untwistedDictNat}\bigl(g_\iota A
  g_\iota^{\mathsf T},\ s_\iota(x),\ 0\bigr).
\end{equation*}
\end{theorem}
\begin{leancode}
theorem liftSignNat_untwisted (A : symMat) (x : SpaceX) :
    liftTwVnatAut ι hι hsurj
        (untwistedDictNat griessShortSign (A, x, 0) : Vnat) =
      (untwistedDictNat griessShortSign
        (liftBlockSymMat ι A, liftSignX ι hι x, 0) : Vnat)
\end{leancode}
\begin{leancode}
theorem liftSignWeight_sign (c : ShortIdx) :
    liftSignWeight ι hι c = 1 ∨ liftSignWeight ι hι c = -1
-- ⋯
theorem liftSignX_apply (x : SpaceX) (c : ShortIdx) :
    liftSignX ι hι x (liftSignPerm ι hι c) = liftSignWeight ι hι c * x c
\end{leancode}
(\leanpath{VOA/Moonshine/MoonshineLRLiftSign.lean}; \lean{liftBlockSymMat ι A} $=g_\iota Ag_\iota^{\mathsf T}$ with
$g_\iota=$ \lean{liftBlockMatrix ι}, \leanpath{VOA/Moonshine/MoonshineLRLiftBlock.lean}.)
\begin{dictionary}
\lean{liftSignX ι hι} is the signed permutation: it moves coordinate $c$ to \lean{liftSignPerm ι hι c} (the short
pair containing $\pm\iota\alpha_c$) with weight \lean{liftSignWeight ι hι c} $=\pm1$.  The ``weights are signs'' clause is
the content: a priori the coefficient is $r_c=\zeta_c\,c(\alpha_c)/\zeta_{p(c)}$, a complex number.  The statement is
in the coordinates \lean{untwistedDictNat griessShortSign}; the transfer to \lean{griessNat} conjugates the short block by
\lean{griessTrX} (Chapter~\ref{ch:21}).
\end{dictionary}
\begin{proof}[Proof on paper]
On quadratic oscillator states $h_a(-1)h_b(-1)\mathbb 1$, $\Phi_\iota$ acts by $O_\iota\otimes O_\iota$, which in
matrix language is $M\mapsto O_\iota MO_\iota^{\mathsf T}$, and Conway's frame converts $O_\iota$ into the rational matrix
$g_\iota$ (\lean{liftBlockMatrix_frame}).  On $x_\alpha$, $\Phi_\iota x_\alpha=c(\alpha)x_{\iota\alpha}$.  Every norm-four
vector is $\pm\alpha_c$ for a unique $c$ (\lean{liftSignShort_existsUnique}), so $\iota$ permutes the short pairs.  The
coefficient $r_c$ has square one: the cocycle relation at $\beta=\alpha$ with $c(2\alpha)=1$ gives
$c(\alpha)^2=\varepsilon(\alpha,\alpha)/\varepsilon(\iota\alpha,\iota\alpha)$ (\lean{liftSignLatCob_sq}), the dictionary phases
satisfy $\zeta_c^2\varepsilon(\alpha_c,\alpha_c)=1$, and $\varepsilon(-\beta,-\beta)=\varepsilon(\beta,\beta)$.  Hence
$r_c^2=1$, $r_c=\pm1$.
\end{proof}
\begin{physics}
The weight-two primaries $\partial X^a\partial X^b$ rotate as a symmetric tensor and the $98280$ ``vertex operator''
primaries $\cos(\alpha\cdot X)$ are permuted with signs: in a real basis the lattice symmetries act by signed
permutations, which is why they preserve a rational structure.
\end{physics}
\begin{status}Unconditional for \emph{every} isometry; audited, standard axioms.  (The extraspecial lifts fix the
$\chtwentytwoA$ block and act by $x_c\mapsto(-1)^{\langle\mu,\alpha_c\rangle}x_c$: \lean{liftBlockExtraNat_symBlock},
\lean{liftBlockExtraNat_shortBlock}, same file family, also audited.)\end{status}

\begin{warning}
\textbf{An open, unneeded question.}  Whether $\varepsilon(\alpha_c,\alpha_c)=1$ for all short classes (which would make
\lean{griessTrX} the identity) was left open: it is equivalent to a divisibility $4\mid\sum_i\langle B_i,B_i\rangle
(B.\mathrm{repr}\,\alpha_c)_i^2$ that depends on the chosen basis of $\Leech$.  It is not needed; all block formulas hold
with the sign conjugation.
\end{warning}

\subsection{The twisted block}

On the twisted weight-two space $(T^+)_2$, spanned by $h_v(-\tfrac12)\otimes t$, the lift of $\iota$ acts by
$h_v(-\tfrac12)\otimes t\mapsto h_{O_\iota v}(-\tfrac12)\otimes \lambda_\iota\rho t$ (\lean{liftWNat_hP}), and the extraspecial lift by
$t\mapsto(\text{scalar})\,e_\mu t$.  In AxiomMath's spin coordinates (the isomorphism \lean{spinIso} of the
$4096$-dimensional spin module with $T$, Chapter~\ref{ch:21}), $e_\mu$ is AxiomMath's phased spin matrix
\lean{spinOp} of the class of $\mu$ (\lean{liftWExtra_spinOp}), and the 24 spin columns transform as $N\mapsto
(\lambda_\iota\rho)N\,O_\iota^{\mathsf T}$ (\lean{liftWNat_coordinates}; \leanpath{VOA/Moonshine/MoonshineLRLiftWBlock.lean}).  These are
complex statements; whether $\lambda_\iota\rho$ is rational is the subject of the criterion below.

\subsection{The rational liftability criterion}

Write, for $\alpha\in\Leech$ with class $(a,\delta)\in\mathcal C\times\mathcal C^*$ in $\Leech/2\Leech$ (code word and cocode
element), $Q_\alpha$ for AxiomMath's \emph{unphased} rational spin matrix of that class (\lean{liftMonoQmat α} $=$
\lean{qmatR ℚ 0 a δ}, a signed permutation matrix), and $\mathrm{spinOp}(\alpha)=Z(\alpha)Q_\alpha$ for the phased one,
$Z$ the sign refinement \lean{signRefinement}, $Z^2=\varepsilon(\alpha,\alpha)(-1)^{\langle\alpha,\alpha\rangle/2}$
(\lean{liftMonoRefinement_sq}).

\begin{theorem}[Rational criterion for lattice lifts]\label{ch22:thm:criterion}
Let $\iota$ be a surjective isometry of $\Leech$.
\begin{enumerate}[label=(\alph*)]
\item $\mathrm{liftTwVnatAut}(\iota)$ is the lift of an element of $\chtwentytwoAutO$ if and only if its normalised
  twisted part maps rational spin-dictionary vectors to rational ones.
\item For every $\alpha$ the number $w(\iota,\alpha)=c(\alpha)Z(\iota\alpha)/Z(\alpha)$ is $\pm1$.
\item If $R$ is a rational $4096\times4096$ matrix with $R^{\mathsf T}R=1$ and $RQ_\alpha=w(\iota,\alpha)\,Q_{\iota\alpha}R$
  for all $\alpha\in\Leech$ (equivalently, for the $24$ basis vectors of $\Leech$), then $\mathrm{liftTwVnatAut}(\iota)$ is
  the lift of an element of $\chtwentytwoAutO$.
\end{enumerate}
\end{theorem}
\begin{leancode}
def liftLiftableSpinRational : Prop :=
  ∀ N : SpaceW, ∃ N' : SpaceW,
    liftTwNormalizedPlusEquiv ι hι (liftWScalar ι hι) (spinDict N) = spinDict N'
-- ⋯
theorem liftLiftableIffSpinRational (hsurj : Function.Surjective ι) :
    (∃ g : liftableSubgroup, monsterLift g = liftTwVnatAut ι hι hsurj) ↔
      liftLiftableSpinRational ι hι
\end{leancode}
\begin{leancode}
def liftMonoPhase (α : leech.carrier) : ℂ :=
  latCob leech leech ι α * signRefinement (coordOfRed (red leech (ι α))) /
    signRefinement (coordOfRed (red leech α))
-- ⋯
theorem liftMonoPhase_sign (α : leech.carrier) :
    liftMonoPhase ι α = 1 ∨ liftMonoPhase ι α = -1
-- ⋯
theorem liftMonoOfUnphasedIntertwiner (hsurj : Function.Surjective ι)
    (R : Matrix SpinIdx SpinIdx ℚ) (horth : Rᵀ * R = 1)
    (hinter : ∀ α : leech.carrier,
      R * liftMonoQmat α = liftMonoWeight ι α • (liftMonoQmat (ι α) * R)) :
    ∃ g : liftableSubgroup, monsterLift g = liftTwVnatAut ι hι hsurj
\end{leancode}
\begin{leancode}
theorem liftGenOfBasisIntertwiner (ι : leech.carrier →ₗ[ℤ] leech.carrier)
    (hι : ∀ α β, leech.form (ι α) (ι β) = leech.form α β)
    (hsurj : Function.Surjective ι) (R : Matrix SpinIdx SpinIdx ℚ)
    (horth : Rᵀ * R = 1)
    (hB : ∀ j, R * liftMonoQmat (latticeBasis leech j) =
      liftMonoWeight ι (latticeBasis leech j) •
        (liftMonoQmat (ι (latticeBasis leech j)) * R)) :
    ∃ g : liftableSubgroup, monsterLift g = liftTwVnatAut ι hι hsurj
\end{leancode}
(\leanpath{VOA/Moonshine/MoonshineLRLiftLiftable.lean}, \leanpath{VOA/Moonshine/MoonshineLRLiftLiftable2.lean},
\leanpath{VOA/Moonshine/MoonshineLRLiftLiftable3.lean}; section variables \lean{ι}, \lean{hι} as before.)
\begin{dictionary}
\textbf{Decoding.}  ``Lift of an element of $\chtwentytwoAutO$'' is \lean{∃ g : liftableSubgroup, monsterLift g = …}.
\lean{spinDict N} is the twisted part of the dictionary for $N\in\chtwentytwoW$; \lean{liftWScalar ι hι} is the scalar $\lambda_\iota$
chosen in Definition~\ref{ch22:def:liftTw}.  \lean{liftMonoWeight ι α} is the rational number $\pm1$ equal to
\lean{liftMonoPhase ι α} $=w(\iota,\alpha)$.  \lean{R.map (Rat.castHom ℂ)} would be $R$ read in $\C$; in (c) everything
is rational.  \lean{latticeBasis leech j} is the chosen $\Z$-basis of $\Leech$.

\textbf{Hypotheses.}  In (c): orthogonality of $R$ and the intertwining equation; no hypothesis relates $R$ to the
chosen $\rho$ --- the theorem shows that the chosen normalised $\rho$ \emph{is} $\pm R$.

\textbf{Strength.}  (a) is an equivalence; (c) is a sufficient condition, which is all that is used.
\end{dictionary}
\begin{proof}[Proof on paper]
(a) The untwisted block is already rational (Theorem~\ref{ch22:thm:untwistedblock}), so the lift maps
$\mathrm{griessNat}(\chtwentytwoB)$ into itself iff the twisted block does; then \lean{isVnatLiftable_of_rational}
applies.  Conversely a liftable element maps rational vectors to rational ones.
(b) $c(\alpha)^2=\varepsilon(\alpha,\alpha)/\varepsilon(\iota\alpha,\iota\alpha)$ and
$Z(\beta)^2=\varepsilon(\beta,\beta)(-1)^{\langle\beta,\beta\rangle/2}$, with $\langle\iota\alpha,\iota\alpha\rangle=
\langle\alpha,\alpha\rangle$; so $w^2=1$.  Note that $c(\alpha)$ alone can be $\pm i$ (exactly when the diagonal cocycles of
$\alpha$ and $\iota\alpha$ differ, \lean{liftMonoLatCob_imaginary_iff}); the spin phases absorb it.
(c) Multiply the equation by the phases: $R\,\mathrm{spinOp}(\alpha)=c(\alpha)\,\mathrm{spinOp}(\iota\alpha)R$, which is the
defining equation of $\rho$.  Any two invertible solutions differ by an operator commuting with all $\mathrm{spinOp}(\alpha)$,
hence by a scalar (Schur's lemma for the irreducible module $T$, \lean{liftLiftableSpinCommuteScalar}); so the chosen
$\rho$ is $z R$.  The normalisation $\lambda_\iota$ makes $\lambda_\iota\rho$ orthogonal for the dot product (\lean{liftLiftableChargeDot}), and
$R$ is orthogonal, so $(\lambda_\iota z)^2=1$ and $\lambda_\iota\rho=\pm R$ is rational (\lean{liftLiftableChargeSign}).  Then (a) applies.  The
reduction to the 24 basis vectors holds because the set of $\alpha$ for which the equation holds is closed under
addition and negation (\lean{liftGenIntertwines_add}, \lean{liftGenIntertwines_neg}).
\end{proof}
\begin{physics}
A lattice symmetry acts on the $4096$ twisted ground states by a ``spinor'' matrix.  The theorem says: the symmetry
preserves AxiomMath's rational structure on the full Griess algebra exactly when this spinor matrix can be chosen
rational; and a rational solution of the Clifford intertwining equation, if one exists, \emph{is} the spinor matrix.
\end{physics}
\begin{status}Unconditional; all four declarations audited in \leanpath{VOA/Moonshine/MoonshineLRLiftAudit.lean},
standard axioms.\end{status}

%======================================================================
\section{The generators: Golay signs, Mathieu permutations and Conway's \texorpdfstring{$\xi$}{xi}}\label{ch22:sec:generators}
%======================================================================

In Conway's coordinates $\Leech\subset\Z^{24}$ (vectors $u$ with $u_i\equiv m+2d_i\pmod 4$ for a common $m$ and a Golay
codeword $d$, and $\sum u_i\equiv4m\pmod 8$; \lean{codeLattice}), $\chtwentytwoCo_0$ is generated by three kinds of
elements: sign changes on a Golay codeword, permutations of the $24$ coordinates preserving the Golay code ($M_{24}$),
and one non-monomial element $\xi$.

\begin{definition}[Monomial generators]\label{ch22:def:monomial}
For a Golay codeword $d$, \lean{liftGenGolay c} (Lean's variable is called \lean{c}) negates the coordinates in the
support of $d$; for a permutation $\pi$ of $\{0,\dots,23\}$ preserving the Golay code, \lean{liftGenMathieu σ hσ}
(Lean's variable is called \lean{σ}) permutes the coordinates.
\end{definition}
\begin{leancode}
def liftGenSign (c : Fin 24 → ZMod 2) (u : Fin 24 → ℤ) : Fin 24 → ℤ :=
  fun i => if c i = 0 then u i else -u i
-- ⋯
def liftGenGolay (c : MonsterGriess.Code.C) : leech.carrier →ₗ[ℤ] leech.carrier :=
  ((leechEquiv.trans (liftGenSignEquiv c)).trans leechEquiv.symm).toAddMonoidHom.toIntLinearMap
-- ⋯
def liftGenPreservesCode (σ : Equiv.Perm (Fin 24)) : Prop :=
  ∀ d, permLin (ZMod 2) σ d ∈ MonsterGriess.Code.C ↔ d ∈ MonsterGriess.Code.C
-- ⋯
def liftGenMathieu (σ : Equiv.Perm (Fin 24)) (hσ : liftGenPreservesCode σ) :
    leech.carrier →ₗ[ℤ] leech.carrier :=
  ((leechEquiv.trans (liftGenPermEquiv σ hσ)).trans
    leechEquiv.symm).toAddMonoidHom.toIntLinearMap
\end{leancode}
(\leanpath{VOA/Moonshine/MoonshineLRLiftLiftable3.lean}.)
\begin{dictionary}
\lean{leechEquiv} identifies the abstract Leech lattice of the library with Conway's coordinate model;
\lean{liftGenSignEquiv}, \lean{liftGenPermEquiv} are the coordinate maps, proved to preserve
\lean{codeLattice} (\lean{liftGenSign_mem}: the sum of a Leech vector over a codeword support is divisible by $4$;
\lean{liftGenPerm_mem}).  \lean{permLin (ZMod 2) σ d} is the permuted binary word; the code-preservation condition is
an ``iff'', i.e.\ $\pi$ maps the code onto itself.  Isometry and surjectivity are separate lemmas
(\lean{liftGenGolay_isometry}, \lean{liftGenMathieu_surjective}, \dots).  Exactly the textbook monomial group
$2^{12}{:}M_{24}$, with no choice of generators.
\end{dictionary}
\begin{physics}
Sign flips and permutations of the $24$ bosons that preserve the lattice: the ``monomial'' symmetries.
\end{physics}
\begin{status}Unconditional; audited, standard axioms.\end{status}

\begin{theorem}[The monomial group lifts]\label{ch22:thm:monomial}
For every Golay codeword $d$ and every code-preserving permutation $\pi$, the chosen lifts
$\mathrm{liftTwVnatAut}(\mathrm{liftGenGolay}\,d)$ and $\mathrm{liftTwVnatAut}(\mathrm{liftGenMathieu}\,\pi)$ are lifts of
elements of $\chtwentytwoAutO$.
\end{theorem}
\noindent File \leanpath{VOA/Moonshine/MoonshineLRLiftLiftable4.lean}:
\begin{leancode}
theorem liftSpinGolay_liftable (c : MonsterGriess.Code.C) :
    ∃ g : liftableSubgroup, monsterLift g =
      liftTwVnatAut (liftGenGolay c) (liftGenGolay_isometry c) (liftGenGolay_surjective c)
-- ⋯
theorem liftSpinMathieu_liftable (σ : Equiv.Perm (Fin 24)) (hσ : liftGenPreservesCode σ) :
    ∃ g : liftableSubgroup, monsterLift g = liftTwVnatAut (liftGenMathieu σ hσ)
      (liftGenMathieu_isometry σ hσ) (liftGenMathieu_surjective σ hσ)
\end{leancode}
\begin{leancode}
def liftSpinFrameLattice (α : leech.carrier) : Prop :=
  (∃ (n : ℤ) (k : Fin 24 → ℤ),
    ∀ i, (leechEquiv α : Fin 24 → ℤ) i = 2 * n + 4 * k i) ∧
    8 ∣ ∑ i, (leechEquiv α : Fin 24 → ℤ) i
-- ⋯
theorem liftSpinRational_of_diagonal {T : Type*} (E : (SpinIdx → ℂ) ≃ₗ[ℂ] (SpinIdx → ℂ))
    (horth : ∀ v w, E v ⬝ᵥ E w = v ⬝ᵥ w)
    (d e : T → SpinIdx → ℂ)
    (hint : ∀ t v, E (Matrix.diagonal (d t) *ᵥ v) =
      Matrix.diagonal (e t) *ᵥ E v)
    (hsep : ∀ j k, (∀ t, d t j = d t k) → j = k) :
    ∃ (σ : Equiv.Perm SpinIdx) (ε : SpinIdx → ZMod 2),
      ∀ v, E v = (liftGenSignedMatrix σ ε).map (Rat.castHom ℂ) *ᵥ v
\end{leancode}
(\leanpath{VOA/Moonshine/MoonshineLRLiftLiftable4.lean}.)
\begin{dictionary}
The two theorems have no hypotheses beyond the data ($d$ a codeword, $\pi$ code-preserving; Lean's \lean{c}, \lean{σ}).  The auxiliary
\lean{liftSpinFrameLattice α} says: all Conway coordinates of $\alpha$ are $\equiv2n\pmod4$ for one $n$, and their sum
is $\equiv0\pmod8$.  \lean{liftSpinRational_of_diagonal} is a general lemma: an orthogonal bijection $E$ of $\C^{4096}$
that carries a separating family of diagonal operators $\mathrm{diag}(d_t)$ to diagonal operators is a rational signed
permutation matrix (\lean{liftGenSignedMatrix σ ε}, entries $\pm1$ placed by a permutation).
\end{dictionary}
\begin{proof}[Proof on paper]
The spin matrices $Q_\alpha$ of the vectors of the ``frame sublattice'' \lean{liftSpinFrameLattice} commute with each
other (their pairings are even) and include the frame involution and all even cocode classes, whose matrices are
diagonal in AxiomMath's spin basis and separate the $4096$ basis vectors (\lean{liftSpinTest_separates}).  Golay sign
changes and code-preserving permutations preserve the frame sublattice (\lean{liftSpinGolay_frame},
\lean{liftSpinMathieu_frame}).  The normalised charge operator $\lambda_\iota\rho$ intertwines $Q_\alpha$ with $\pm Q_{\iota\alpha}$, so it
carries the separating diagonal family to a diagonal family; a linear map with this property sends each basis line to a
basis line, i.e.\ it is monomial, and orthogonality forces the nonzero entries to be $\pm1$.  So $\lambda_\iota\rho$ is a rational
signed permutation (\lean{liftSpinSigned_exists}) satisfying the intertwining equation for every $\alpha$, and
Theorem~\ref{ch22:thm:criterion}(c) applies.  No enumeration of the $4096$-dimensional module is used; the permutation
and signs are obtained by choice (\lean{liftSpinGolayPerm}, \lean{liftSpinMathieuSigns}, \dots), not by closed formulas.
\end{proof}
\begin{physics}
Monomial symmetries permute the twisted ground states (up to signs) because they preserve a maximal commuting set of
Clifford generators --- the analogue of a rotation that maps coordinate axes to coordinate axes acting on spinors by a
signed permutation of the weight basis.
\end{physics}
\begin{status}Unconditional; both theorems audited, standard axioms.\end{status}


\begin{definition}[Conway's $\xi$]\label{ch22:def:xi}
Number the coordinates by MOG columns, $\mathrm{coord}(r,j)=r+4j$ ($r\in\{0,1,2,3\}$ the row, $j\in\{0,\dots,5\}$ the
column).  On each column $\xi$ acts by the $4\times4$ matrix $\mathrm{diag}(1,-1,-1,-1)\,\tfrac12(J-2I)$, $J$ the all-ones
matrix:
\begin{equation*}
  (\xi u)_{(0,j)}=\tfrac12 s_j-u_{(0,j)},\qquad (\xi u)_{(r,j)}=u_{(r,j)}-\tfrac12 s_j\ (r\ne0),\qquad
  s_j=\textstyle\sum_{r}u_{(r,j)}.
\end{equation*}
\end{definition}
\begin{leancode}
def liftXiBlock (u : Fin 4 → ℚ) : Fin 4 → ℚ :=
  fun r => if r = 0 then (∑ s, u s) / 2 - u r else u r - (∑ s, u s) / 2
-- ⋯
theorem liftXi_cube (α : leech.carrier) : liftXi (liftXi (liftXi α)) = α
-- ⋯
theorem liftXi_isometry (α β : leech.carrier) :
    leech.form (liftXi α) (liftXi β) = leech.form α β
-- ⋯
theorem liftXiMatrix_nonmonomial :
    liftXiMatrix 0 0 = -1 / 2 ∧ liftXiMatrix 0 1 = 1 / 2
\end{leancode}
(\leanpath{VOA/Moonshine/MoonshineLRLiftXi.lean}.)
\begin{dictionary}
\lean{liftXiBlock} is the column formula above; \lean{liftXi : leech.carrier →ₗ[ℤ] leech.carrier} is the induced
integral map, defined through \lean{liftXiInt}, which uses that column sums of Leech vectors are even
(\lean{liftXiCol_parity}) and preserves the lattice (\lean{liftXiInt_mem}, from the hexacode description of the
Golay code).  $\xi^3=1$, $\xi$ is an isometry and surjective (\lean{liftXi_surjective}); \lean{liftXiMatrix} is its
rational $24\times24$ matrix, and the two entries $-\tfrac12,\tfrac12$ in one row witness that it is not monomial.
\end{dictionary}
\begin{physics}
A genuinely ``non-classical'' lattice symmetry that mixes the four bosons of each MOG column by a Hadamard-type
rotation; it is what makes $\chtwentytwoCo_0$ much bigger than the monomial group.
\end{physics}
\begin{status}Unconditional; audited, standard axioms.\end{status}

\begin{theorem}[$\xi$ lifts]\label{ch22:thm:xi}
$\mathrm{liftTwVnatAut}(\xi)$ is the lift of an element of $\chtwentytwoAutO$.
\end{theorem}
\begin{leancode}
theorem liftXiCharge_projective :
    ∃ (T : Matrix SpinIdx SpinIdx ℚ) (z : ℂ), z ≠ 0 ∧
      ∀ v : SpinIdx → ℂ, T.map (Rat.castHom ℂ) *ᵥ v = z • liftXiCharge v
-- ⋯
theorem liftXi_liftable : ∃ g : liftableSubgroup,
    monsterLift g = liftTwVnatAut liftXi liftXi_isometry liftXi_surjective
\end{leancode}
(\leanpath{VOA/Moonshine/MoonshineLRLiftXi.lean}.)
\begin{dictionary}
\lean{liftXiCharge} is the normalised charge operator $\lambda_\iota\rho$ of $\xi$ in spin coordinates.  The first theorem says it
is a complex multiple of a rational matrix; the second is the conclusion, with no hypotheses.
\end{dictionary}
\begin{proof}[Proof on paper]
\emph{Rational descent.}  Let $M$ be the complex matrix of $\lambda_\iota\rho$.  The intertwining equation
$MQ_\alpha=w(\xi,\alpha)Q_{\xi\alpha}M$ has rational coefficients.  For a $\Q$-linear functional $f:\C\to\Q$, applying $f$
entrywise commutes with multiplication by rational matrices (\lean{liftXiDual_intertwine}), so $f(M)$ is a rational
solution of the same equation; for a suitable $f$ it is nonzero.  By Schur's lemma for $T$, $f(M)=zM$ with
$z\in\C^\times$ (\lean{liftXiCharge_projective}).  Orthogonality of $M$ makes $f(M)^{\mathsf T}f(M)=z^2$, so
$z^2\in\Q$.  Since $\xi^3=1$, $M^3$ intertwines each $Q_\alpha$ with $\pm Q_\alpha$, hence normalises the separating diagonal
family, hence is a rational signed permutation (\lean{liftXiCharge_cube_signed}); then $f(M)^3=z^3M^3$ gives
$z^3\in\Q$.  So $z=z^3/z^2\in\Q$ and $M=z^{-1}f(M)$ is rational.  Theorem~\ref{ch22:thm:criterion}(c) applies with
$R=M$ (\lean{liftXiSpinMatrix}, chosen noncomputably; no entrywise Hadamard formula is asserted).
\end{proof}
\begin{physics}
The spinor matrix of the Hadamard rotation is itself a (normalised) Hadamard-type matrix with entries in
$2^{-k}\Z$; the argument proves its rationality without ever writing it down.
\end{physics}
\begin{status}Unconditional; audited, standard axioms.\end{status}

%======================================================================
\section{The visible group acts on \texorpdfstring{$V^\natural$}{V-natural}}\label{ch22:sec:visible}
%======================================================================

AxiomMath's visible group is defined uniformly, as the image of \lean{thetaEquiv} on compatible pairs
$p=(s,g)\in\mathrm{Pvis}$, not by generators.  The lifting is therefore done uniformly too.

\begin{theorem}[Every visible element is liftable]\label{ch22:thm:visible}
Every element of $\chtwentytwoGvis$, regarded as an element of $\chtwentytwoAutO$, is liftable.  Consequently
Conway--Norton holds for every element of $\chtwentytwoAutO$ as soon as some liftable $\tau$ together with
$\chtwentytwoGvis$ generates $\chtwentytwoAutO$.
\end{theorem}
\begin{leancode}
theorem liftVisRAllVisible : ∀ (g : SpaceB ≃ₗ[ℚ] SpaceB) (hg : g ∈ Gvis),
    IsVnatLiftable ⟨g, Gvis_le_AutO hg⟩
\end{leancode}
(\leanpath{VOA/Moonshine/MoonshineLRLiftVisible2.lean}.)
\begin{leancode}
def visibleSetAutO : Set AutO := {g | (g : SpaceB ≃ₗ[ℚ] SpaceB) ∈ Gvis}
-- ⋯
theorem isVnatLiftable_of_mem_visibleSetAutO : ∀ g ∈ visibleSetAutO, IsVnatLiftable g
-- ⋯
theorem conwayNorton_autO_of_tauB {τ : AutO} (hτ : IsVnatLiftable τ)
    (hGmat : Subgroup.closure (insert τ visibleSetAutO) = ⊤) (g : AutO) :
    ∃ h : AutVnat, vnatWeightTwoRestrict h = griessConj (g : SpaceB ≃ₗ[ℚ] SpaceB) ∧
      ∃ (Γ : Subgroup (GL (Fin 2) ℝ)) (N : ℕ), IsMoonshineGroup Γ N ∧
        IsNormalizedHauptmodul Γ (mckayThompsonFun h)
\end{leancode}
(\leanpath{VOA/Moonshine/MoonshineLRMonsterVisible.lean}; full names \chtwentytwolean{VOA.Moonshine.ConwayNorton.M8.liftVisRAllVisible},
\lean{...M8.conwayNorton_autO_of_tauB}.)
\begin{dictionary}
\textbf{Decoding.}  \lean{⟨g, Gvis_le_AutO hg⟩ : AutO} is $g$ regarded as an element of $\chtwentytwoAutO$ (using
$\chtwentytwoGvis\le\chtwentytwoAutO$).  \lean{visibleSetAutO} is the same set inside \lean{AutO}.  In
\lean{conwayNorton_autO_of_tauB}: \lean{hτ} says $\tau$ is liftable; \lean{hGmat} says
$\langle\tau,\chtwentytwoGvis\rangle=\chtwentytwoAutO$ (as subgroups of \lean{AutO}, \lean{⊤} the whole group).  The
conclusion for $g$: an automorphism $h$ of $\Vnat$ acting on $\Vnat_2$ as $g$ does, whose $T_h$ is a normalised
Hauptmodul for a moonshine group --- the clauses of \lean{ConwayNortonFor h} spelled out.

\textbf{Hypotheses.}  \lean{liftVisRAllVisible}: none.  \lean{conwayNorton_autO_of_tauB}: \lean{hτ} and
\lean{hGmat}, both genuine inputs at this point; they are discharged for $\tau=\tau_B$ in \S\ref{ch22:sec:tauB}
(\lean{hτ} unconditionally, \lean{hGmat} only under \lean{GmatEqAutO}).
\end{dictionary}
\begin{proof}[Proof on paper]
Let $p=(s,g)\in\mathrm{Pvis}$ and $\iota_g$ the Leech isometry with matrix $g$ (\lean{liftVisIsometry}; AxiomMath's
matrices act on row vectors, the library's on columns, hence a transpose, \lean{liftVisBlockMatrix}).  Compatibility of
$p$ says $s$ conjugates the extraspecial group as $\iota_g$ does: $s^{\mathsf T}Q_\alpha=w'(p,\alpha)Q_{\iota_g\alpha}s^{\mathsf T}$ with
signs $w'$ (\lean{liftVisIntertwine}).  Compare with the chosen charge operator $\lambda_\iota\rho$ of $\iota_g$, which satisfies the same intertwining
equation with the signs $w(\iota_g,\alpha)$ (Theorem~\ref{ch22:thm:criterion}).  The relative operator
$D=(s^{\mathsf T})^{-1}\lambda_\iota\rho$ is a priori a \emph{complex} orthogonal matrix, and it commutes with every $Q_\alpha$
up to the sign $e_\alpha=w(\iota_g,\alpha)w'(p,\alpha)$ (\lean{liftVisRelativeCharge_intertwine}).  Three steps:
\begin{enumerate}[label=(\roman*)]
\item \emph{$D$ is rational.}  Among the $Q_\alpha$ are the diagonal matrices of the frame-sublattice vectors, which separate
  the $4096$ basis vectors (as in Theorem~\ref{ch22:thm:monomial}).  $D$ conjugates each of them to $\pm$ itself, so it maps
  each basis line to a basis line: $D$ is monomial, and orthogonality makes its entries $\pm1$.  So $D$ is a rational signed
  permutation matrix (\lean{liftVisRelativeCharge_signed}).
\item \emph{$D=\pm Q_\mu$.}  The $2^{24}$ matrices $Q_\mu$ ($\mu$ running over $\Leech/2\Leech$) form a basis of the
  matrix algebra $M_{4096}(\Q)$, and conjugation by $Q_q$ multiplies $Q_\mu$ by the commutator sign $\kappa(\mu,q)=\pm1$;
  distinct $\mu$ give distinct characters $q\mapsto\kappa(\mu,q)$.  Expand $D=\sum_\mu d_\mu Q_\mu$.  The relation
  $DQ_q=e_qQ_qD$ says that every $\mu$ with $d_\mu\ne0$ has $\kappa(\mu,\cdot)=e_\cdot$; so only one $\mu$ occurs (in Lean:
  some $\mu$ has $\tr(Q_\mu^{\mathsf T}D)\ne0$, \lean{liftVisRTrace_nonzero}, and then its character is $e$,
  \lean{liftVisRTrace_sign}).  Hence $D=d\,Q_\mu$, and $D^{\mathsf T}D=1$ forces $d=\pm1$ (\lean{liftVisRExtraspecial},
  \lean{liftVisRRelative}).
\end{enumerate}
\emph{Cancelling the discrepancy.}  So the chosen lift of $\iota_g$ is liftable with blocks
$(g^{\mathsf T}Ag,\ \cdot\ ,\ s^{\mathsf T}(\pm Q_\mu)Ng)$ (\lean{liftVisRChosen}).  Extraspecial lifts realise every
$\pm Q_\mu$ on $\chtwentytwoW$ while fixing $\chtwentytwoA$ (\lean{liftVisRSpinLift}; the central $-1$ is the square of an
extraspecial lift of odd quadratic norm, \lean{liftVisRCentral}), so multiplying by one of them cancels the discrepancy.
The product agrees with $\mathrm{thetaEquiv}(p)$ on the $\chtwentytwoA$ and $\chtwentytwoW$ blocks; two automorphisms
commuting with $x_B$ that agree on those blocks are equal, because the $\chtwentytwoX$ block is determined by the action on
$\chtwentytwoW$ through the product $\chtwentytwoX\times\chtwentytwoW\to\chtwentytwoW$, which detects $x$
(\lean{liftVisParity_ext}).  Hence $\mathrm{thetaEquiv}(p)$ is liftable (\lean{liftVisRLiftable}).  The second part is
\lean{conwayNorton_autO_of_generators} of Chapter~\ref{ch:21} with its \lean{hvis} discharged.
\end{proof}
\begin{physics}
Every element of $C_{\chtwentytwoM}(z)=2^{1+24}.\chtwentytwoCo_1$, in AxiomMath's presentation, is a symmetry of the full
moonshine CFT.  The bookkeeping of signs is the usual ambiguity of lifting $\chtwentytwoCo_0$ to the extraspecial
extension, resolved by multiplying with a momentum-phase symmetry.
\end{physics}
\begin{status}
\lean{liftVisRAllVisible}: unconditional.  \lean{conwayNorton_autO_of_tauB}: conditional on its arguments \lean{hτ},
\lean{hGmat}.  Both audited in \leanpath{VOA/Moonshine/MoonshineLRLiftAudit.lean}, standard axioms.
\end{status}


%======================================================================
\section{Triality: \texorpdfstring{$\sigma$}{sigma} is rational on the Griess algebra}\label{ch22:sec:sigma}
%======================================================================

\subsection{Frame sectors}

Recall from Chapters~\ref{ch:11}--\ref{ch:14} the Golay frame sublattice $\Leech_0\subset\Leech$ of index two.  It splits
each half of $\Vnat$ in two:
\begin{itemize}
\item $V_\Leech^+=V_0^+\oplus V_3^+$, where $V_0^+$ (Lean \lean{frameSector 0}) is spanned by the states with momentum in
  $\Leech_0$ and $V_3^+$ (Lean \lean{frameSector 1}) by those with momentum in the other coset.  (The names $V_0,V_3$ are
  inherited from Chapters~\ref{ch:11}--\ref{ch:14}, where the frame sectors are labelled as in FLM; ``3'' is a label, not
  a number of anything.)  Equivalently, they are the $\pm1$ eigenspaces of the frame involution \lean{plusFrame}.
\item $T^+=T_0^+\oplus T_1^+$ (Lean \lean{TPframe 0}, \lean{TPframe 1}), the two eigenspaces of the frame involution on the
  twisted sector, indexed by $\Z/2$.
\end{itemize}
FLM's triality $\sigma$ is assembled from \emph{triality data} $\mathfrak D$ (Lean \lean{SigmaData}, Chapter~\ref{ch:14}):
an automorphism $\sigma_0$ of $V_0^+$ (\lean{v0Sigma}), a label $s=\mathfrak D.s\in\Z/2$, a $\sigma_0$-twisted intertwiner
$\Phi:V_3^+\to T_s^+$ with inverse $\Psi$, and a map $\Theta$ of the other twisted sector $T_{s+1}^+$ to itself.  So $\sigma$
preserves $V_0^+$ and $T_{s+1}^+$ and exchanges $V_3^+$ with $T_s^+$.  Which of the two twisted frame sectors is exchanged
with $V_3^+$ (the label $s$) is fixed by the construction of $\Phi$ in Chapter~\ref{ch:13}; nothing below depends on its
value.  $\sigma$ is an involution.  Chapter~\ref{ch:14} proved that the assembled operator is an
automorphism of $\Vnat$; in Lean it is \lean{sigmaAut : AutVnat} (\leanpath{VOA/Moonshine/AutSigma.lean},
Chapter~\ref{ch:20}).  In Conway's coordinates the short vectors split by parity: $\alpha_c$ lies in $\Leech_0$ iff all its
coordinates are even (shapes $(\pm4)^2$ and $(\pm2)^8$), and in the other coset iff all are odd (shape
$(\mp3,\pm1^{23})$) (\lean{sigmaBlockShort_even_iff}, \lean{sigmaBlockShort_odd_iff}).

\begin{theorem}[Sector split and the conjugate of $z$]\label{ch22:thm:sector}
(a) For every $u=(A,x,N)$, $\sigma$ sends the untwisted dictionary vector of $u$ to the pair
$\bigl(\sigma_0(\text{matrix part}+\text{even short part}),\ \Phi(\text{odd short part})\bigr)$.
(b) For any triality data $\mathfrak D$, $\sigma\,z\,\sigma$ is the frame-sector sign operator, $+1,-1,+1,-1$ on
$V_0^+,V_3^+,T_s^+,T_{s+1}^+$.
\end{theorem}
\begin{leancode}
theorem sigmaBlockNat_split (ζ : ShortIdx → ℂ) (u : SpaceB) :
    sigmaAut (untwistedDictNat ζ u : Vnat) =
      (((v0Sigma (sigmaBlockEvenState ζ u) : V0Sub.carrier) : VPlus),
        Dnat.Φ (shortDict ζ (sigmaBlockShortPart 1 u.2.1)))
-- ⋯
def sigmaBlockZ : Vnat →ₗ[ℂ] Vnat where
  toFun v := (v.1, -v.2)
-- ⋯
theorem sigmaBlockConjugate (D : SigmaData) :
    sigNat D ∘ₗ sigmaBlockZ ∘ₗ sigNat D = sigmaBlockFrameSign D.s
\end{leancode}
(\leanpath{VOA/Moonshine/MoonshineLRSigmaBlock.lean}.)
\begin{dictionary}
\lean{sigmaBlockEvenState ζ u} is the $V_0^+$-component (matrix part and even short coordinates);
\lean{sigmaBlockShortPart 1} the rational projection onto odd short coordinates; \lean{Dnat} the triality data used to
build \lean{sigmaAut}, \lean{Dnat.Φ} its twisted half.  \lean{sigNat D} is the involution assembled from data \lean{D} (our $\mathfrak D$);
\lean{sigmaBlockZ} is $z$; \lean{sigmaBlockFrameSign D.s} is the sign operator described in (b) (its definition, a block
operator built from the frame involution \lean{plusFrame} and the two projections of $T^+$, is in the same file).
\lean{∘ₗ} is composition of linear maps.  Both statements are the textbook ones.
\end{dictionary}
\begin{proof}[Proof on paper]
(a) is the definition of $\sigma$ sector by sector, plus the parity criterion for short vectors.  (b): $\sigma$ preserves
each of $V_0^+$ and $T_{s+1}^+$ and exchanges $V_3^+$ with $T_s^+$; $z$ is $+1$ on $V_0^+\oplus V_3^+$ and $-1$ on
$T_s^+\oplus T_{s+1}^+$; conjugating, $\sigma z\sigma$ is $+1$ on $V_0^+$, $-1$ on $\sigma(T_s^+)=V_3^+$, $+1$ on
$\sigma(V_3^+)=T_s^+$, $-1$ on $T_{s+1}^+$.
\end{proof}
\begin{physics}
\textbf{Triality exchanges the orbifold sectors.}  $\Vnat$ has two commuting $\Z/2$ symmetries visible in the
Leech description: $z$ (untwisted vs.\ twisted; $+1$ on $V_0^+\oplus V_3^+$, $-1$ on $T_0^+\oplus T_1^+$) and the frame sign
$z'$ of part (b) ($+1$ on $V_0^+\oplus T_s^+$, $-1$ on $V_3^+\oplus T_{s+1}^+$).  Their product $zz'$ is a third.  Each of the three involutions in the four-group
$\{1,z,z',zz'\}$ exhibits $\Vnat$ as a $\Z/2$-orbifold of a Leech-type theory, with ``untwisted'' and ``twisted'' sectors
that are different unions of the four frame sectors.  Triality is the symmetry that permutes these three orbifold
descriptions; part (b) is the statement that $\sigma$ conjugates $z$ into $z'$.  (That each of the three involutions exhibits $\Vnat$
as such an orbifold is the classical picture, not a Lean statement.)  In AxiomMath's language the four-group is
$\{1,x_B,y_B,z_B\}$ and $\tau_B$ permutes it cyclically; through the dictionary $x_B$ is $z$, and $\{y_B,z_B\}=\{z',zz'\}$,
which of the two being $z'$ is decided in Theorem~\ref{ch22:thm:tauB}.
\end{physics}
\begin{status}Unconditional; audited (\leanpath{VOA/Moonshine/MoonshineLRLiftAudit.lean}), standard axioms.\end{status}

\subsection{\texorpdfstring{$\sigma$}{sigma} on the even block}

\begin{theorem}[$\sigma_0$ is rational]\label{ch22:thm:sigmaeven}
There is a $\Q$-linear map $F:\chtwentytwoB\to\chtwentytwoB$ such that for every $A\in\chtwentytwoA$ and every $x\in\chtwentytwoX$
supported on even short classes,
$\sigma\bigl(\mathrm{untwistedDictNat}(A,x,0)\bigr)=\mathrm{untwistedDictNat}(F(A,x,0))$.  Explicitly: diagonal matrices are
fixed; for $j\ne k$ the off-diagonal matrix unit $U=U_{jk}$ and the two short coordinates $P,M$ of the pairs
$\pm(4e_j+4e_k)$, $\pm(4e_j-4e_k)$ span a $\sigma$-stable \emph{three-dimensional} space, on which
\begin{equation*}
  U\mapsto\tfrac12(pP+mM),\quad P\mapsto pU+\tfrac12P-\tfrac{pm}2M,\quad M\mapsto mU-\tfrac{pm}2P+\tfrac12M,\qquad
  p,m\in\{\pm1\};
\end{equation*}
this $3\times3$ matrix squares to $1$ and has trace $1$, so it is a reflection (eigenvalues $1,1,-1$), as it must be for the
involution $\sigma$;
and on the $64$ short classes of shape $(\pm2)^8$ supported on one octad, $\sigma$ acts by a Hadamard matrix with entries
$\pm\tfrac18$.
\end{theorem}
\begin{leancode}
theorem sigmaBlockRWeight_values {j k : Fin 24} (hjk : j ≠ k) :
    (sigmaBlockRPWeight hjk = 1 ∨ sigmaBlockRPWeight hjk = -1) ∧
      (sigmaBlockRMWeight hjk = 1 ∨ sigmaBlockRMWeight hjk = -1)
-- ⋯
theorem sigmaBlockOKS (S T : sigmaBlockOSubsets D) :
    griessShortSign (sigmaBlockOIndex hDA hD8 S) * KS hDA S.val T.val =
      ((sigmaBlockOWeight hDA hD8 S T : ℂ) / 8) *
        griessShortSign (sigmaBlockOIndex hDA hD8 T)
-- ⋯
theorem sigmaBlockOFull_nat (A : symMat) (x : SpaceX)
    (hx : ∀ c, sigmaBlockParity c ≠ 0 → x c = 0) :
    sigmaAut (untwistedDictNat griessShortSign (A, x, 0) : Vnat) =
      (untwistedDictNat griessShortSign (sigmaBlockOFull (A, x, 0)) : Vnat)
\end{leancode}
(\leanpath{VOA/Moonshine/MoonshineLRSigmaBlock2.lean}, \leanpath{VOA/Moonshine/MoonshineLRSigmaBlock3.lean}.)
\begin{dictionary}
\lean{sigmaBlockRPWeight hjk}, \lean{sigmaBlockRMWeight hjk} are $p,m$; $U_{jk}$ is \lean{sigmaBlockRH j k}; the formulas for $U,P,M$ are
\lean{sigmaBlockRNat_H}, \lean{sigmaBlockRNat_P}, \lean{sigmaBlockRNat_M}.  In \lean{sigmaBlockOKS}, \lean{KS} is the
complex transport coefficient of FLM's $\sigma$ between two octad classes $S,T$ (even subsets of one octad, called \lean{D} in Lean), and the
statement says that in the rational dictionary it becomes \lean{sigmaBlockOWeight S T}$/8$ with weight $\pm1$
(\lean{sigmaBlockOWeight_values}).  \lean{sigmaBlockOFull} is the assembled rational map $F$; the hypothesis \lean{hx}
says $x$ vanishes on odd classes.  All signs are relative to the noncomputable phases \lean{griessShortSign}; no sign
table is asserted.
\end{dictionary}
\begin{proof}[Proof on paper]
$\sigma_0=$ \lean{v0Sigma} is defined in Chapter~\ref{ch:11} on explicit generators of $V_0^+$ with explicit complex
coefficients.  The content here is that, measured against the dictionary phases $\zeta$, every coefficient is rational.
For the two-root vectors: the transport coefficients $c_Q t_c$ have the same ``cocycle-normalised square''
$z^2\varepsilon(a,a)=1$ as $\zeta^2$, so their ratio is a sign (\lean{sigmaBlockRRatio_sq}).  For octads: all even subsets of
one octad have the same normalised transport square (\lean{sigmaBlockOTransport_sq}), so the ratios of phases are signs and
FLM's coefficient $K_S(T)=\pm\tfrac18$ times a ratio of phases becomes a rational $\pm\tfrac18$.  Every even short class is of
two-root or octad type (\lean{sigmaBlockOColumn_exists}); assembling the columns gives $F$.
\end{proof}
\begin{physics}
On the subalgebra $V_0^+$ common to two orbifold descriptions, $\sigma$ is the change of basis between them.  It mixes
each ``current bilinear'' $\partial X^j\partial X^k$ with the vertex operators of the root vectors $\pm4e_j\pm4e_k$ by a
\emph{reflection} of the three-dimensional space they span (an order-two operation, since $\sigma^2=1$), and it acts on
each octad's $64$ vertex operators by a normalised Hadamard (discrete Fourier) matrix, again an involution.  The order-three
element that cyclically rotates such triples is AxiomMath's $\tau_B$, not $\sigma$ (\S\ref{ch22:sec:tauB}).
\end{physics}
\begin{status}Unconditional; audited, standard axioms.\end{status}


\subsection{The odd columns and the twisted blocks}

\begin{theorem}[$\sigma$ is rational on $\Vnat_2$, and is a Monster element]\label{ch22:thm:sigmarational}
(a) $\sigma$ is rational on all of $\Vnat_2$ if and only if, for every odd short class $c$,
$\zeta'_c\,\Phi(x_{\alpha_c})$ is the spin-dictionary vector of a rational $N\in\chtwentytwoW$.
(b) This holds, with an explicit column.  (c) Hence $\sigma$ maps $\mathrm{griessNat}(\chtwentytwoB)$ into itself, and $\sigma$
is the lift of an element $\sigma_B$ of $\chtwentytwoAutO$.
\end{theorem}
\begin{leancode}
theorem sigmaB_rational_iff_odd_columns :
    (∀ u : SpaceB, ∃ v : SpaceB, sigmaAut (griessNat u : Vnat) = griessNat v) ↔
      (∀ c : ShortIdx, sigmaBlockParity c = 1 → ∃ N : SpaceW,
        fullShortSign c • Dnat.Φ (frameX0P (shortVec c)) = spinDict N)
-- ⋯
theorem sigmaBOdd_columns (c : ShortIdx) (hc : sigmaBlockParity c = 1) :
    fullShortSign c • Dnat.Φ (frameX0P (shortVec c)) =
      spinDict (sigmaBOddColumn sigmaBOddSeedQ c)
-- ⋯
theorem sigmaBOdd_rational :
    ∀ u : SpaceB, ∃ v : SpaceB, sigmaAut (griessNat u : Vnat) = griessNat v
-- ⋯
theorem sigmaBOdd_liftable : ∃ g : liftableSubgroup, monsterLift g = sigmaAut
\end{leancode}
(\leanpath{VOA/Moonshine/MoonshineLRSigmaB.lean}, \leanpath{VOA/Moonshine/MoonshineLRSigmaB2.lean}.)
\begin{dictionary}
\lean{sigmaBlockParity c = 1}: $c$ is an odd class.  \lean{fullShortSign c} is $\zeta'_c$; \lean{frameX0P (shortVec c)}
is $x_{\alpha_c}$; \lean{spinDict N} the twisted dictionary vector of $N$; \lean{sigmaBOddColumn sigmaBOddSeedQ c} the
explicit rational column, built from a rational ``seed'' vector \lean{sigmaBOddSeedQ} with a single nonzero entry
$\pm1$, AxiomMath's rational spin matrix \lean{Rmat c} and a sign.  (c) is stated twice: as rationality, and as
\lean{∃ g : liftableSubgroup, monsterLift g = sigmaAut}.  The element $g$ is called $\sigma_B$ below (Lean:
\lean{tauBLiftSigma}, chosen by \lean{Classical.choose}; it is unique by injectivity of \lean{monsterLift}).
\end{dictionary}
\begin{proof}[Proof on paper]
(a) The even untwisted block is rational by Theorem~\ref{ch22:thm:sigmaeven}.  The odd untwisted block goes to $T_s^+$, and
its rationality is the hypothesis.  The rest follows from structure: $\sigma$ is an involution, so rational images on
$V_3^+$ give rational preimages on $T_s^+$ (\lean{sigmaBTs_rational}); and $\sigma$ preserves the Griess product, while
the product with one odd short class is a rational bijection exchanging the two halves $T_s^+$ and $T_{s+1}^+$ of the
twisted weight-two space (\lean{sigmaBShortAction_frame}).  It is invertible (\lean{sigmaBShortAction_injective}): its oscillator factor is built from
AxiomMath's rank-one matrix $\widehat U_c=u_cu_c^{\mathsf T}$ (\lean{Umat c}, $u_c$ the Conway vector of the short class, $u_c\cdot u_c=32$,
so $\widehat U_c^2=32\,\widehat U_c$, \lean{sigmaBUmat_sq}), and the Lean docstring records that the factor has eigenvalues $1$ and $-7$, neither
zero.  Hence the $\Theta$ block is rational without computing $\Theta$ (\lean{sigmaBTs1_rational_of_phi}).
(b) By the top-level formula for $\Phi$ (Chapter~\ref{ch:13}), $\zeta'_c\Phi(x_{\alpha_c})$ is a single oscillator column
$h_{f_b}(-\tfrac12)\otimes(\text{sign})\,R_c\,v$ with $v$ one fixed ``seed'' spinor.  The minus-root characters (diagonal
cocode operators) and the frame sector force $v$ to be supported on one spin index
(\lean{sigmaBOddSeed_support}); the invariant form and the isometry property of $\sigma$ give $\Phi$-images of norm
$2$, which forces the square of that one coordinate to be $1$ (\lean{sigmaBOddSeed_form}).  So $v$ is rational.
(c) The rationality criterion \lean{isVnatLiftable_of_rational} (\S\ref{ch22:sec:recall}) applied to $h=\sigma$.
\end{proof}
\begin{physics}
The triality symmetry of the moonshine CFT, which mixes untwisted vertex operators with twisted-sector fields, preserves
AxiomMath's rational form of the $196884$ weight-two operators: in the rational basis its matrix has entries in
$\Z[\tfrac12]$.  Physically this is the statement that the triality-transformed twisted ground state is again a ``real''
state of the same lattice of normalisations.
\end{physics}
\begin{status}Unconditional; all four declarations audited, standard axioms.\end{status}

%======================================================================
\section{The triality element \texorpdfstring{$\tau_B$}{tauB} and Conway--Norton for \texorpdfstring{$G_{\mathrm{mat}}$}{Gmat}}\label{ch22:sec:tauB}
%======================================================================

$\sigma_B$ is a Monster element of order $2$ that exchanges the orbifold descriptions; AxiomMath's $\tau_B$ is an element of
order $3$ with the same effect, written as an explicit matrix.  They are \emph{different} elements, and no attempt is made
to compute a word expressing one through the other.  Instead:

\begin{theorem}[$\tau_B$ is liftable]\label{ch22:thm:tauB}
$\tau_B$, as an element of $\chtwentytwoAutO$, is liftable.  More precisely:
(a) $\sigma_Bx_B\sigma_B^{-1}\in\{y_B,z_B\}$;
(b) if $g\in\chtwentytwoAutO$ satisfies $gx_Bg^{-1}=\sigma_Bx_B\sigma_B^{-1}$, then $g$ is liftable;
(c) hence $\tau_B$ is liftable.
\end{theorem}
\begin{leancode}
def tauBLiftConjugate : AutO :=
  (tauBLiftSigma : AutO) * ⟨xB, xB_mem_AutO⟩ * (tauBLiftSigma : AutO)⁻¹
-- ⋯
theorem tauBLiftConjugate_dict (u : SpaceB) :
    (griessNat ((tauBLiftConjugate : SpaceB ≃ₗ[ℚ] SpaceB) u) : Vnat) =
      sigmaBlockFrameSign Dnat.s (griessNat u : Vnat)
-- ⋯
theorem tauBLiftSigma_conj_xB (hτ : tauB ∈ AutO) :
    (tauBLiftConjugate : SpaceB ≃ₗ[ℚ] SpaceB) = yB ∨
      (tauBLiftConjugate : SpaceB ≃ₗ[ℚ] SpaceB) = zB
-- ⋯
theorem tauBLift_of_conjugate (g : AutO)
    (hg : (g : SpaceB ≃ₗ[ℚ] SpaceB) * xB * (g : SpaceB ≃ₗ[ℚ] SpaceB)⁻¹ =
      (tauBLiftConjugate : SpaceB ≃ₗ[ℚ] SpaceB)) : IsVnatLiftable g
-- ⋯
theorem tauBLift_isVnatLiftable_tauB_of_mem_AutO (hτ : tauB ∈ AutO) :
    IsVnatLiftable ⟨tauB, hτ⟩
\end{leancode}
(\leanpath{VOA/Moonshine/MoonshineLRTauB.lean}.)
\begin{leancode}
theorem isVnatLiftable_tauB : IsVnatLiftable ⟨tauB, tauB_mem_AutO⟩ :=
  tauBLift_isVnatLiftable_tauB_of_mem_AutO tauB_mem_AutO
\end{leancode}
(\leanpath{VOA/Moonshine/MoonshineLRMonsterFinal.lean}; full name \chtwentytwolean{VOA.Moonshine.ConwayNorton.M8.isVnatLiftable_tauB}.)
\begin{dictionary}
\lean{tauBLiftSigma : liftableSubgroup} is $\sigma_B$; \lean{tauBLiftConjugate} $=\sigma_Bx_B\sigma_B^{-1}$, computed in
\lean{AutO} (where $x_B\in\chtwentytwoAutO$ is \lean{xB_mem_AutO}).  \lean{hτ : tauB ∈ AutO} is a hypothesis of the
intermediate theorems; it is discharged by AxiomMath's (ported) \lean{tauB_mem_AutO}, which is how
\lean{isVnatLiftable_tauB} has no hypotheses.  \lean{⟨tauB, tauB_mem_AutO⟩} is $\tau_B$ as an element of
\lean{AutO}.
\end{dictionary}
\begin{proof}[Proof on paper]
(a) \emph{Three ingredients from AxiomMath's coordinates} (Chapter~\ref{ch:21}).
\begin{itemize}
\item \emph{Spin index and row halves.}  The rows of a matrix $N\in\chtwentytwoW$ are indexed by
  $\mathrm{SpinIdx}=\Z/2\times C_0$, pairs $(\eta,a)$ with $\eta\in\Z/2$ and $a$ in a fixed $2^{11}$-element subspace $C_0$ of
  the Golay code (so $|\mathrm{SpinIdx}|=2\cdot2048=4096$).  The \emph{row half}
  $\eta$ of a row is its first component.  Through the dictionary the two row halves $\eta=0,1$ go to the two twisted frame
  sectors $T_0^+$, $T_1^+$, up to a fixed shift (\lean{spinDict_mem_TPframe_shift}).
\item \emph{The triality basis.}  \lean{tbasis} is AxiomMath's basis of $\chtwentytwoB$ adapted to $\tau_B$: the $24$ diagonal
  matrix units (fixed by $\tau_B$), $276$ ``pair'' triples, $2048\cdot24$ ``odd'' triples, and $759$ octad blocks of size $64$ ($24+3\cdot276+3\cdot2048\cdot24+64\cdot759=196884$).
  The odd triple of $(a,i)\in C_0\times\{0..23\}$ consists of $k=0$: an odd short line of $\chtwentytwoX$; $k=1,2$: the spin
  vectors $\pm E(\eta,a)\otimes e_i\in\chtwentytwoW$ with $\eta=1$ and $\eta=0$ respectively; $\tau_B$ cycles $k\mapsto k+1$.  So
  all of $\chtwentytwoW$ lies in the odd triples.
\item \emph{Sign patterns.}  For $s\in\{\pm1\}^3$, \lean{oddSign s} is the diagonal operator that multiplies the $k$-th vector of
  every odd triple by $s_k$ and fixes every other basis vector.  $x_B$ is the pattern $(+,-,-)$ ($+1$ on $\chtwentytwoA\oplus
  \chtwentytwoX$, $-1$ on $\chtwentytwoW$).  Since $\tau_B$ shifts $k$, $y_B=\tau_Bx_B\tau_B^{-1}$ is the pattern $(-,+,-)$ and
  $z_B=\tau_B^2x_B\tau_B^{-2}$ is $(-,-,+)$.
\end{itemize}
\emph{The argument.}  Through the dictionary $x_B$ is $z$ (\lean{tauBLiftParity_dict}) and $\sigma_B$ is $\sigma$, so
$\sigma_Bx_B\sigma_B^{-1}$ is the frame sign $z'$ of Theorem~\ref{ch22:thm:sector}(b) (\lean{tauBLiftConjugate_dict}).  In
AxiomMath's coordinates $z'$ fixes $\chtwentytwoA$ and the even short lines (they lie in $V_0^+$), negates the odd short lines
($V_3^+$), and on $\chtwentytwoW$ is $+1$ on one row half and $-1$ on the other (\lean{tauBLiftConjugate_W}).  On the odd
triples it is therefore the pattern $(-,+,-)$ or $(-,-,+)$, according to which row half is positive.  A parity automorphism
(one commuting with $x_B$) is determined by its action on $\chtwentytwoA$ and on the spin vectors $k=1,2$
(\lean{tauBLiftParity_ext}); $z'$, $y_B$, $z_B$ all fix $\chtwentytwoA$, so $z'=y_B$ or $z'=z_B$
(\lean{tauBLiftConjugate_eq_of_signs}).  The hypothesis $\tau_B\in\chtwentytwoAutO$ is used to know that $y_B,z_B$ lie in
$\chtwentytwoAutO$ and commute with $x_B$.

(b) Put $r=\sigma_B^{-1}g$.  Then $rx_Br^{-1}=x_B$, so $r\in\chtwentytwoAutO\cap C(x_B)=\chtwentytwoGvis$ by the
parity-centralizer theorem \eqref{ch22:eq:parity}, so $r$ is liftable (Theorem~\ref{ch22:thm:visible}); $\sigma_B$ is liftable;
the liftable elements form a group; so $g=\sigma_Br$ is liftable.

(c) If $\sigma_Bx_B\sigma_B^{-1}=y_B=\tau_Bx_B\tau_B^{-1}$, apply (b) to $g=\tau_B$.  If it is $z_B=\tau_B^2x_B\tau_B^{-2}$, then
since $\tau_B^3=1$ we have $\tau_B^{-1}=\tau_B^2$ and $\tau_B^{-1}x_B\tau_B=z_B$; apply (b) to $g=\tau_B^{-1}$, and invert.
\end{proof}
\begin{physics}
The four-group $\{1,z,z',zz'\}$ consists of three ``2B'' involutions, each the quantum symmetry of an orbifold
description.  Any symmetry that permutes them in the same way as $\tau_B$ differs from $\tau_B$ by a symmetry commuting with
$z$, i.e.\ by an element of the orbifold-visible group $2^{1+24}.\chtwentytwoCo_1$, which is already known to act.  So
it suffices to exhibit one CFT symmetry ($\sigma$) that moves $z$ to $z'$ or $zz'$.
\end{physics}
\begin{status}
\lean{isVnatLiftable_tauB}: unconditional.  The intermediate theorems carry the hypothesis \lean{hτ : tauB ∈ AutO}, which is a
theorem of the port.  All audited (\leanpath{VOA/Moonshine/MoonshineLRLiftAudit.lean}; \lean{isVnatLiftable_tauB} also in
\leanpath{VOA/Moonshine/MoonshineCNMonsterFinalAudit.lean}), standard axioms.
\end{status}


\subsection{Conway--Norton for the generated group and for the Monster}

\begin{theorem}[Every element of $\chtwentytwoGmat$ acts on $\Vnat$; Conway--Norton for $\chtwentytwoGmat$]\label{ch22:thm:gmat}
(a) $\chtwentytwoGmat$ is the image in $\mathrm{GL}_\Q(\chtwentytwoB)$ of the subgroup of $\chtwentytwoAutO$ generated by $\tau_B$ and
the visible elements.
(b) Every $g\in\chtwentytwoAutO$ with $g\in\chtwentytwoGmat$ is liftable.
(c) For every $g\in\chtwentytwoAutO$ with $g\in\chtwentytwoGmat$ there is an automorphism $h$ of $\Vnat$ acting on
$\Vnat_2\cong\C\otimes\chtwentytwoB$ as $1\otimes g$, whose McKay--Thompson function $T_h$ is the normalised Hauptmodul of a
genus-zero group $\Gamma\supseteq\Gamma_0(N)$ commensurable with $\mathrm{SL}_2(\Z)$.
\end{theorem}
\begin{leancode}
noncomputable def tauBAutO : AutO := ⟨tauB, tauB_mem_AutO⟩
-- ⋯
theorem gmat_eq_map_closure_insert_tauB :
    Generated.Gmat = (Subgroup.closure (insert tauBAutO M8.visibleSetAutO)).map AutO.subtype
-- ⋯
theorem isVnatLiftable_of_mem_gmat (g : AutO)
    (hg : (g : SpaceB ≃ₗ[ℚ] SpaceB) ∈ Generated.Gmat) : IsVnatLiftable g
-- ⋯
theorem conwayNorton_gmat (g : AutO) (hg : (g : SpaceB ≃ₗ[ℚ] SpaceB) ∈ Generated.Gmat) :
    ∃ h : AutVnat, vnatWeightTwoRestrict h = griessConj (g : SpaceB ≃ₗ[ℚ] SpaceB) ∧
      ∃ (Γ : Subgroup (GL (Fin 2) ℝ)) (N : ℕ), IsMoonshineGroup Γ N ∧
        IsNormalizedHauptmodul Γ (mckayThompsonFun h)
\end{leancode}
and its unfolded form, an anonymous \lean{example} in the same file:
\begin{leancode}
example (g : SpaceB ≃ₗ[ℚ] SpaceB)
    (hA : (∀ u v, g (gmul u v) = gmul (g u) (g v)) ∧ ∀ u v, bform (g u) (g v) = bform u v)
    (hg : g ∈ Subgroup.closure ((Visible.Gvis : Set (SpaceB ≃ₗ[ℚ] SpaceB)) ∪ {tauB})) :
    ∃ h : AutVnat, vnatWeightTwoRestrict h
        = (griessIso.symm ≪≫ₗ LinearEquiv.baseChange ℚ ℂ SpaceB SpaceB g ≪≫ₗ griessIso) ∧
      ∃ (Γ : Subgroup (GL (Fin 2) ℝ)) (N : ℕ), IsMoonshineGroup Γ N ∧
        IsNormalizedHauptmodul Γ (mckayThompsonFun h) :=
\end{leancode}
(\leanpath{VOA/Moonshine/MoonshineCNMonsterFinalAudit.lean}; full names \chtwentytwolean{VOA.Moonshine.ConwayNorton.conwayNorton_gmat},
etc.)
\begin{dictionary}
\textbf{Decoding.}  \lean{AutO.subtype} is the inclusion $\chtwentytwoAutO\hookrightarrow\mathrm{GL}_\Q(\chtwentytwoB)$ and
\lean{.map} the image of a subgroup.  \lean{(g : SpaceB ≃ₗ[ℚ] SpaceB) ∈ Generated.Gmat} is membership of the underlying linear
map.  In the \lean{example}, every defined term is unfolded: $g$ is any $\Q$-linear automorphism of $\chtwentytwoB$
preserving \lean{gmul} and \lean{bform} (\lean{hA}: membership in $\chtwentytwoAutO$) and lying in the subgroup generated by
$\chtwentytwoGvis\cup\{\tau_B\}$ (\lean{hg}); the conclusion's first clause writes \lean{griessConj g} out as
\lean{griessIso.symm ≪≫ₗ LinearEquiv.baseChange ℚ ℂ SpaceB SpaceB g ≪≫ₗ griessIso}, i.e.\ $\mathcal G\circ(1_\C\otimes
g)\circ\mathcal G^{-1}$ with $\mathcal G=$ \lean{griessIso}.

\textbf{Hypotheses.}  \lean{conwayNorton_gmat}: only the membership \lean{hg}, which is the intended restriction to
$\chtwentytwoGmat$.  No \lean{GmatEqAutO}.

\textbf{Same statement, and not vacuous.}  The conclusion has two parts.  The first (existence of a lift $h$ with prescribed
action on $\Vnat_2$) is the new content of this chapter; the second is \lean{conwayNorton h} from Chapter~\ref{ch:20}.  The
statement is not vacuous: $\chtwentytwoGmat$ contains $x_B$, $\tau_B$, and all of $\chtwentytwoGvis$ (a group of order
$2^{25}\cdot|\chtwentytwoCo_1|$ in the classical identification), and for each such $g$ the theorem produces a specific
automorphism of $\Vnat$, unique by \lean{vnatWeightTwoRestrict_injective}.  Since the first clause pins $h|_{\Vnat_2}$, the
traces of $h$ on $\Vnat_2$ are those of $g$ on $\chtwentytwoB$; e.g.\ for $g=\tau_B$, $\tr(h|\Vnat_2)=\tr\tau_B=783$
(\lean{trace_tauB}); classically (not a Lean statement) this is $1+782$, the value on the Monster class 3A of the
character $1+\chi_{196883}$.
\end{dictionary}
\begin{proof}[Proof on paper]
(a) $\chtwentytwoGmat$ is by definition the subgroup generated by $\chtwentytwoGvis\cup\{\tau_B\}$; all generators lie in
$\chtwentytwoAutO$, and images of closures are closures of images.  (b) The generators of that subgroup of
$\chtwentytwoAutO$ are liftable (Theorems~\ref{ch22:thm:visible} and~\ref{ch22:thm:tauB}), and liftable elements form a
subgroup (\lean{closure_insert_tauB_le_liftableSubgroup}).  (c) For liftable $g$, \lean{conwayNorton_monster_of_isVnatLiftable}
(Chapter~\ref{ch:21}) produces $h=\mathrm{monsterLift}(g)$ and applies \lean{conwayNorton}.
\end{proof}
\begin{physics}
Every element of the group generated by the orbifold-visible symmetries and triality is a symmetry of the moonshine CFT,
and its twisted partition function is a genus-zero Hauptmodul.  If one already knows (from group theory) that this
generated group is the Monster, this is monstrous moonshine; the theorem itself does not need to know it.
\end{physics}
\begin{status}
\lean{conwayNorton_gmat}, \lean{isVnatLiftable_of_mem_gmat}, \lean{gmat_eq_map_closure_insert_tauB}: \textbf{unconditional} (no
hypothesis beyond the membership premise $g\in\chtwentytwoGmat$, which is part of the statement).
\texttt{\#print axioms} in \leanpath{VOA/Moonshine/MoonshineCNMonsterFinalAudit.lean} (build record of the release, 9,441 jobs):
\lean{[propext, Classical.choice, Quot.sound]}, $0$ \lean{sorryAx}.
\end{status}

\begin{theorem}[Conway--Norton for all of $\chtwentytwoAutO$, given $\chtwentytwoGmat=\chtwentytwoAutO$]\label{ch22:thm:cnmonster}
If $\chtwentytwoGmat=\chtwentytwoAutO$, then for every $g\in\chtwentytwoAutO$ the conclusion of
Theorem~\ref{ch22:thm:gmat}(c) holds.
\end{theorem}
\begin{leancode}
theorem conwayNorton_monster_of_Gmat (hGmat : Generated.Gmat = AutO) (g : AutO) :
    ∃ h : AutVnat, vnatWeightTwoRestrict h = griessConj (g : SpaceB ≃ₗ[ℚ] SpaceB) ∧
      ∃ (Γ : Subgroup (GL (Fin 2) ℝ)) (N : ℕ), IsMoonshineGroup Γ N ∧
        IsNormalizedHauptmodul Γ (mckayThompsonFun h)
\end{leancode}
\begin{leancode}
theorem conwayNorton_monster (hGmat : Generated.Gmat = AutO) (g : AutO) :
    ∃ h : AutVnat, vnatWeightTwoRestrict h = griessConj (g : SpaceB ≃ₗ[ℚ] SpaceB) ∧
      ∃ (Γ : Subgroup (GL (Fin 2) ℝ)) (N : ℕ), IsMoonshineGroup Γ N ∧
        IsNormalizedHauptmodul Γ (mckayThompsonFun h)
\end{leancode}
(\leanpath{VOA/Moonshine/MoonshineLRMonsterFinal.lean}, \leanpath{VOA/Moonshine/MoonshineCNMonsterFinalAudit.lean}; the
second restates the first in the namespace \chtwentytwolean{VOA.Moonshine.ConwayNorton}.)
\begin{dictionary}
The hypothesis \lean{hGmat : Generated.Gmat = AutO} is \emph{literally} the body of \lean{GmatEqAutO}
(\S\ref{ch22:sec:headline}); the two theorems predate that name.  Equality of subgroups of $\mathrm{GL}_\Q(\chtwentytwoB)$.
\end{dictionary}
\begin{proof}[Proof on paper]
Under \lean{hGmat}, $\tau_B$ and the visible elements generate $\chtwentytwoAutO$ (\lean{tauBLift_closure_eq_top}), and
\lean{conwayNorton_autO_of_tauB} with \lean{isVnatLiftable_tauB} gives the claim.
\end{proof}
\begin{physics}Monstrous moonshine for every element of AxiomMath's Monster, modulo the group-theoretic identification
$\chtwentytwoGmat=\chtwentytwoAutO$.\end{physics}
\begin{status}\textbf{Conditional} on \lean{hGmat}.  Audited in \leanpath{VOA/Moonshine/MoonshineCNMonsterFinalAudit.lean},
standard axioms (which do not report \lean{hGmat}).\end{status}

%======================================================================
\section{The converse, I: every graded trace is rational}\label{ch22:sec:traces}
%======================================================================

\subsection{Why the converse needs a new idea}

So far: $\chtwentytwoGmat\le\mathrm{liftable}\le\chtwentytwoAutO$, and $\mathrm{monsterLift}$ embeds the liftable subgroup
in $\Aut(\Vnat)$.  The converse asks whether \emph{every} automorphism of $\Vnat$ arises this way, i.e.\ preserves the
rational lattice $\mathrm{griessNat}(\chtwentytwoB)\subset\Vnat_2$ (\lean{mem_range_monsterLift_iff},
Chapter~\ref{ch:21}).  Classically one argues: $\Aut(\Vnat)$ acts faithfully on the Griess algebra $\Vnat_2$ by algebra
automorphisms, and by Tits' theorem the automorphism group of the complex Griess algebra is $\chtwentytwoM$.  Tits'
theorem is not formalised.  The route taken instead (``route (e)'' of the authors' notes) is:

\begin{quote}
\emph{Let $h$ be an automorphism of $\Vnat$.  If the traces on $\Vnat_2$ of all elements of the coset
$h\cdot\chtwentytwoGmat$ (more precisely of $\mathrm{gmatLift}(s)^{-1}h$, $s\in\chtwentytwoGmat$) are rational, and the already
rational group $\chtwentytwoGmat$ is absolutely irreducible on $\mathbf \mathbf{e}_1^\perp$, then the matrix of $h$ on $\Vnat_2$ is
rational.}
\end{quote}
Note that the hypothesis is about the traces of \emph{many} automorphisms, not only the character value of $h$ itself;
this is why rationality of the traces of \emph{every} automorphism of $\Vnat$ is proved first
(Theorem~\ref{ch22:thm:rationaltraces}).

The arithmetic input (rational characters) comes from moonshine itself; the algebraic input (absolute irreducibility) is
the commutant lemma of \S\ref{ch22:sec:commutant}; Schur's orthogonality relations combine them.

\begin{warning}
\textbf{Recorded dead ends.}
(1) \emph{Galois descent on $\Aut(\Vnat)$ alone} gives only $N_G(M)=M$ for the rational subgroup $M$ (from $C_G(M)=1$); it
cannot exclude a finite overgroup in which $M$ is self-normalising.  The missing ingredient is rationality of characters.
(2) \emph{``Rational character $\Rightarrow$ rational matrices'' is false in general}: the quaternion group $Q_8$ acting on
$\C^2$ has rational character but no rational model (Schur index $2$).  The argument needs an \emph{already rational,
absolutely irreducible} subgroup, which is exactly the commutant lemma; it cannot be dropped.
(3) \emph{Counting} $|\Aut(\Vnat)|\le|\chtwentytwoGmat|$ needs $|\chtwentytwoGmat|$, which AxiomMath obtains only through the
Baby Monster recognition (the $\approx16{,}000$-file closure).  Dead.
(4) \emph{Realness first, then real $\Rightarrow$ rational}: the second step is Tits-hard; superseded.
(5) \emph{Classifying rigid idempotents / Ising vectors}: the known proofs that every Ising vector of $\Vnat$ is a 2A-axis go
through $\Aut(\Vnat)=\chtwentytwoM$ (circular here).  Unnecessary.
\end{warning}

\subsection{Galois invariance and coefficients}

The following three theorems are in \leanpath{VOA/Moonshine/MoonshineCNRationalTraces.lean} (namespace
\chtwentytwolean{VOA.Moonshine.ConwayNorton}); they assemble inputs proved in Chapters~\ref{ch:16}, \ref{ch:19} and \ref{ch:20}
into unconditional statements.

\begin{theorem}[Complete replicability, the $\Gamma_1$ theorem, Galois invariance]\label{ch22:thm:galois}
(a) Every McKay--Thompson series of $\Vnat$ is completely replicable.  (b) The $\Gamma_1$ form of Conway--Norton holds.
(c) For every $g\in\Aut(\Vnat)$ and every $u$ coprime to the order of $g$, $T_{g^u}=T_g$ as functions on $\mathbb H$.
\end{theorem}
\begin{leancode}
theorem completelyReplicable_holds : CompletelyReplicable :=
  completelyReplicable_of_simpleRootGeneration M8.monsterLieSimpleRootGeneration
    M8.autVnat_isOfFinOrder
-- ⋯
theorem conwayNortonGammaOne_holds : ConwayNortonGammaOne :=
  conwayNortonGammaOne_of_completelyReplicable completelyReplicable_holds
    (fun _ _ k => M8.ReplicableFamily_primeHeckeIdentity _ k)
    (fun _ hg => M8.ReplicableFamily_notInjectiveModTranslations hg.orderOf_pos _)
-- ⋯
theorem mckayThompsonFun_pow_eq_of_coprime_holds (g : AutVnat) {u : ℕ}
    (hu : Nat.Coprime u (orderOf g)) : mckayThompsonFun (g ^ u) = mckayThompsonFun g
\end{leancode}
(\leanpath{VOA/Moonshine/MoonshineCNRationalTraces.lean}.)
\begin{dictionary}
\lean{CompletelyReplicable} and \lean{ConwayNortonGammaOne} are the definitions of Chapter~\ref{ch:16}
(\leanpath{VOA/Moonshine/MoonshineCNStatement.lean}).  The proof terms are shown because they are informative: (a) is
Chapter~\ref{ch:19}'s \lean{completelyReplicable_of_simpleRootGeneration} fed with the generation theorem
\lean{M8.monsterLieSimpleRootGeneration} and finite order \lean{M8.autVnat_isOfFinOrder} (Chapter~\ref{ch:20}); (b) is the
case-free genus-zero theorem fed with (a) and the two unconditional inputs on replicable families (Hecke identity at primes,
non-injectivity modulo translations); (c) is the Galois-invariance theorem of Chapter~\ref{ch:16}
(\lean{mckayThompsonFun_pow_eq_of_coprime_of_completelyReplicable}) with all its hypotheses discharged.
\lean{Nat.Coprime u (orderOf g)}: $\gcd(u,o(g))=1$; \lean{g ^ u} the $u$-th power in $\Aut(\Vnat)$.
\end{dictionary}
\begin{proof}[Proof on paper]
Assembly.  For (c), recall the argument of Chapter~\ref{ch:16} (Galois invariance): let $(\Gamma,M)$ come from the
$\Gamma_1$ form for $g$; by Dirichlet's theorem there is a prime $p>M$ with $p\equiv u\pmod{o(g)}$, so $g^p=g^u$; replication
at $p$ and the main lemma of the $\Gamma_1\Rightarrow\Gamma_0$ section of Chapter~\ref{ch:16} give $T_{g^p}=T_g$.
\end{proof}
\begin{physics}
$g$ and $g^u$ generate the same cyclic group; their twisted partition functions are related by a Galois conjugation of
the $o(g)$-th roots of unity in the characters, and the theorem says the partition functions coincide.  In a CFT where the
twisted sectors are governed by genus-zero modular curves, this is the expected ``rationality'' of the twisted partition
functions.
\end{physics}
\begin{status}Unconditional.  \texttt{\#print axioms} in \leanpath{VOA/Moonshine/MoonshineCNRationalAudit.lean}: standard
axioms.\end{status}

\begin{theorem}[Coefficients are determined by the function]\label{ch22:thm:coeff}
If $g,h\in\Aut(\Vnat)$ have finite order and $T_g=T_h$ on $\mathbb H$, then their McKay--Thompson power series are equal.
\end{theorem}
\begin{leancode}
theorem mckayThompson_eq_of_mckayThompsonFun_eq {g h : AutVnat} (hg : IsOfFinOrder g)
    (hh : IsOfFinOrder h) (hfun : mckayThompsonFun g = mckayThompsonFun h) :
    mckayThompson g = mckayThompson h
\end{leancode}
(\leanpath{VOA/Moonshine/MoonshineCNRationalTraces.lean}.)
\begin{dictionary}
\lean{mckayThompson g} is the power series $\sum_{n\ge0}\tr(g|\Vnat_{n+1})q^n$ (Chapter~\ref{ch:16}); \lean{IsOfFinOrder g}
is $g^n=1$ for some $n\ge1$ (always true here, \lean{autVnat_isOfFinOrder}, but stated as a hypothesis in this lemma).
\end{dictionary}
\begin{proof}[Proof on paper]
The difference of the regular parts is a power series in $q$ whose coefficients are bounded by those of $J$
(\lean{isJDominated_mckayThompson}), hence converges on the punctured unit disc, and it sums to $0$ at every
$q=e^{2\pi i\tau}$; the Cauchy estimate \lean{M8.norm_coeff_le_of_hasSum_qPar} with bound $0$ kills every coefficient.
\end{proof}
\begin{physics}A $q$-expansion is determined by the function it sums to: the twisted partition function determines the
twisted degeneracies at every level.\end{physics}
\begin{status}Unconditional; audited in \leanpath{VOA/Moonshine/MoonshineCNRationalAudit.lean}, standard axioms.\end{status}

\subsection{Ramanujan sums}

\begin{theorem}[Galois-invariant traces of finite-order operators are rational]\label{ch22:thm:ramanujan}
Let $K$ be a field of characteristic $0$, $V$ a finite-dimensional $K$-vector space, $A\in\End_K V$ with $A^N=1$ for some
$N\ge1$, and suppose $\tr(A^u)=\tr(A)$ for every $u$ coprime to $N$.  Then $\tr A\in\Q$.
\end{theorem}
\begin{leancode}
theorem exists_rat_trace_of_trace_pow_coprime [CharZero K] (A : Module.End K V) {N : ℕ}
    (hN : 0 < N) (hA : A ^ N = 1)
    (hgal : ∀ u : ℕ, N.Coprime u → LinearMap.trace K V (A ^ u) = LinearMap.trace K V A) :
    ∃ r : ℚ, LinearMap.trace K V A = r
\end{leancode}
(\leanpath{VOA/Moonshine/MoonshineCNRationalRamanujan.lean}; section variables
\lean{{K V : Type*} [Field K] [AddCommGroup V] [Module K V] [FiniteDimensional K V]}.)
\begin{dictionary}
\lean{N.Coprime u} is $\gcd(N,u)=1$; the conclusion \lean{∃ r : ℚ, LinearMap.trace K V A = r} says the trace is the image of a
rational number under $\Q\to K$.  Hypotheses exactly as stated; note $u$ ranges over \emph{all} naturals coprime to $N$, so
(\lean{hgal}) is the usual Galois-invariance hypothesis.
\end{dictionary}
\begin{proof}[Proof on paper]
No eigenvalues are needed.  Möbius inversion of the coprimality indicator gives
\begin{equation*}
  \sum_{\substack{0\le u<N\\(u,N)=1}}A^u=\sum_{e\mid N}\mu(e)\sum_{k<N/e}(A^e)^k
\end{equation*}
(\lean{sum_coprime_eq_sum_divisors_moebius}).  For $B$ with $B^M=1$, $P=\sum_{k<M}B^k$ satisfies $P^2=MP$, so $P/M$ is
an idempotent and $\tr P=M\cdot\mathrm{rank}(P/M)\in M\N$ (\lean{exists_trace_geom_sum_eq_of_pow_eq_one}).  Hence
$\varphi(N)\tr A=\sum_{(u,N)=1}\tr(A^u)\in\Z$, and $\tr A\in\Q$.  (Read on an eigenvector, this is the classical fact that
Ramanujan sums $\sum_{(u,N)=1}\lambda^u$ are integers.)
\end{proof}
\begin{physics}Pure linear algebra; the physical content is in the hypothesis.\end{physics}
\begin{status}Unconditional; audited in \leanpath{VOA/Moonshine/MoonshineCNRationalAudit.lean}, standard axioms.\end{status}

\begin{theorem}[Rational traces]\label{ch22:thm:rationaltraces}
For every $g\in\Aut(\Vnat)$ and $n\in\N$, $\tr(g|\Vnat_{n+1})\in\Q$.  Equivalently, every coefficient of every
McKay--Thompson series of $\Vnat$ is rational.
\end{theorem}
\begin{leancode}
theorem exists_rat_gradedAutTrace (g : AutVnat) (n : ℕ) :
    ∃ r : ℚ, gradedAutTrace VnatVOA g ((n : ℤ) + 1) = r
-- ⋯
theorem rationalTraces : RationalTraces
\end{leancode}
(\leanpath{VOA/Moonshine/MoonshineCNRationalTraces.lean}; full name \chtwentytwolean{VOA.Moonshine.ConwayNorton.rationalTraces}.)
\begin{dictionary}
\lean{gradedAutTrace VnatVOA g k} is $\tr(g|\Vnat_k)$ (Chapter~\ref{ch:16}); \lean{((n : ℤ) + 1)} is the weight $n+1$ as
an integer.  \lean{RationalTraces} (quoted in \S\ref{ch22:sec:recall}) quantifies over all $g$ and all coefficients $n$ of
\lean{mckayThompson g}, whose $n$-th coefficient is $\tr(g|\Vnat_{n+1})$ (\lean{mckayThompson_coeff}).  No hypotheses.
\textbf{Strength.}  Rationality, not integrality: integrality of the coefficients is classical (they are virtual characters
evaluated at $g$) but is not a statement of the library (Chapter~\ref{ch:16}).
\end{dictionary}
\begin{proof}[Proof on paper]
Fix $g$ of order $N$ (finite by Chapter~\ref{ch:20}) and $A=g|_{\Vnat_{n+1}}$, so $A^N=1$.  For $u$ coprime to $N$,
Theorem~\ref{ch22:thm:galois}(c) gives $T_{g^u}=T_g$, Theorem~\ref{ch22:thm:coeff} gives equality of coefficients, i.e.\
$\tr(A^u)=\tr(A)$.  Theorem~\ref{ch22:thm:ramanujan} concludes.
\end{proof}
\begin{physics}
\textbf{Why rational characters matter.}  For a finite-order symmetry $g$ of a theory with a rational structure, the trace on
each energy level is a sum of roots of unity; Galois invariance of the twisted partition function forces it to be rational.
Rational characters are the first indication that the symmetry group can be ``defined over $\Q$''; combined with an
absolutely irreducible rational subgroup (next sections), they pin the group down to the rational one.
\end{physics}
\begin{status}\textbf{Unconditional.}  \texttt{\#print axioms}: standard, in \leanpath{VOA/Moonshine/MoonshineCNRationalAudit.lean}
and \leanpath{VOA/Moonshine/MoonshineCNSymmetryFinalAudit.lean}.\end{status}

%======================================================================
\section{The converse, II: Schur's span formula and rational descent}\label{ch22:sec:schur}
%======================================================================

\begin{theorem}[Schur's span formula]\label{ch22:thm:schurspan}
Let $K$ be a field, $S$ a finite group, $\sigma:S\to M_n(K)$ a multiplicative map, $E\in M_n(K)$ an idempotent commuting with
$\sigma(S)$ with $\tr E\ne0$, and suppose every $T$ commuting with $\sigma(S)$ satisfies $ETE\in K\cdot E$.  Then for every
$Y\in M_n(K)$,
\begin{equation*}
  \sum_{s\in S}\tr\bigl(\sigma(s^{-1})EY\bigr)\,\sigma(s)E=\frac{|S|}{\tr E}\,EYE .
\end{equation*}
\end{theorem}
\begin{leancode}
theorem schur_span_formula (σ : S →* Matrix n n K) (E : Matrix n n K) (hE : E * E = E)
    (hEσ : ∀ s, E * σ s = σ s * E) (hr : E.trace ≠ 0)
    (hirr : ∀ T : Matrix n n K, (∀ s, T * σ s = σ s * T) → ∃ c : K, E * T * E = c • E)
    (Y : Matrix n n K) :
    ∑ s, (σ s⁻¹ * E * Y).trace • (σ s * E)
      = ((Fintype.card S : K) / E.trace) • (E * Y * E)
\end{leancode}
(\leanpath{VOA/Moonshine/MoonshineCNRationalSchur.lean}; section variables
\lean{{K n S : Type*} [Field K] [Fintype n] [DecidableEq n] [Group S] [Fintype S]}.)
\begin{dictionary}
\lean{S →* Matrix n n K} is a monoid homomorphism into the matrix ring (so $\sigma(1)=1$ and $\sigma(st)=\sigma(s)\sigma(t)$;
invertibility follows); \lean{σ s⁻¹} is $\sigma(s^{-1})$, which equals $\sigma(s)^{-1}$.  \lean{hirr} is the commutant
hypothesis (``absolute irreducibility on the image of $E$''); \lean{Fintype.card S} is $|S|$.  No hypothesis on the characteristic is needed; the theorem is used only over $\Q$.
\end{dictionary}
\begin{proof}[Proof on paper]
For any $X$, the group average $A(X)=\sum_s\sigma(s)EXE\sigma(s)^{-1}$ commutes with $\sigma(S)$ and satisfies $A=EAE$, so
by the hypothesis $A(X)=c(X)E$; taking traces, $c(X)\tr E=|S|\tr(EX)$ (\lean{groupAverage_eq_smul}).  Applied to matrix
units, this is the orthogonality relation $\sum_s(\sigma(s)E)_{ij}(E\sigma(s^{-1}))_{kl}=\frac{|S|}{\tr E}E_{kj}E_{il}$
(\lean{sum_mul_coeff_eq}), and summing it against the entries of $Y$ gives the formula.
\end{proof}
\begin{physics}The Schur orthogonality relations for an irreducible representation: the matrix coefficients of the group
span the whole matrix algebra of the irreducible block.\end{physics}
\begin{status}Unconditional; audited, standard axioms.\end{status}

\begin{theorem}[Rationality from rational traces]\label{ch22:thm:descent}
(a) If $H\in M_{m\times n}(\C)$ and $H_0\in M_{m\times n}(\Q)$ satisfy $f(H)=f(1)H_0$ entrywise for every $\Q$-linear
$f:\C\to\Q$, then $H=H_0$.
(b) In the setting of Theorem~\ref{ch22:thm:schurspan} over $K=\Q$, if $H\in M_n(\C)$ and all traces
$\tr(\sigma(s^{-1})EH)$ are rational numbers $q(s)$, then $EHE$ is the rational matrix
$\frac{\tr E}{|S|}\sum_sq(s)\sigma(s)E$.
\end{theorem}
\begin{leancode}
theorem eq_map_ratCast_of_forall_dual (H : Matrix m n ℂ) (Z : Matrix m n ℚ)
    (h : ∀ f : ℂ →ₗ[ℚ] ℚ, H.map f = f 1 • Z) : H = Z.map (Rat.castHom ℂ)
-- ⋯
theorem eq_map_ratCast_of_rational_traces (σ : S →* Matrix n n ℚ) (E : Matrix n n ℚ)
    (hE : E * E = E) (hEσ : ∀ s, E * σ s = σ s * E) (hr : E.trace ≠ 0)
    (hirr : ∀ T : Matrix n n ℚ, (∀ s, T * σ s = σ s * T) → ∃ c : ℚ, E * T * E = c • E)
    (H : Matrix n n ℂ) (q : S → ℚ)
    (hq : ∀ s, ((σ s⁻¹ * E).map (Rat.castHom ℂ) * H).trace = (q s : ℂ)) :
    E.map (Rat.castHom ℂ) * H * E.map (Rat.castHom ℂ)
      = ((E.trace / Fintype.card S) • ∑ s, q s • (σ s * E)).map (Rat.castHom ℂ)
\end{leancode}
(\leanpath{VOA/Moonshine/MoonshineCNRationalSchur.lean}.)
\begin{dictionary}
\lean{H.map f} applies $f$ entrywise; \lean{f 1 • Z} is $f(1)H_0$ (Lean calls $H_0$ \lean{Z}).  \lean{Z.map (Rat.castHom ℂ)} is $H_0$ read as a complex
matrix.  In (b), \lean{(σ s⁻¹ * E).map (Rat.castHom ℂ) * H} is the complex product $\sigma(s^{-1})EH$.
\end{dictionary}
\begin{proof}[Proof on paper]
(a) For each entry, $f(H_{ij}-(H_0)_{ij})=0$ for every $\Q$-linear $f:\C\to\Q$; since $\C$ is a $\Q$-vector space, an element
killed by all functionals is $0$ (\lean{Module.forall_dual_apply_eq_zero_iff}).  (b) Apply $f$ entrywise to $EHE$: this
commutes with multiplication by the rational $E$, so $f(EHE)=Ef(H)E$, and Theorem~\ref{ch22:thm:schurspan} over $\Q$
writes it as $\frac{\tr E}{|S|}\sum_s\tr(\sigma(s^{-1})Ef(H))\sigma(s)E$; and $\tr(\sigma(s^{-1})Ef(H))=f(\tr(\sigma(s^{-1})EH))=
f(q(s))=q(s)f(1)$.  So $f(EHE)=f(1)H_0$ with $H_0$ independent of $f$, and (a) gives $EHE=H_0$.  This is Galois descent with the
field automorphisms of $\C$ (which need transcendence bases) replaced by $\Q$-linear functionals.
\end{proof}
\begin{physics}A symmetry whose traces against all elements of an irreducible rational symmetry group are rational is
itself rational on that irreducible block.\end{physics}
\begin{status}Unconditional; audited, standard axioms.\end{status}

%======================================================================
\section{The converse, III: rationality of \texorpdfstring{$\mathrm{Aut}(V^\natural)$}{Aut(V-natural)} given the commutant}\label{ch22:sec:rationalgriess}
%======================================================================

\begin{definition}[The identity line and its projection]\label{ch22:def:P0}
$\mathbf \mathbf{e}_1=(I,0,0)\in\chtwentytwoB$ ($I$ the $24\times24$ identity matrix) and $P_0u=\frac{b(u,\mathbf \mathbf{e}_1)}{24}\,\mathbf \mathbf{e}_1$, the
$b$-orthogonal projection onto $\Q\mathbf \mathbf{e}_1$ ($b=$ \lean{bform}, $b(\mathbf \mathbf{e}_1,\mathbf \mathbf{e}_1)=\tr I=24$).
\end{definition}
\begin{leancode}
noncomputable def griessIdentityVec : SpaceB :=
  (⟨1, mem_symMat.2 Matrix.transpose_one⟩, 0, 0)
-- ⋯
noncomputable def griessUnitProj : SpaceB →ₗ[ℚ] SpaceB :=
  LinearMap.smulRight ((24 : ℚ)⁻¹ • bform.flip griessIdentityVec) griessIdentityVec
-- ⋯
theorem griessUnitProj_apply (u : SpaceB) :
    griessUnitProj u = (bform u griessIdentityVec / 24) • griessIdentityVec
\end{leancode}
(\leanpath{VOA/Moonshine/MoonshineCNRationalGriess.lean}, namespace \chtwentytwolean{VOA.Moonshine.MonsterAction}.)
\begin{dictionary}
\lean{⟨1, mem_symMat.2 Matrix.transpose_one⟩} is the identity matrix with the proof that it is symmetric.
\lean{LinearMap.smulRight φ v} is $u\mapsto\varphi(u)v$; \lean{bform.flip griessIdentityVec} is $u\mapsto b(u,\mathbf{e}_1)$.  The
\lean{_apply} lemma states the textbook formula.
\end{dictionary}
\begin{physics}
$\mathbf{e}_1$ corresponds to (twice) the stress tensor: AxiomMath's unit is $\frac14 \mathbf{e}_1$, which acts on $\chtwentytwoB$ by the
identity; $\omega$ acts on $\Vnat_2$ by $L_0=2$.  $P_0$ projects onto the stress-tensor direction, and $1-P_0$ onto the
$196883$ further primaries.
\end{physics}
\begin{status}Unconditional definitions.\end{status}

\begin{theorem}[The identity line is the conformal line, and is fixed]\label{ch22:thm:omega}
$\mathrm{griessNat}(\mathbf{e}_1)=2\omega$; hence every automorphism of $\Vnat$ fixes $\mathrm{griessNat}(\mathbf{e}_1)$, every element of
$\chtwentytwoGmat$ fixes $\mathbf{e}_1$, and $P_0$ commutes with $\chtwentytwoGmat$.  Moreover $\chtwentytwoGmat$ is finite.
\end{theorem}
\begin{leancode}
theorem griessNat_griessIdentityVec :
    griessNat griessIdentityVec = (2 : ℂ) • vnatConformalTwo
-- ⋯
theorem autVnat_apply_griessNat_identityVec (h : AutVnat) :
    h (griessNat griessIdentityVec : Vnat) = griessNat griessIdentityVec
-- ⋯
theorem finite_gmat : Finite Generated.Gmat
\end{leancode}
(\leanpath{VOA/Moonshine/MoonshineCNRationalGriess.lean}.)
\begin{dictionary}
\lean{vnatConformalTwo} is $\omega$ as an element of $\Vnat_2$.  \lean{Finite Generated.Gmat}: the subgroup, as a type, is
finite.
\end{dictionary}
\begin{proof}[Proof on paper]
$\frac14\mathbf{e}_1$ is the unit of AxiomMath's algebra (\lean{gmul_oneB_left}, a direct computation from the three block products),
so $\mathrm{griessNat}(\mathbf{e}_1)$ acts on $\Vnat_2$ by multiplication by $4$; $\omega$ acts by $L_0=2$; the Griess product is
commutative; hence $4\omega=\mathrm{griessNat}(\mathbf{e}_1)\cdot\omega=\omega\cdot\mathrm{griessNat}(\mathbf{e}_1)=2\,\mathrm{griessNat}(\mathbf{e}_1)$.
Automorphisms fix $\omega$.  For $g\in\chtwentytwoGmat$, its lift fixes $\mathrm{griessNat}(\mathbf{e}_1)$ and acts as $g$ on the
rational lattice, so $ge_1=\mathbf{e}_1$; $g$ preserves $b$, so it commutes with $P_0$.  Finiteness: $\chtwentytwoGmat$ embeds in
$\Aut(\Vnat)$ by \lean{gmatLift} (Theorem~\ref{ch22:thm:gmat}), which is finite (\lean{M8.finite_autVnat}, Chapter~\ref{ch:20}).
\end{proof}
\begin{physics}
The stress tensor is invariant under every symmetry.  The finiteness of $\chtwentytwoGmat$ is a notable by-product: it is proved
here from the finiteness of the CFT's symmetry group, without any group-theoretic computation of $|\chtwentytwoGmat|$.
\end{physics}
\begin{status}Unconditional; audited in \leanpath{VOA/Moonshine/MoonshineCNRationalAudit.lean}, standard axioms.\end{status}

\begin{definition}[The commutant property]\label{ch22:def:commutant}
\emph{$\chtwentytwoGmat$ has scalar commutant} if every $\Q$-linear $\mathcal K:\chtwentytwoB\to\chtwentytwoB$ with $\mathcal Kg=g\mathcal K$ for
all $g\in\chtwentytwoGmat$ is of the form $a\cdot1+b\cdot P_0$, $a,b\in\Q$.
\end{definition}
\begin{leancode}
def GmatCommutantScalar : Prop :=
  ∀ T : SpaceB →ₗ[ℚ] SpaceB, (∀ g ∈ Generated.Gmat, ∀ u : SpaceB, T (g u) = g (T u)) →
    ∃ a b : ℚ, T = a • LinearMap.id + b • griessUnitProj
-- ⋯
theorem gmat_commute_griessUnitProj (a b : ℚ) {g : SpaceB ≃ₗ[ℚ] SpaceB}
    (hg : g ∈ Generated.Gmat) (u : SpaceB) :
    (a • LinearMap.id + b • griessUnitProj : SpaceB →ₗ[ℚ] SpaceB) (g u)
      = g ((a • LinearMap.id + b • griessUnitProj : SpaceB →ₗ[ℚ] SpaceB) u)
\end{leancode}
(\leanpath{VOA/Moonshine/MoonshineCNRationalGriess.lean}, namespace \chtwentytwolean{VOA.Moonshine.MonsterAction}.)
\begin{dictionary}
\textbf{Decoding.}  \lean{T : SpaceB →ₗ[ℚ] SpaceB} is our $\mathcal K$, any $\Q$-linear endomorphism (not assumed invertible or
form-preserving).  The hypothesis \lean{∀ g ∈ Generated.Gmat, ∀ u, T (g u) = g (T u)} is $\mathcal Kg=g\mathcal K$ for all $g$.
\lean{a • LinearMap.id + b • griessUnitProj} is $a\cdot1+bP_0$.

\textbf{Meaning.}  Since $1$ and $P_0$ commute with $\chtwentytwoGmat$ (\lean{gmat_commute_griessUnitProj}), the proposition
says the commutant is \emph{exactly} the two-dimensional algebra $\Q1\oplus\Q P_0$.  Equivalently,
$\chtwentytwoB=\Q \mathbf{e}_1\oplus \mathbf{e}_1^\perp$ (both summands $\chtwentytwoGmat$-stable; $\chtwentytwoGmat$ is finite,
Theorem~\ref{ch22:thm:omega}, so its representations are completely reducible, which is what makes the two formulations
equivalent) with $\chtwentytwoGmat$ \emph{absolutely} irreducible on the $196883$-dimensional
$\mathbf{e}_1^\perp$ (over $\Q$, the commutant of an irreducible representation is a division algebra; ``scalar'' means it is $\Q$,
which is equivalent to irreducibility over $\C$).

\textbf{Could it be vacuous?}  It is a $\forall$ over a nonempty class ($T=1$); it is false for a too-small group (e.g.\ for
the monomial subgroup, \S\ref{ch22:sec:commutant}), so it has content.
\end{dictionary}
\begin{physics}
The $196883$ non-stress-tensor primaries of weight $2$ form a single irreducible multiplet of the symmetry group
$\chtwentytwoGmat$, and remain irreducible over $\C$ --- the smallest faithful representation of the Monster.
\end{physics}
\begin{status}A definition (a \lean{Prop}); it is proved in \S\ref{ch22:sec:commutant} (\lean{M8.gmatCommutantScalar}).
\lean{gmat_commute_griessUnitProj}: unconditional, audited.\end{status}

\begin{theorem}[Rationality, given the commutant property]\label{ch22:thm:rationalcomm}
Assume $\chtwentytwoGmat$ has scalar commutant.  Then every automorphism $h$ of $\Vnat$ maps
$\mathrm{griessNat}(\chtwentytwoB)$ into itself; consequently \lean{monsterLift} is surjective, $\Aut(\Vnat)\cong$ the
liftable subgroup of $\chtwentytwoAutO$, and, if moreover $\chtwentytwoGmat=\chtwentytwoAutO$, $\Aut(\Vnat)\cong\chtwentytwoAutO$.
\end{theorem}
\begin{leancode}
theorem autVnat_rational_of_commutant (hcomm : GmatCommutantScalar) :
    ∀ h : AutVnat, ∀ u : SpaceB, ∃ u' : SpaceB, h (griessNat u : Vnat) = griessNat u'
-- ⋯
theorem monsterLift_surjective_of_commutant (hcomm : GmatCommutantScalar) :
    Function.Surjective monsterLift
-- ⋯
noncomputable def autVnat_mulEquiv_liftable_of_commutant (hcomm : GmatCommutantScalar) :
    AutVnat ≃* liftableSubgroup
-- ⋯
theorem autVnat_mulEquiv_autO_of_commutant_of_Gmat (hcomm : GmatCommutantScalar)
    (hGmat : Generated.Gmat = AutO) : Nonempty (AutVnat ≃* AutO)
\end{leancode}
(\leanpath{VOA/Moonshine/MoonshineCNRationalGriess.lean}.)
\begin{dictionary}
\lean{(griessNat u : Vnat)} is the weight-two vector regarded in $\Vnat$; \lean{AutVnat ≃* liftableSubgroup} a group
isomorphism (its inverse is \lean{monsterLift}, \lean{autVnat_mulEquiv_liftable_of_commutant_symm_apply});
\lean{Nonempty (AutVnat ≃* AutO)} asserts that an isomorphism exists.  Hypotheses: \lean{hcomm} (and \lean{hGmat} in the last).
\end{dictionary}
\begin{proof}[Proof on paper]
Work with matrices in the dictionary bases (\lean{griessBasisQ} of $\chtwentytwoB$, \lean{griessBasisC} its image in $\Vnat_2$);
in these bases the Griess conjugate of a rational map has the same rational matrix (\lean{toMatrix_griessLinear}).  Let $H$
be the complex matrix of $h|_{\Vnat_2}$, $\sigma:\chtwentytwoGmat\to M(\Q)$ the rational representation, $E=1-P_0$ (rational
idempotent, trace $196883\ne0$, commuting with $\sigma$).
\begin{enumerate}[label=(\arabic*)]
\item $P_0H=P_0$ and $HP_0=P_0$: $h$ fixes $\mathrm{griessNat}(\mathbf{e}_1)$ and preserves $\mathrm{natForm}=2b$
  (\lean{griessUnitMatrix_map_mul_vnatMatrix}, \lean{vnatMatrix_mul_griessUnitMatrix_map}).  So $H=P_0+EHE$.
\item The commutant property in matrix form gives the hypothesis of Theorem~\ref{ch22:thm:schurspan} for $E$
  (\lean{griessComplMatrix_commutant}).
\item The traces are rational: $\tr(\sigma(s^{-1})EH)=\tr(\sigma(s^{-1})H)-\tr(\sigma(s^{-1})P_0)$, the first term is the trace
  on $\Vnat_2$ of the automorphism $\mathrm{gmatLift}(s)^{-1}h$, rational by Theorem~\ref{ch22:thm:rationaltraces}, and the
  second is a rational number (\lean{exists_rat_trace_gmat_compl_vnatMatrix}).
\item Theorem~\ref{ch22:thm:descent}(b) shows $EHE$ rational, so $H$ is rational (\lean{exists_vnatMatrix_eq_map_of_commutant}).
\end{enumerate}
A rational matrix maps the rational lattice to itself.  Then \lean{isVnatLiftable_of_rational} makes every $h$ a lift, so
\lean{monsterLift} is bijective; with \lean{hGmat}, the liftable subgroup is all of $\chtwentytwoAutO$ (it contains
$\chtwentytwoGmat$).
\end{proof}
\begin{physics}
\textbf{Schur's lemma pins down the symmetry group.}  An arbitrary symmetry $h$ of the CFT has rational traces (moonshine)
and fixes the stress tensor.  Its matrix on the $196883$-multiplet, Fourier-expanded along the irreducible rational group
$\chtwentytwoGmat$ (Schur orthogonality), has rational coefficients; so $h$ preserves the rational Griess algebra and is
one of AxiomMath's $\chtwentytwoAutO$.  No classification of finite groups enters: irreducibility of one explicit multiplet and
the arithmetic of the partition functions suffice.
\end{physics}
\begin{status}\textbf{Conditional} on \lean{hcomm} (and \lean{hGmat} for the last); \lean{hcomm} is discharged in
\S\ref{ch22:sec:commutant}.  All four audited in \leanpath{VOA/Moonshine/MoonshineCNRationalAudit.lean} ($94$
\texttt{\#print axioms} commands), standard axioms.\end{status}

%======================================================================
\section{The commutant lemma}\label{ch22:sec:commutant}
%======================================================================

\begin{theorem}[The commutant lemma]\label{ch22:thm:commutant}
Every $\Q$-linear endomorphism of $\chtwentytwoB$ commuting with every element of $\chtwentytwoGmat$ is $a\cdot1+b\cdot P_0$ for
some $a,b\in\Q$.  Equivalently, $\chtwentytwoGmat$ is absolutely irreducible on the $196883$-dimensional $\mathbf{e}_1^\perp$.
\end{theorem}
\begin{leancode}
theorem gmatCommutantScalar : GmatCommutantScalar := by
\end{leancode}
(\leanpath{VOA/Moonshine/MoonshineLRCommutantFinal.lean}; full name
\chtwentytwolean{VOA.Moonshine.ConwayNorton.M8.gmatCommutantScalar}.)  The coordinate form, from which it is derived:
\begin{leancode}
def commutF_P₀ : SpaceB →ₗ[ℚ] SpaceB :=
  ((1 / 24 : ℚ) • bform.flip commutF_eOne).smulRight commutF_eOne
-- ⋯
theorem commutF_commutant_Gmat (T : SpaceB →ₗ[ℚ] SpaceB)
    (hT : ∀ g ∈ Generated.Gmat, ∀ u, T (g u) = g (T u)) :
    ∃ a b : ℚ, T = a • LinearMap.id + b • commutF_P₀
\end{leancode}
\begin{leancode}
theorem commutF_P₀_eq_griessUnitProj : commutF_P₀ = griessUnitProj
\end{leancode}
(\leanpath{VOA/Moonshine/MoonshineLRCommutant4.lean}, \leanpath{VOA/Moonshine/MoonshineLRCommutantFinal.lean}.)
\begin{dictionary}
\textbf{Decoding.}  \lean{GmatCommutantScalar} is Definition~\ref{ch22:def:commutant}, quoted in full there.
\lean{commutF_P₀} is the same projection as \lean{griessUnitProj} written independently ($\frac1{24}$ instead of
$24^{-1}$, a separately defined $\mathbf{e}_1$ called \lean{commutF_eOne}); \lean{commutF_P₀_eq_griessUnitProj} proves them equal, and
\lean{commutF_eOne_eq_griessIdentityVec} holds by \lean{rfl}.  So the theorem is exactly the named input of
Theorem~\ref{ch22:thm:rationalcomm}.

\textbf{Hypotheses.}  None.  \textbf{Same statement.}  Yes, and \lean{gmat_commute_griessUnitProj} gives the converse
inclusion, so the commutant is exactly $\Q1\oplus\Q P_0$.

\textbf{Is it a strong statement?}  Yes: the analogous statement for proper subgroups of $\chtwentytwoGmat$ can fail.  For
the subgroup generated by $Q$, $x_B$, $\tau_B$ and the monomial lifts it is false (the warning after Step 7 below exhibits an
extra commuting projection), so the theorem genuinely uses the non-monomial element $\xi$.
\end{dictionary}

\begin{physics}
The $196883$ weight-two primary fields of the moonshine CFT form one irreducible multiplet under the subgroup of symmetries
generated by the orbifold-visible symmetries and triality --- even over $\C$.  The proof is a Clifford-theory analysis: the
extraspecial ``momentum-phase'' symmetries split the multiplet into many lines with distinct charges, and the remaining
symmetries (lattice rotations, $\xi$, triality) are shown to connect all the lines.
\end{physics}

\begin{status}
\textbf{Unconditional.}  \lean{gmatCommutantScalar}: \texttt{\#print axioms} standard in
\leanpath{VOA/Moonshine/MoonshineLRLiftAudit.lean} and \leanpath{VOA/Moonshine/MoonshineCNSymmetryFinalAudit.lean};
\lean{commutF_commutant_Gmat} audited in the former.  The finite certificates are kernel \lean{decide} (84 occurrences in
\leanpath{VOA/Moonshine/MoonshineLRCommutant4.lean}); no \lean{native_decide}.  Rebuild of the audit: 9,455 jobs,
0 \lean{sorryAx}.
\end{status}

\subsection{The proof on paper}

\paragraph{Setting and notation.}  Let $\mathcal K:\chtwentytwoB\to\chtwentytwoB$ be $\Q$-linear and commute with every element
of $\chtwentytwoGmat$.  We use the following elements of $\chtwentytwoGmat$: the extraspecial group $Q\cong2^{1+24}$ acting through
AxiomMath's \lean{rhoEquiv} (all in $\chtwentytwoGvis$, \lean{commutRho_mem_Gmat}); $x_B$; $\tau_B$; the visible lifts of the
code-preserving permutations and of $\xi$ (\S\ref{ch22:sec:generators}).  $\mathcal D,\mathcal D_0,\mathcal O\subset\chtwentytwoA$ are
the diagonal, traceless diagonal and off-diagonal symmetric matrices; $E_{ii}$ the diagonal matrix units, $U_{jk}$ ($j\ne k$)
the off-diagonal ones.

\emph{Short elements and lines of $\chtwentytwoX$.}  $Q$ is AxiomMath's group of $2^{25}$ signed permutation matrices on the
spin index; its centre is $\{1,z_Q\}$ and $Q/\{1,z_Q\}\cong\Leech/2\Leech$.  An element $t\in Q$ is \emph{short} if its class
is the class of a minimal vector.  For short $t$, \lean{Xlift t} $\in\chtwentytwoX$ is $\pm$ the basis vector of the
corresponding short coordinate, with $X_{z_Qt}=-X_t$; so the $2\cdot98280$ short elements give the $98280$ lines $\Q X_t$ of
$\chtwentytwoX$, each twice.

\emph{Spin index.}  Rows of $N\in\chtwentytwoW$ are indexed by $(\eta,a)\in\Z/2\times C_0$, $C_0$ a $2^{11}$-element subspace of
the Golay code (\S\ref{ch22:sec:tauB}); $a$ is the \emph{code component}, $E(\eta,a)$ the corresponding basis vector of
$\Q^{4096}$, and $E(\eta,a)\otimes e_i$ the matrix with a single $1$ in row $(\eta,a)$, column $i$.

\emph{MOG, sextets, transverse partitions.}  In the MOG numbering the $24$ coordinates are $\mathrm{coord}(r,j)=r+4j$, and
\emph{column} $j$ is $\{4j,\dots,4j+3\}$.  A \emph{sextet} is a partition of the $24$ points into six $4$-sets (tetrads) such
that the union of any two tetrads is an octad (a weight-$8$ Golay codeword); the six MOG columns form a sextet.  Two
partitions of $\{0,\dots,23\}$ are \emph{transverse} if the only functions constant on the blocks of both are the constant
functions (their join is the one-block partition).

\paragraph{Roadmap.}  Steps 1--7 show that $\mathcal K$ is diagonal in AxiomMath's natural basis, with one scalar $d_i$ per column
$i$ of $\chtwentytwoW$ (shared by the odd short lines attached to it), one scalar per pair $\{j,k\}$ (shared by $U_{jk}$ and two
short lines), one scalar per octad block, and that $\mathcal K$ preserves $\mathcal D$.  Steps 8--9 show that all these scalars are
equal to one number $d$: permutation lifts identify scalars of the same type, and $\xi$ connects different types (an odd line to
an octad line, a pair line to an octad line).  Step 10 shows that $\mathcal K-d\cdot1$ vanishes except on $\mathcal D$, where it is
a multiple of $A\mapsto\tr(A)\,I$, i.e.\ of $P_0$.

\paragraph{Step 1 (parity blocks).}  $x_B$ is $-1$ on $\chtwentytwoW$ and $+1$ on $\chtwentytwoA\oplus\chtwentytwoX$, so
$\mathcal K$ preserves both (\lean{commutPreservesAX}, \lean{commutPreservesW}).

\paragraph{Step 2 ($Q$-characters).}  $Q$ fixes $\chtwentytwoA$ pointwise and acts on $X_t$ by the commutator sign
$\kappa(r,t)=\pm1$ (because $X_{z_Qt}=-X_t$, \lean{commutXlift_zQ_mul}, a lemma of AxiomMath's library that was not in the port
and was reproved).  The characters $\kappa(\cdot,t)$ of different lines are different (\lean{commutChar_injective}) and
nontrivial (\lean{commutChar_nontrivial}).  Hence $\chtwentytwoA$ is the joint fixed space of $Q$ on
$\chtwentytwoA\oplus\chtwentytwoX$ (\lean{commutFixed_eq_A}), so $\mathcal K(\chtwentytwoA)\subseteq\chtwentytwoA$; and each line $\Q
X_t$ is a joint eigenline with its own character, so $\mathcal KX_t=c_tX_t$ for a scalar $c_t$ (\lean{commutXlift_eigen}).  No
averaging is needed: compare components of a joint eigenvector.

\paragraph{Step 3 ($\tau_B$ diagonalises $\mathcal K$ on $\chtwentytwoW$).}  The odd triple of $(a,i)$ in the triality basis is
$(X_{r(a,i)},\ \pm E(1,a)\otimes e_i,\ \pm E(0,a)\otimes e_i)$ with $X_{r(a,i)}$ an odd short line, and $\tau_B$ cycles it
(\lean{tauB_odd}).  $\mathcal K$ is scalar on the first vector (Step 2) and commutes with $\tau_B$, so it is the same scalar on the
other two (\lean{commutOdd_eigen}, \lean{commutSpin_eigen}).

\paragraph{Step 4 (columns).}  The elements of $Q$ coming from codewords translate the code component $a$:
$E(\eta,a)\otimes e_i\mapsto E(\eta,a')\otimes e_i$ for any $a,a'$ (\lean{commutB_rho_translate}).  So the scalar of Step 3
depends only on the column $i$: $\mathcal K(0,0,N)=(0,0,N\,\mathrm{diag}(d_0,\dots,d_{23}))$ (\lean{commutB_columns}).

\paragraph{Step 5 (pairs).}  For $j<k$, $\tau_B$ cycles $U_{jk}$ and the two combinations $X_m\mp X_p$ of the short lines of
$4e_j-4e_k$ and $4e_j+4e_k$ (\lean{tauB_pair}).  $\mathcal K$ is scalar on $X_m$ and on $X_p$ (Step 2); writing out
$\mathcal K\tau_B=\tau_B\mathcal K$ on the triple forces the two scalars to be equal and equal to the eigenvalue on $U_{jk}$
(\lean{commutB_pair_lines}).

\paragraph{Step 6 (octads).}  On the $64$ short lines of an octad, $\tau_B$ acts by the matrix
$\mathrm{octHD}=\mathrm{octH}\cdot\mathrm{octD}$, a Fourier matrix times a sign diagonal, every entry of which is $\pm\frac18\ne0$
(\lean{commutB_octHD_ne_zero}).  A diagonal map commuting with a matrix without zero entries is scalar
(\lean{commutB_oct_scalar}).

\paragraph{Step 7 (the diagonal is preserved).}  Let $\Delta\in\mathcal D$ and $\Delta'=\mathcal K(\Delta)$, which lies in
$\chtwentytwoA$ by Step 2.  Two facts about $\tau_B$ on $\chtwentytwoA$: it fixes every diagonal matrix
(\lean{commutB_tau_diagonal}), and for every $A\in\chtwentytwoA$ the $\chtwentytwoA$-component of $\tau_B(A)$ is diagonal
(\lean{commutB_tau_A_offdiag}: $\tau_B$ sends each $U_{jk}$ into $\chtwentytwoX$, the first vector of a pair triple going to the
second).  Since $\tau_B\Delta=\Delta$, we get $\tau_B\Delta'=\mathcal K(\tau_B\Delta)=\Delta'$.  Taking $\chtwentytwoA$-components,
$\Delta'=(\tau_B\Delta')_{\chtwentytwoA}$, which is diagonal.  So $\mathcal K(\mathcal D)\subseteq\mathcal D$
(\lean{commutB_preserves_diagonal}).

\begin{warning}
\textbf{Dead end: the monomial subgroup is not enough.}  Steps 1--7 use $Q$, $x_B$
and $\tau_B$.  Adding the monomial lifts $2^{12}{:}M_{24}$ still does not suffice: this subgroup preserves the $23$-dimensional
space $\mathcal D_0$ ($\tau_B$ fixes each $E_{ii}$, monomial elements permute them, $Q$ and $x_B$ act trivially on
$\chtwentytwoA$), so its commutant contains the projection onto $\mathcal D_0$, which is not of the form $a+bP_0$.  It also
preserves the decomposition into octad, pair and odd$\cup\chtwentytwoW$ types up to $M_{24}$.  A non-monomial element of
$\chtwentytwoCo_0$, such as $\xi$, is indispensable.
\end{warning}

\paragraph{Step 8 (the linking elements are in $\chtwentytwoGmat$, with explicit blocks).}  A liftable element $g$ of
$\chtwentytwoAutO$ of block form
\begin{equation}\label{ch22:eq:blockform}
  g(A,x,N)=\bigl(\alpha(A),\ \beta(x),\ R\,N\,M\bigr)\qquad(R\in\mathrm{GL}_{4096}(\Q),\ M\in\mathrm{GL}_{24}(\Q))
\end{equation}
commutes with $x_B$, hence lies in $\chtwentytwoAutO\cap C(x_B)=\chtwentytwoGvis$ \eqref{ch22:eq:parity} (\lean{commutL_mem_Gvis}).
This applies to the lifts of code-preserving permutations $\pi$ (with $\alpha(A)=P_\pi^{\mathsf T}AP_\pi$, $M=P_\pi$ the permutation
matrix, $R$ a signed permutation) and of $\xi$ (\lean{commutL_mathieu_blocks}, \lean{commutL_xi_blocks}):
\begin{leancode}
theorem commutL_xi_blocks :
    ∃ g : liftableSubgroup,
      monsterLift g = liftTwVnatAut liftXi liftXi_isometry liftXi_surjective ∧
      ((g : AutO) : SpaceB ≃ₗ[ℚ] SpaceB) ∈ Gvis ∧
      ((g : AutO) : SpaceB ≃ₗ[ℚ] SpaceB) ∈ Generated.Gmat ∧
      ∀ u : SpaceB, ((g : AutO) : SpaceB ≃ₗ[ℚ] SpaceB) u =
        (thetaA liftXiMatrixᵀ u.1,
          griessTrX (liftSignX liftXi liftXi_isometry (griessTrX u.2.1)),
          liftXiSpinMatrix * u.2.2 * liftXiMatrixᵀ)
\end{leancode}
(\leanpath{VOA/Moonshine/MoonshineLRCommutant3.lean}; \lean{thetaA M A} $=M^{\mathsf T}AM$, so the $\chtwentytwoA$-block of $\xi$ is
$A\mapsto\Xi A\Xi^{\mathsf T}$ with $\Xi=$ \lean{liftXiMatrix}.)  On $\chtwentytwoW$, $g\mathcal K(0,0,N)=(0,0,RN\,\mathrm{diag}(d)M)$ and
$\mathcal Kg(0,0,N)=(0,0,RNM\,\mathrm{diag}(d))$; since $R$ is invertible and $N$ arbitrary, $\mathrm{diag}(d)M=M\mathrm{diag}(d)$
(\lean{commutL_column_matrix}).  For $M=P_\pi$ this says $d_{\pi(i)}=d_i$ (\lean{commutL_column_perm}).

\paragraph{Step 9 (certificates and transport).}  Three explicit code-preserving permutations are used: $\pi_A$
(\lean{commutL_pointPermA}), which swaps MOG columns $2\leftrightarrow3$ and $4\leftrightarrow5$ as blocks; $\pi_B$
(\lean{commutL_pointPermB}), which maps each column onto a column, permuting the points within it (e.g.\ $0,1,2,3\mapsto8,10,11,9$);
and $\pi_C$ (\lean{commutF_genC}), an involution that does \emph{not} preserve the columns ($0\mapsto23$ but $1\mapsto1$).  Their
tables, and the derived $\pi_{\mathrm{tr}}=\pi_C\pi_A\pi_C$ (\lean{commutF_transverse}):
\begin{leancode}
def commutL_pointPermA : Equiv.Perm (Fin 24) where
  toFun := ![0, 1, 2, 3, 4, 5, 6, 7, 12, 13, 14, 15,
    8, 9, 10, 11, 20, 21, 22, 23, 16, 17, 18, 19]
  invFun := ![0, 1, 2, 3, 4, 5, 6, 7, 12, 13, 14, 15,
    8, 9, 10, 11, 20, 21, 22, 23, 16, 17, 18, 19]
  left_inv := by decide
  right_inv := by decide
-- ⋯
def commutL_pointPermB : Equiv.Perm (Fin 24) where
  toFun := ![8, 10, 11, 9, 14, 12, 13, 15, 20, 22, 23, 21,
    18, 16, 17, 19, 7, 5, 4, 6, 1, 3, 2, 0]
  invFun := ![23, 20, 22, 21, 18, 17, 19, 16, 0, 3, 1, 2,
    5, 6, 4, 7, 13, 14, 12, 15, 8, 11, 9, 10]
  left_inv := by decide
  right_inv := by decide
\end{leancode}
\begin{leancode}
def commutF_genC : Equiv.Perm (Fin 24) where
  toFun := ![23, 1, 11, 7, 20, 8, 6, 3, 5, 22, 10, 2,
    12, 17, 14, 15, 19, 13, 18, 16, 4, 21, 9, 0]
  invFun := ![23, 1, 11, 7, 20, 8, 6, 3, 5, 22, 10, 2,
    12, 17, 14, 15, 19, 13, 18, 16, 4, 21, 9, 0]
  left_inv := by decide
  right_inv := by decide
-- ⋯
def commutF_transverse : Equiv.Perm (Fin 24) :=
  commutF_genC * commutL_pointPermA * commutF_genC
-- ⋯
def commutF_generator : Fin 3 → Equiv.Perm (Fin 24) :=
  ![commutL_pointPermA, commutL_pointPermB, commutF_genC]
\end{leancode}
(\leanpath{VOA/Moonshine/MoonshineLRCommutant3.lean}, \leanpath{VOA/Moonshine/MoonshineLRCommutant4.lean}; \lean{![…]} lists the
images of $0,\dots,23$.)  Each is checked to preserve the Golay code by kernel \lean{decide} on the twelve basis words
(\lean{commutL_preservesCode_basis}, \lean{commutF_genC_code}).
\begin{itemize}
\item \emph{Points.}  $\langle\pi_A,\pi_B\rangle$ is transitive on the $24$ points (\lean{commutL_point_transitive}); so $d$ is
  constant, $d_i=d$: $\mathcal K=d$ on $\chtwentytwoW$ and on every odd short line (\lean{commutL_spin_odd_scalar}).
\item \emph{One octad.}  $\xi$ maps an odd minimal vector to an octad-shaped vector
  (\leanpath{VOA/Moonshine/MoonshineLRCommutant4.lean}):
\begin{leancode}
theorem commutF_odd_octad_coordinates :
    liftXiInt (oddMin 0 (B 1)) = vOct (B 0) commutF_seedSigns := by decide
\end{leancode}
  (proved by \lean{decide}; \lean{oddMin 0 (B 1)} has entries $\pm1$ and one $\mp3$, signs given by the basis codeword $B_1$;
  \lean{vOct (B 0) S} has entries $\pm2$ on the octad $B_0$, minus signs on $S$.)  Transporting the eigenvalue through the lift
  of $\xi$ (\lean{commutL_eigen_transport}), the line of that octad vector has eigenvalue $d$, hence the whole octad block of $B_0$
  (Step 6; \lean{commutF_seed_octad_scalar}).
\item \emph{All octads.}  A spanning tree of the $759$ octads under $\pi_A,\pi_B,\pi_C$, rooted at $B_0$, is stored as a table: for
  each octad (a $24$-bit mask, sorted) its parent, the generator carrying parent to child, and a breadth-first rank.  The check,
  done by kernel \lean{decide} in $24$ chunks of $32$ (\leanpath{VOA/Moonshine/MoonshineLRCommutant4.lean}):
\begin{leancode}
def commutF_octadCheck (n : ℕ) : Prop :=
  if hn : n < 759 then
    let i : Fin 759 := ⟨n, hn⟩
    (∀ hs : n + 1 < 759, commutF_octadMask i < commutF_octadMask ⟨n + 1, hs⟩) ∧
    (i = 130 ∨
      (commutF_octadLink (commutF_octadLink i).1).2.2 < (commutF_octadLink i).2.2 ∧
      permLin (ZMod 2) (commutF_generator (commutF_octadLink i).2.1)
        (bvec 24 (commutF_octadMask (commutF_octadLink i).1)) =
          bvec 24 (commutF_octadMask i))
  else True
\end{leancode}
  i.e.\ masks strictly increase (no repetitions), and every non-root entry is the image of its parent under its generator, with
  the parent of strictly smaller rank.  So every listed mask is reachable from $B_0$; each is an octad (\lean{commutF_octad_valid});
  and since the $759$ masks are distinct octads and there are exactly $759$ octads (\lean{card_octadWord}, ported), every octad
  is listed (\lean{commutF_octadEnum_surjective}).  A permutation lift carries octad blocks to octad blocks with the same
  eigenvalue (\lean{commutF_octad_eigen_perm}), so every octad block has eigenvalue $d$ (\lean{commutF_octads_scalar}).
\item \emph{Pairs.}  $\xi$ carries $4e_0-4e_4$ (columns $0$ and $1$) to an octad vector on the first two MOG columns
  (\lean{commutF_pair_octad_coordinates}), so the pair $\{0,4\}$ has eigenvalue $d$; a spanning tree of the $276$ unordered
  pairs under the three generators (\lean{commutF_pairCheck}, same format) and the transport along permutation lifts give
  eigenvalue $d$ on every $U_{jk}$ and both short lines of every pair (\lean{commutF_offdiag_scalar},
  \lean{commutF_pair_short_scalar}).  Hence $\mathcal K=d$ on all of $\chtwentytwoX$ (\lean{commutF_X_scalar}).
\end{itemize}

\begin{warning}
\textbf{Recorded dead end: two generators that preserve the MOG columns.}  The first attempt at Step 9 used only
$\pi_A$ and $\pi_B$.  Both map columns onto columns, so $\langle\pi_A,\pi_B\rangle$ preserves the MOG column partition (it lies in
the stabiliser of the MOG sextet).  It is transitive on points, but it can never move a within-column pair to a cross-column
pair, cannot be transitive on the $759$ octads, and cannot produce a partition transverse to the MOG columns.  The fix is the
code-preserving involution $\pi_C=$ \lean{commutF_genC},
$[23,1,11,7,20,8,6,3,5,22,10,2,12,17,14,15,19,13,18,16,4,21,9,0]$, which breaks the columns; $\langle\pi_A,\pi_B,\pi_C\rangle$ is
transitive on octads and pairs (certified above), and the columns of $\pi_{\mathrm{tr}}=\pi_C\pi_A\pi_C$ are transverse to the MOG
columns.
\end{warning}

\paragraph{Step 10 (the diagonal and assembly).}  Let $\mathcal R=\mathcal K-d\cdot1$ (\lean{commutF_residual}); it commutes with
$\chtwentytwoGmat$, vanishes on $\chtwentytwoX$, $\chtwentytwoW$ and $\mathcal O$, and maps $\chtwentytwoA$ into $\mathcal D$
(\lean{commutF_residual_diagonal}).
\begin{itemize}
\item \emph{$\xi$ forces column-constancy.}  The lift of $\xi$ acts on $\chtwentytwoA$ by $A\mapsto\Xi A\Xi^{\mathsf T}$.  If
  $\Delta$ is diagonal with entries $\Delta_0,\dots,\Delta_3$ in a MOG column $j$, the $(\mathrm{coord}(0,j),\mathrm{coord}(r,j))$
  entry of $\Xi\Delta\Xi^{\mathsf T}$ is $\tfrac14(\Delta_0+\Delta_r-\Delta_s-\Delta_t)$, $\{r,s,t\}=\{1,2,3\}$
  (\lean{commutF_xi_correlation}).  If $\Xi\Delta\Xi^{\mathsf T}$ is also diagonal, the three equations
  $\Delta_0+\Delta_r=\Delta_s+\Delta_t$ force $\Delta_0=\Delta_1=\Delta_2=\Delta_3$ (\lean{commutF_xi_diagonal_columns}).  Since
  $\mathcal R(\Xi A\Xi^{\mathsf T})=\Xi\mathcal R(A)\Xi^{\mathsf T}$ and both residual images are diagonal, every $\mathcal R(A)$ is
  constant on MOG columns (\lean{commutF_residual_columns}).
\item \emph{The transverse partition.}  $\mathcal R$ also commutes with the lift of $\pi_{\mathrm{tr}}$, so $\mathcal R(A)$ is
  constant on the images of the MOG columns under $\pi_{\mathrm{tr}}$; a function constant on the blocks of both partitions is
  constant (\lean{commutF_transverse_constant}, \leanpath{VOA/Moonshine/MoonshineLRCommutant4.lean}:
\begin{leancode}
theorem commutF_transverse_constant {K : Type*} (f : Fin 24 → K)
    (hc : ∀ i j, col i = col j → f i = f j)
    (hp : ∀ i j, col i = col j → f (commutF_transverse i) =
      f (commutF_transverse j)) (i : Fin 24) : f i = f 0
\end{leancode}
  proved by an explicit chain of equalities, each step justified by a kernel \lean{decide} on the column data).  So
  $\mathcal R(A)\in\Q I$ (\lean{commutF_residual_identity}).
\item \emph{Trace.}  Permutation lifts make $\mathcal R(E_{ii})$ independent of $i$ (\lean{commutF_residual_diag_equal}), and
  $\mathcal R$ vanishes on $\mathcal O$, so $\mathcal R(A)=\tr(A)\,\mathcal R(E_{00})=\tr(A)\,\ell I$ for a scalar $\ell$
  (\lean{commutF_residual_trace}), and then $\mathcal R(u)=\tr(u_{\chtwentytwoA})\,\ell\,\mathbf{e}_1=24\ell\,P_0(u)$ for all $u$
  (\lean{commutF_residual_full}; $b(u,\mathbf{e}_1)=\tr u_{\chtwentytwoA}$, \lean{commutF_bform_identity}).
\end{itemize}
Therefore $\mathcal K=d\cdot1+24\ell\cdot P_0$. \qed

\begin{physics}
Step by step: the momentum-phase symmetries $2^{1+24}$ are ``charges'' that resolve the multiplet into lines; triality glues
each odd vertex-operator line to two twisted ground states and each root line to a current bilinear; lattice permutations
and the Hadamard rotation $\xi$ glue all lines of the same type and then the different types together; the only thing they
cannot move is the stress tensor.
\end{physics}

\begin{remark}[On the role of computation]
The kernel-checked tables are \emph{certificates} of finite combinatorial facts about the Golay code (orbit
transitivity, one transverse sextet, two witness vectors); every other step is a structural argument.  The tables were
generated by a script (verified first in Python); only the kernel
check counts.
\end{remark}


%======================================================================
\section{The converse, assembled: \texorpdfstring{$\mathrm{Aut}(V^\natural)$}{Aut(V-natural)} inside AxiomMath's Monster}\label{ch22:sec:converse}
%======================================================================

\begin{theorem}[The converse, without Tits]\label{ch22:thm:converse}
(a) Every automorphism $h$ of $\Vnat$ maps the rational Griess lattice into itself: for each $u\in\chtwentytwoB$ there is
$u'\in\chtwentytwoB$ with $h(\mathrm{griessNat}\,u)=\mathrm{griessNat}\,u'$.
(b) $\Aut(\Vnat)$ is isomorphic to the subgroup of $\chtwentytwoAutO$ of liftable elements; the isomorphism is inverse to
\lean{monsterLift}.
(c) If $\chtwentytwoGmat=\chtwentytwoAutO$, then $\Aut(\Vnat)\cong\chtwentytwoAutO$.
\end{theorem}
\begin{leancode}
noncomputable def autVnat_mulEquiv_liftable : AutVnat ≃* liftableSubgroup :=
  autVnat_mulEquiv_liftable_of_commutant gmatCommutantScalar
-- ⋯
theorem autVnat_mulEquiv_liftable_symm_apply (g : liftableSubgroup) :
    autVnat_mulEquiv_liftable.symm g = monsterLift g
\end{leancode}
(\leanpath{VOA/Moonshine/MoonshineLRCommutantFinal.lean}, namespace \chtwentytwolean{VOA.Moonshine.ConwayNorton.M8}.)
\begin{leancode}
theorem autVnat_rational (h : AutVnat) (u : SpaceB) :
    ∃ u' : SpaceB, h (griessNat u : Vnat) = griessNat u' :=
  autVnat_rational_of_commutant M8.gmatCommutantScalar h u
-- ⋯
theorem nonempty_autVnat_mulEquiv_liftable : Nonempty (AutVnat ≃* liftableSubgroup) :=
  ⟨M8.autVnat_mulEquiv_liftable⟩
-- ⋯
theorem autVnat_mulEquiv_autO (hGmat : Generated.Gmat = AutO) : Nonempty (AutVnat ≃* AutO) :=
  autVnat_mulEquiv_autO_of_commutant_of_Gmat M8.gmatCommutantScalar hGmat
\end{leancode}
(\leanpath{VOA/Moonshine/MoonshineCNSymmetryFinalAudit.lean}, namespace \chtwentytwolean{VOA.Moonshine.ConwayNorton}.)
\begin{dictionary}
\textbf{Decoding.}  \lean{AutVnat ≃* liftableSubgroup}: a multiplicative equivalence (group isomorphism) from $\Aut(\Vnat)$ onto
the subgroup \lean{liftableSubgroup} of \lean{AutO}, regarded as a group.  The \lean{_symm_apply} lemma identifies its inverse
with \lean{monsterLift}, so the isomorphism is the canonical one: $h\mapsto$ the unique $g\in\chtwentytwoAutO$ with
$h|_{\Vnat_2}=\mathcal G(1\otimes g)\mathcal G^{-1}$.  \lean{Nonempty (…)} is the propositional form (``there exists an isomorphism'').

\textbf{Hypotheses.}  (a), (b): none.  (c): \lean{hGmat : Generated.Gmat = AutO}, the body of \lean{GmatEqAutO}.

\textbf{Same statement.}  (b) is the precise form of ``$\Aut(\Vnat)$ is a subgroup of the Monster containing
$\chtwentytwoGmat$''; since \lean{liftableSubgroup} is defined as the set of liftable elements, (b) is not a tautology: its
content is that \emph{every} automorphism of $\Vnat$ is a lift, i.e.\ (a).

\textbf{Could it be vacuous?}  No: $\Aut(\Vnat)$ is nontrivial (it contains the image of $\chtwentytwoGmat$), and
(a) is a statement about every element of it.
\end{dictionary}
\begin{proof}[Proof on paper]
Theorem~\ref{ch22:thm:rationalcomm} with the hypothesis discharged by Theorem~\ref{ch22:thm:commutant}.
\end{proof}
\begin{physics}
\textbf{The Monster contains the full symmetry group.}  Every symmetry of the moonshine CFT is, on the $196884$ weight-two
fields, one of AxiomMath's rational Griess-algebra automorphisms; together with \S\S\ref{ch22:sec:visible}--\ref{ch22:sec:tauB},
$\chtwentytwoGmat\subseteq\Aut(\Vnat)\subseteq\chtwentytwoAutO$.  The identification of the symmetry group is reduced to
the purely group-theoretic statement $\chtwentytwoGmat=\chtwentytwoAutO$, which no longer mentions the CFT.
\end{physics}
\begin{status}
(a), (b): \textbf{unconditional}.  (c): conditional on \lean{hGmat}.  Coordinator audit
\leanpath{VOA/Moonshine/MoonshineCNSymmetryFinalAudit.lean} (9,455 jobs): \lean{[propext, Classical.choice, Quot.sound]},
$0$ \lean{sorryAx}, for all three and for \lean{M8.gmatCommutantScalar}, \lean{M8.autVnat_mulEquiv_liftable},
\lean{rationalTraces}, \lean{conwayNorton}, \lean{conwayNorton_gmat}.
\end{status}


%======================================================================
\section{The headline: monstrous moonshine for AxiomMath's Monster}\label{ch22:sec:headline}
%======================================================================

The file \leanpath{VOA/Moonshine/MoonshineLRMonsterHeadline.lean} (namespace \chtwentytwolean{VOA.Moonshine.ConwayNorton}) states the
final form of the part's theorems in terms of the Monster itself.  It has four declarations of substance.

\subsection{The named hypothesis}

\begin{definition}[The generation theorem, as a proposition]\label{ch22:def:GmatEqAutO}
$\mathrm{GmatEqAutO}$ is the proposition $\chtwentytwoGmat=\chtwentytwoAutO$: the subgroup of $\mathrm{GL}_\Q(\chtwentytwoB)$
generated by the visible group and $\tau_B$ is the full group of $\Q$-linear automorphisms of AxiomMath's rational Griess
algebra preserving its product and its form.  It is AxiomMath's statement \lean{gmat_eq_griess} after renaming, with one
definition (\lean{SpinLattice.phi0}) restated in a definitionally equal form; that statement is proved in AxiomMath's library (Lean~4.35, commit \texttt{69cb972f}); the literature milestones on its proof path that were
inspected (\lean{cl_lit_stroth_leon_sims_baby_monster_recognition}, its \lean{_models} variant, and
\lean{cl_lit_goldschmidt_strongly_closed_abelian}) are theorems there, not axioms; its full axiom footprint has not been
re-checked (that needs \texttt{\#print axioms} on a Lean~4.35 build of AxiomMath's library).
\end{definition}
\begin{leancode}
def GmatEqAutO : Prop :=
  Generated.Gmat = AutO
\end{leancode}
(\leanpath{VOA/Moonshine/MoonshineLRMonsterHeadline.lean}; full name \chtwentytwolean{VOA.Moonshine.ConwayNorton.GmatEqAutO}.)

\begin{dictionary}
\textbf{What it says, symbol by symbol.}  \lean{Generated.Gmat} is
\lean{Subgroup.closure ((Gvis : Set (SpaceB ≃ₗ[ℚ] SpaceB)) ∪ {tauB})} (\S\ref{ch22:sec:recall}), a subgroup of the group
\lean{SpaceB ≃ₗ[ℚ] SpaceB} of $\Q$-linear automorphisms; \lean{AutO} is the subgroup of those $g$ with
$g(u\cdot v)=gu\cdot gv$ and $b(gu,gv)=b(u,v)$ (\S\ref{ch22:sec:recall}).  \lean{=} is equality of subgroups, i.e.\ equality of
the underlying sets.  Since $\chtwentytwoGmat\le\chtwentytwoAutO$ is a theorem (\lean{Generated.Gmat_le_AutO}, with
\lean{tauB_mem_AutO}), the content of \lean{GmatEqAutO} is the reverse inclusion: \emph{every} automorphism of the rational
Griess algebra is a word in the visible elements and $\tau_B$.

\textbf{It is AxiomMath's statement after renaming, with \lean{phi0} restated definitionally equal.}  AxiomMath's CFSG
library, at commit \texttt{69cb972f587a3ecc035e092253167ff2e303c962}, contains the theorem
\begin{plaincode}
theorem gmat_eq_griess :
    CFSG.Monster.Generated.Gmat = CFSG.Monster.Griess.AutO :=
  CFSG.YGroups.Y443.Gmat_eq_AutO_of_stroth
\end{plaincode}
(full name \chtwentytwolean{CFSG.Blueprint.NativeSporadic.MonsterModel.gmat_eq_griess}, file
\texttt{CFSG/Sporadic/Bridge/MonsterModelsGriess.lean}; quoted in \texttt{plaincode} because it is not a file of this
library).  Both sides are written with fully qualified names, they are the declarations the port copied, and both have type
\lean{Subgroup (SpaceB ≃ₗ[ℚ] SpaceB)}.  The project's fidelity audit (\texttt{Notes/monster\_aut/CFSG\_PORT\_FIDELITY.md})
compared the definitional closure of the two sides --- $199$ definition-like declarations and $104$ theorem
statements that fix data (through implicit arguments or \lean{Classical.choose} predicates), $46$ instances, and the
\lean{variable} contexts --- with AxiomMath's source after normalising comments, attributes, visibility keywords and the
namespace prefixes \chtwentytwolean{CFSG.Monster.Griess.}, \chtwentytwolean{CFSG.Monster.} $\mapsto$ \chtwentytwolean{VOA.Moonshine.MonsterGriess.}.  Verdict:
$194$ definition bodies are textually identical; three differ only in proof arguments of \lean{Prop} type (\lean{bform},
\lean{gmul}) or by a parser artefact (\lean{signMatrix}); the one genuine difference is
\begin{itemize}
\item \lean{SpinLattice.phi0}: ours is the plain function $p\mapsto\mathrm{lamT}(p_1)+\mathrm{lamCocode}(p_2)$; AxiomMath's is
  the bundled $\Z/2$-linear map \lean{(AddMonoidHom.toZModLinearMap 2 lamTHom).coprod lamCocode}.  The two are
  \emph{definitionally equal} as functions (\lean{LinearMap.coprod f g} is $f\circ\mathrm{fst}+g\circ\mathrm{snd}$ pointwise,
  and \lean{lamTHom} has \lean{toFun := lamT}; both libraries have \lean{phi0_apply} by \lean{rfl}).  It enters only through
  \lean{IsShortCoord} (and hence the index type \lean{ShortIdx} and the \lean{choose} in \lean{uvec}), where the two
  predicates are kernel-convertible.
\end{itemize}
So \lean{GmatEqAutO} is AxiomMath's statement up to renaming and definitional unfolding of one definition; it is not
syntactically the same term.  (Chapter~\ref{ch:21} gives the port audit in detail.)

\textbf{What the audit does not cover} (stated in the audit itself).  (i) Mathlib drift: the closure bottoms out in Mathlib
constants (\lean{Subgroup.closure}, \lean{Matrix.toLinOfInv}, \lean{LinearMap.mk₂}, \lean{Classical.choose},
\dots), assumed to mean the same in Mathlib v4.30 and in AxiomMath's pinned Mathlib; not compared.  (ii) No cross-toolchain
build: the comparison is source text, not a Lean proof that the two terms are equal.  (iii) The axiom footprint of
\lean{gmat_eq_griess} inside AxiomMath's library.  Its proof passes through
\lean{CFSG.Milestones.cl_lit_stroth_leon_sims_baby_monster_recognition_models}.  Inspection of the
source at the pinned commit (reading the files as data) shows that this declaration and
\lean{cl_lit_stroth_leon_sims_baby_monster_recognition} (file \texttt{CFSG/Milestones/CL\_LIT\_Stroth\_Leon\_Sims\_Baby\_Monster\_recognition.lean})
are \emph{theorems}, proved from \lean{CFSG.Milestones.cl_lit_goldschmidt_strongly_closed_abelian}, itself a theorem
(\texttt{CFSG/Milestones/CL\_LIT\_Goldschmidt\_strongly\_closed\_abelian.lean}); in AxiomMath's naming, \lean{CL_LIT} marks a
literature milestone that is \emph{proved}, not assumed.  Nobody in this project has run \texttt{\#print axioms} on
\lean{gmat_eq_griess}, which needs a Lean~4.35 build of AxiomMath's library; so its full axiom footprint is not re-checked.

\textbf{Why a named hypothesis, not an axiom.}  The library's policy (\texttt{README.md}, \texttt{CLAUDE.md} \S3) is:
\emph{no project axioms} in production code.  An \lean{axiom} would silently enter the \texttt{\#print axioms} output of
every theorem using it and would be global; a definition \lean{GmatEqAutO : Prop} used as an explicit argument
\lean{(hM : GmatEqAutO)} makes the dependence visible in the type of every theorem that needs it, and leaves every other
theorem untouched.  The theorem itself cannot be imported: AxiomMath's library is on Lean v4.35.0-rc3 with a Mathlib about
four months newer than this library's Lean 4.30 / Mathlib v4.30, and porting the proof down is infeasible --- its
file-level import closure is about $15{,}500$ CFSG files ($\approx16{,}000$ with the TauCeti dependency; $\approx197$--$208$\,MB,
$4.2$\,M lines), dominated by the Stroth--Leon--Sims recognition of the Baby Monster, with $1{,}720$ certificate files.  No Stroth-free proof of $\chtwentytwoGmat=\chtwentytwoAutO$ exists in that library.

\textbf{Could it be false, or trivially true?}  It is not trivially true: the inclusion $\chtwentytwoAutO\le\chtwentytwoGmat$
is the deep content of the identification of the Monster with the automorphism group of the Griess algebra (classically:
Conway's group generated by $C$ and the triality element is $\chtwentytwoM$, and Tits proved that $\chtwentytwoM$ is the full
automorphism group of the Griess algebra).  It is not refuted by anything in this library: everything proved here is
consistent with it (e.g.\ \lean{finite_gmat} shows $\chtwentytwoGmat$ finite, as $\chtwentytwoAutO$ must then be).
\end{dictionary}

\begin{physics}
The proposition is purely group-theoretic: it says that the symmetries of the $196884$ weight-two OPE coefficients that one can
write down explicitly (orbifold-visible ones and triality) generate \emph{all} symmetries of those OPE coefficients.  It no
longer refers to the CFT.
\end{physics}

\begin{status}
A definition (a \lean{Prop}); \texttt{\#print axioms} \chtwentytwolean{VOA.Moonshine.ConwayNorton.GmatEqAutO} is printed in
\leanpath{VOA/Moonshine/MoonshineLRLiftAudit.lean} (standard axioms; a definition of a proposition has no proof content).
\textbf{Not proved in this library.}  It is AxiomMath's statement after renaming (\lean{phi0} restated definitionally equal),
proved in AxiomMath's library (Lean~4.35, commit \texttt{69cb972f}); the literature milestones on its proof path that were
inspected (\lean{cl_lit_stroth_leon_sims_baby_monster_recognition}, its \lean{_models} variant, and
\lean{cl_lit_goldschmidt_strongly_closed_abelian}) are theorems there, not axioms; its full axiom footprint has not been
re-checked (that needs \texttt{\#print axioms} on a Lean~4.35 build of AxiomMath's library).
\end{status}

\begin{warning}
\textbf{Minor inconsistency in the recorded closure size.}  The docstring of \lean{GmatEqAutO}'s file says ``about 15,500
files''; the project's records say ``15.5k files''; Chapter~\ref{ch:16} says ``about $16{,}000$ files and $208$\,MB''.  All are
consistent with the following: $15{,}502$ CFSG files ($196.9$\,MB) for the minimal roots, plus $859$ TauCeti
files; $16{,}361$ files and $208$\,MB at file level.  No defect; recorded so that a reader does not mistake the figures for a
discrepancy.
\end{warning}

\subsection{The isomorphism}

\begin{theorem}[$\chtwentytwoAutO\cong\Aut(\Vnat)$, given $\mathrm{GmatEqAutO}$]\label{ch22:thm:equiv}
Assume $\chtwentytwoGmat=\chtwentytwoAutO$.  Then every element of $\chtwentytwoAutO$ is liftable, the lift
$\chtwentytwoAutO\to\Aut(\Vnat)$ is bijective, and the resulting group isomorphism $\varphi$ satisfies
$\varphi(g)|_{\Vnat_2}=\mathcal G\circ(1_\C\otimes g)\circ\mathcal G^{-1}$ for every $g$, $\mathcal G=$ \lean{griessIso}.
\end{theorem}
\begin{leancode}
theorem isVnatLiftable_of_gmatEqAutO (hM : GmatEqAutO) (g : AutO) : IsVnatLiftable g
-- ⋯
theorem monsterLiftOfAll_bijective (hM : GmatEqAutO) :
    Function.Bijective (monsterLiftOfAll (isVnatLiftable_of_gmatEqAutO hM))
-- ⋯
noncomputable def monsterEquivAutVnat (hM : GmatEqAutO) : AutO ≃* AutVnat :=
  MulEquiv.ofBijective (monsterLiftOfAll (isVnatLiftable_of_gmatEqAutO hM))
    (monsterLiftOfAll_bijective hM)
-- ⋯
theorem monsterEquivAutVnat_restrict (hM : GmatEqAutO) (g : AutO) :
    vnatWeightTwoRestrict (monsterEquivAutVnat hM g) =
      griessConj (g : SpaceB ≃ₗ[ℚ] SpaceB) :=
  monsterLiftOfAll_restrict (isVnatLiftable_of_gmatEqAutO hM) g
\end{leancode}
(\leanpath{VOA/Moonshine/MoonshineLRMonsterHeadline.lean}.)
\begin{dictionary}
\textbf{Decoding.}  \lean{monsterLiftOfAll hlift : AutO →* AutVnat} (\S\ref{ch22:sec:recall}) is \lean{monsterLift} composed with
the inclusion of $\chtwentytwoAutO$ into the liftable subgroup, available once every element is liftable;
\lean{Function.Bijective} is injective and surjective; \lean{MulEquiv.ofBijective} turns a bijective homomorphism into a group
isomorphism \lean{AutO ≃* AutVnat}.  In \lean{monsterEquivAutVnat_restrict}, \lean{monsterEquivAutVnat hM g} is $\varphi(g)$ (the
coercion of a \lean{MulEquiv} to a function), and the right side is the Griess conjugate (quoted in \S\ref{ch22:sec:recall}),
so the identity says: $\varphi(g)$ acts on $\Vnat_2\cong\C\otimes\chtwentytwoB$ as $1\otimes g$.  The proof term
\lean{monsterLiftOfAll_restrict …} is shown because it is informative: the compatibility holds by construction.

\textbf{Hypothesis.}  Only \lean{hM : GmatEqAutO}.  Injectivity of the lift is unconditional (\lean{monsterLiftOfAll_injective});
surjectivity uses the unconditional converse (\lean{monsterLift_surjective_of_commutant M8.gmatCommutantScalar}, visible in the
proof of \lean{monsterLiftOfAll_bijective}); \lean{hM} is used only to make every element of $\chtwentytwoAutO$ liftable
(\lean{isVnatLiftable_of_gmatEqAutO}: an element of $\chtwentytwoAutO$ is in $\chtwentytwoGmat$, hence liftable by
Theorem~\ref{ch22:thm:gmat}(b)).

\textbf{Uniqueness, i.e.\ why the isomorphism is canonical.}  An automorphism of $\Vnat$ is determined by its restriction to
$\Vnat_2$ (\lean{vnatWeightTwoRestrict_injective}), so $\varphi$ is the \emph{only} map $\chtwentytwoAutO\to\Aut(\Vnat)$ with
the compatibility property.

\textbf{Strength.}  This is the expected statement ``the Monster acts faithfully on $\Vnat$ and every automorphism of $\Vnat$
comes from it'', for AxiomMath's model of the Monster, compatibly with the identification of Griess algebras.
\end{dictionary}
\begin{proof}[Proof on paper]
Under \lean{hM}, $\chtwentytwoAutO=\chtwentytwoGmat\subseteq\mathrm{liftable}$ (Theorem~\ref{ch22:thm:gmat}); \lean{monsterLift}
is injective (Chapter~\ref{ch:21}) and surjective onto $\Aut(\Vnat)$ (Theorem~\ref{ch22:thm:converse}).
\end{proof}
\begin{physics}
\textbf{The Monster is exactly the symmetry group of the moonshine CFT}, realised on the $196884$ weight-two fields by
AxiomMath's explicit rational matrices --- conditional on the group-theoretic identification.
\end{physics}
\begin{status}
\textbf{Conditional} on \lean{GmatEqAutO}.  \texttt{\#print axioms} for all four in \leanpath{VOA/Moonshine/MoonshineLRLiftAudit.lean}:
\lean{[propext, Classical.choice, Quot.sound]} (rebuild of the audit, 9,457 jobs, 926 reports, $0$ \lean{sorryAx}).  \texttt{\#print axioms} does not report \lean{hM}.
\end{status}

\subsection{Monstrous moonshine}

\begin{theorem}[Monstrous moonshine for the Monster; Conway--Norton, Borcherds]\label{ch22:thm:monstrous}
Assume $\chtwentytwoGmat=\chtwentytwoAutO$.  Then there is a group isomorphism $\varphi:\chtwentytwoAutO\to\Aut(\Vnat)$ such that
\begin{enumerate}[label=(\roman*)]
\item for every $g\in\chtwentytwoAutO$, $\varphi(g)$ acts on $\Vnat_2\cong\C\otimes_\Q\chtwentytwoB$ as $1\otimes g$; and
\item for every $g\in\chtwentytwoAutO$ there are a subgroup $\Gamma\le\mathrm{SL}_2(\mathbb{R})$ (in Lean: a subgroup of
  $\mathrm{GL}_2(\mathbb{R})$ all of whose elements have determinant $1$) and $N\ge1$ with $\Gamma_0(N)\le\Gamma$, $\Gamma$ commensurable with $\mathrm{SL}_2(\Z)$, such that
  \begin{equation*}
    T_{\varphi(g)}(\tau)=q^{-1}+\sum_{n\ge1}\tr\bigl(\varphi(g)\big|\Vnat_{n+1}\bigr)q^n,\qquad q=e^{2\pi i\tau},
  \end{equation*}
  is a normalised Hauptmodul for $\Gamma$.
\end{enumerate}
\end{theorem}
\begin{leancode}
theorem monstrousMoonshine (hM : GmatEqAutO) :
    ∃ φ : AutO ≃* AutVnat,
      (∀ g : AutO, vnatWeightTwoRestrict (φ g) = griessConj (g : SpaceB ≃ₗ[ℚ] SpaceB)) ∧
        ConwayNortonVia (φ : AutO →* AutVnat) :=
  ⟨monsterEquivAutVnat hM, monsterEquivAutVnat_restrict hM, fun _ => conwayNorton _⟩
\end{leancode}
(\leanpath{VOA/Moonshine/MoonshineLRMonsterHeadline.lean}; full name \chtwentytwolean{VOA.Moonshine.ConwayNorton.monstrousMoonshine}.)

\begin{dictionary}
\textbf{Clause by clause.}
\begin{enumerate}[label=(\arabic*)]
\item \lean{(hM : GmatEqAutO)}: the one hypothesis, Definition~\ref{ch22:def:GmatEqAutO}.
\item \lean{∃ φ : AutO ≃* AutVnat}: a group isomorphism from AxiomMath's Monster onto the automorphism group of the vertex
  operator algebra \lean{VnatVOA} (Chapter~\ref{ch:16}).  \lean{AutO} is a subgroup, used here as a group (its elements are
  pairs of a linear automorphism and a membership proof; \lean{(g : SpaceB ≃ₗ[ℚ] SpaceB)} is the underlying map).
\item \lean{∀ g : AutO, vnatWeightTwoRestrict (φ g) = griessConj (g : SpaceB ≃ₗ[ℚ] SpaceB)}: clause (i).  Both sides are
  $\C$-linear automorphisms of $\Vnat_2$; the right side is $\mathcal G\circ(1\otimes g)\circ\mathcal G^{-1}$.  By
  \lean{vnatWeightTwoRestrict_injective}, this clause determines $\varphi$ uniquely, so the existential is not a loophole:
  it is witnessed by \lean{monsterEquivAutVnat hM} and by nothing else.
\item \lean{ConwayNortonVia (φ : AutO →* AutVnat)}: by definition (\S\ref{ch22:sec:recall}),
  \lean{∀ g : AutO, ConwayNortonFor (φ g)}, i.e.\ clause (ii): \lean{IsMoonshineGroup Γ N} encodes $N\ge1$,
  $\Gamma\le\mathrm{SL}_2(\mathbb{R})$ (\lean{HasDetOne}), $\Gamma_0(N)\le\Gamma$ (\lean{gamma0GL N ≤ Γ}) and commensurability with
  $\mathrm{SL}_2(\Z)$; \lean{IsNormalizedHauptmodul Γ (mckayThompsonFun (φ g))} is the eight-clause Hauptmodul condition of
  Chapter~\ref{ch:16} for $T_{\varphi(g)}$.  The coercion \lean{(φ : AutO →* AutVnat)} forgets that $\varphi$ is invertible.
\end{enumerate}

\textbf{Comparison with the textbook statement.}  The Conway--Norton conjecture (proved by Borcherds) says: there is a graded
faithful representation $V=\bigoplus V_n$ of $\chtwentytwoM$ with $\dim V_n$ the coefficients of $J$, such that for every
$g\in\chtwentytwoM$, $T_g$ is the normalised Hauptmodul of a genus-zero group $\Gamma_g$ containing $\Gamma_0(N)$.  The Lean
statement has: \emph{the Monster} $=$ AxiomMath's \lean{AutO} (a \emph{name}: LeanVOA does not prove that \lean{AutO} is simple,
nor its order $2^{46}3^{20}5^9\cdots71\approx8\cdot10^{53}$; these are classical facts about the Monster, and AxiomMath's library
proves finiteness and order results for its model on Lean~4.35, none of which is imported here); \emph{the representation} $=\Vnat$, through $\varphi$ (the graded
dimension $J$ and simplicity of $\Vnat$ are Chapter~\ref{ch:14}'s \lean{moonshine}); \emph{faithful} $=\varphi$ injective; and
a property not in the conjecture but in FLM's theorem: the action is by vertex-operator-algebra automorphisms and is
\emph{all} of $\Aut(\Vnat)$, compatibly with the Griess algebras.  It is \emph{weaker} than the full Conway--Norton table in one
recorded respect: $\Gamma$ and $N$ are existentially quantified; their identification with the groups of the Conway--Norton
list (and $N=o(g)h$) is not stated (Chapter~\ref{ch:16}).  It is \emph{stronger} than ``moonshine for some finite group'' in that
the group is pinned to AxiomMath's independent construction.

\textbf{What the statement would mean if \lean{GmatEqAutO} were false.}  Then the theorem would be vacuous (an implication
with a false premise), and would say nothing about the Monster.  But every \emph{unconditional} result of the part stands, and
together they still say a great deal:
\begin{itemize}
\item \lean{conwayNorton}: Conway--Norton for every automorphism of $\Vnat$, whatever that group is;
\item \lean{conwayNorton_gmat}: every element of the explicit group $\chtwentytwoGmat$ acts on $\Vnat$ compatibly with the
  Griess algebras, and its McKay--Thompson series is a Hauptmodul;
\item \lean{nonempty_autVnat_mulEquiv_liftable}: $\chtwentytwoGmat\le\Aut(\Vnat)\le\chtwentytwoAutO$ (canonically);
\item \lean{finite_autVnat}, \lean{rationalTraces}, \lean{gmatCommutantScalar}.
\end{itemize}
So the only thing \lean{GmatEqAutO} adds is the identification of the sandwiched group $\Aut(\Vnat)$ with the outer group
$\chtwentytwoAutO$.

\textbf{What release v1.1 changes} (the authors' plan: ``v1.1 will then bump the library to 4.35 and
discharge it'').  After moving the library to AxiomMath's toolchain (Lean 4.35 and AxiomMath's Mathlib pin) and importing
AxiomMath's library, \lean{gmat_eq_griess} becomes available; with a bridge identifying the ported definitions with the
imported ones (by \lean{rfl}, given the fidelity audit; or by re-porting \lean{phi0} in bundled form), the hypothesis can be
discharged and \lean{monstrousMoonshine} instantiated to an unconditional theorem --- unconditional relative to AxiomMath's
library, whose own axiom footprint must then be checked.  This is a plan; none of it is in v1.05.
\end{dictionary}

\begin{proof}[Proof on paper]
Take $\varphi=$ \lean{monsterEquivAutVnat hM}; (i) is Theorem~\ref{ch22:thm:equiv}; (ii) is \lean{conwayNorton} (Chapter~\ref{ch:20})
applied to $\varphi(g)$.  The Lean proof is the anonymous-constructor term
\lean{⟨monsterEquivAutVnat hM, monsterEquivAutVnat_restrict hM, fun _ => conwayNorton _⟩}, displayed above.
\end{proof}

\begin{physics}
\textbf{Monstrous moonshine as a CFT statement.}  The moonshine CFT has the Monster as its full global symmetry group, and the
torus partition function twisted by any Monster element $g$ is the Hauptmodul of a genus-zero modular curve.  In string-theory
language (Chapters~\ref{ch:17}--\ref{ch:19}): compactifying the bosonic string on $\Vnat\otimes V_{I\!I_{1,1}}$ produces the
Monster Lie algebra as the BRST cohomology; the Monster acts on it; its twisted denominator identities express that the
twisted partition functions are completely replicable, and complete replicability forces genus zero.  The present theorem
adds the identification of the symmetry group with an independently constructed Monster.
\end{physics}

\begin{status}
\textbf{Conditional on the named hypothesis \lean{GmatEqAutO}} and on nothing else.  \texttt{\#print axioms}
(\leanpath{VOA/Moonshine/MoonshineLRLiftAudit.lean}): \lean{[propext, Classical.choice, Quot.sound]}; $0$ \lean{sorryAx}; wired
into the umbrella \lean{VOA.Moonshine} (umbrella green, 9,504 jobs).  The aggregate
\texttt{lake build} of the whole library for this state is the release step and is not reported in the sources read for this
chapter.  This book never calls \lean{monstrousMoonshine} unconditional.
\end{status}


%======================================================================
\section{Documentation defects recorded while writing this chapter}\label{ch22:sec:docdefects}
%======================================================================

The following module docstrings and notes describe an earlier state of the work; none affects a Lean statement.  They
are collected here rather than interrupting the mathematics, and are classed with severity S4 (the lowest: documentation only).
\begin{warning}
\begin{itemize}
\item \textbf{Stale module docstrings.}  The module docstrings of \leanpath{VOA/Moonshine/MoonshineLRLiftTwisted.lean} (``The
candidate on \lean{TP} is not yet proved to intertwine \lean{plusModule}'') and
\leanpath{VOA/Moonshine/MoonshineLRLiftTwisted2.lean} (``we have not proved that a scalar multiple preserves
\lean{TPForm}'') describe the state of each file in isolation; both gaps are closed in the next file
(\lean{liftTwPlusEquiv_mode}, \lean{liftTwExists_formScalar}).  The definition \lean{liftTwVnatAut} has no form
hypothesis.  (Recorded as an S4 documentation defect.)
\item \textbf{Stale docstrings.}  \leanpath{VOA/Moonshine/MoonshineLRLiftLiftable2.lean} says ``The other Golay sign changes
and the Mathieu permutations remain open''; \leanpath{VOA/Moonshine/MoonshineLRLiftLiftable3.lean} and
\leanpath{VOA/Moonshine/MoonshineLRLiftLiftable4.lean} say ``The generator $\xi$ is not treated''.  These describe the
files in isolation; the gaps are closed by Theorem~\ref{ch22:thm:monomial} and Theorem~\ref{ch22:thm:xi}.  (S4.)
\item \textbf{Stale docstring.}  \leanpath{VOA/Moonshine/MoonshineLRLiftVisible.lean} says ``The remaining gap is to identify
and cancel that discrepancy as an extraspecial action''; it is closed in
\leanpath{VOA/Moonshine/MoonshineLRLiftVisible2.lean}.  (S4.)
\item \textbf{Stale docstrings.}  \leanpath{VOA/Moonshine/MoonshineLRSigmaBlock.lean} (``The full rational dictionary matrix of
\lean{v0Sigma} still requires matching \dots''), \leanpath{VOA/Moonshine/MoonshineLRSigmaBlock2.lean} (``The octad
\lean{KS} coefficients are not treated'') and \leanpath{VOA/Moonshine/MoonshineLRSigmaB.lean} (``isolates the unproved
input'') describe files in isolation; the gaps are closed in \leanpath{VOA/Moonshine/MoonshineLRSigmaBlock3.lean} and
\leanpath{VOA/Moonshine/MoonshineLRSigmaB2.lean}.  (S4.)
\item \textbf{Stale docstring.}  The module docstring of \leanpath{VOA/Moonshine/MoonshineLRTauB.lean} says ``The final
theorem retains only \lean{tauB ∈ AutO} and \lean{Generated.Gmat = AutO}''.  That is true of the file; the first of the two
was discharged afterwards (\lean{tauB_mem_AutO}, ported from the AxiomMath library), and
\leanpath{VOA/Moonshine/MoonshineLRMonsterFinal.lean} carries only the second.  (S4.)
\item \textbf{Route deviation from the paper plan.}  The authors' paper plan (Step 7) planned to use the Golay sign characters on
$\chtwentytwoA$ to show $T(D)\subseteq D$; the Lean proof obtains it from triality directly (Step 7 above), so the sign-character
step was not needed.  (S4, note.)
\item \textbf{Stale note.}  An earlier working note says ``The opposite direction
$\mathrm{AutVnat}\le$ Monster (rationality, Tits) is not needed for this statement'' and a later entry calls it
``open research''; both predate Theorem~\ref{ch22:thm:converse}, which proves it (without Tits).  (S4, note.)
\item \textbf{Namespace.}  The \texttt{MoonshineLR*} files use the namespace \lean{VOA.Moonshine.ConwayNorton.M8}, whose suffix is a
  historical node label (recorded as naming debt).
\end{itemize}
\end{warning}

%======================================================================
\section{The assigned files and how each is covered}\label{ch22:sec:filemap}
%======================================================================

The rules of \texttt{Notes/Textbook/assignment.txt} assign to this chapter the $39$ files
\texttt{VOA/\allowbreak Moonshine/\allowbreak MoonshineLR\{Lift*,\allowbreak Sigma*,\allowbreak TauB,\allowbreak MonsterVisible,\allowbreak MonsterFinal,\allowbreak Commutant*,\allowbreak MonsterHeadline\}.lean} and
\texttt{VOA/\allowbreak Moonshine/\allowbreak MoonshineCN\{MonsterFinalAudit,\allowbreak SymmetryFinalAudit,\allowbreak Rational*\}.lean}.  ``Individually'' means the
headline declarations are presented under the Fidelity Protocol; ``collectively'' means the file's content is described and
its role explained.

\begin{center}\small
\begin{longtable}{p{0.36\textwidth}p{0.18\textwidth}p{0.38\textwidth}}
\toprule
file (\texttt{VOA/Moonshine/}) & coverage & where; content\\
\midrule
\endhead
\texttt{MoonshineLRLiftNatural.lean} & individually & Thm~\ref{ch22:thm:natural}; naturality \lean{vnatAutOfPair}\\
\texttt{MoonshineLRLiftLattice.lean} & individually & Def~\ref{ch22:def:latplus}; $\chtwentytwoCo_0$ on $V_\Leech^+$\\
\texttt{MoonshineLRLiftTwisted.lean} & collectively & proof of Thm~\ref{ch22:thm:twcov}; Fock frame change, tensor candidate, coboundary relations\\
\texttt{MoonshineLRLiftTwisted2.lean} & individually & Thms~\ref{ch22:thm:rho}, \ref{ch22:thm:twcov}; $\rho$, word-field covariance\\
\texttt{MoonshineLRLiftTwisted3.lean} & individually & Thm~\ref{ch22:thm:twcov}, Def~\ref{ch22:def:liftTw}; form scalar, \lean{liftTwVnatAut}\\
\texttt{MoonshineLRLiftExtra.lean} & individually & before Def~\ref{ch22:def:extra}; character automorphism, $e_\mu$, form multiplier\\
\texttt{MoonshineLRLiftExtraAut.lean} & individually & Def~\ref{ch22:def:extra}; \lean{extraVnatAut}\\
\texttt{MoonshineLRLiftBlock.lean} & individually & Thm~\ref{ch22:thm:untwistedblock}; $A\mapsto gAg^{\mathsf T}$, extraspecial blocks\\
\texttt{MoonshineLRLiftSign.lean} & individually & Thm~\ref{ch22:thm:untwistedblock}; signed permutation, $r_c=\pm1$\\
\texttt{MoonshineLRLiftWBlock.lean} & collectively & \S\ref{ch22:sec:blocks} (twisted block); spin coordinates of the lifts\\
\texttt{MoonshineLRLiftLiftable.lean} & individually & Thm~\ref{ch22:thm:criterion}(a),(c); rational criterion\\
\texttt{MoonshineLRLiftLiftable2.lean} & individually & Thm~\ref{ch22:thm:criterion}(b),(c); phases, unphased equation, $-1$\\
\texttt{MoonshineLRLiftLiftable3.lean} & individually & Def~\ref{ch22:def:monomial}, Thm~\ref{ch22:thm:criterion}; generators, basis reduction\\
\texttt{MoonshineLRLiftLiftable4.lean} & individually & Thm~\ref{ch22:thm:monomial}; $2^{12}{:}M_{24}$ liftable\\
\texttt{MoonshineLRLiftXi.lean} & individually & Def~\ref{ch22:def:xi}, Thm~\ref{ch22:thm:xi}; $\xi$, rational descent\\
\texttt{MoonshineLRLiftVisible.lean} & collectively & proof of Thm~\ref{ch22:thm:visible}; visible-pair dictionary, \lean{liftVisParity_ext}\\
\texttt{MoonshineLRLiftVisible2.lean} & individually & Thm~\ref{ch22:thm:visible}; \lean{liftVisRAllVisible}\\
\texttt{MoonshineLRMonsterVisible.lean} & individually & Thm~\ref{ch22:thm:visible}; \lean{conwayNorton_autO_of_tauB}\\
\texttt{MoonshineLRSigmaBlock.lean} & individually & Thm~\ref{ch22:thm:sector}; sector split, $\sigma z\sigma$\\
\texttt{MoonshineLRSigmaBlock2.lean} & individually & Thm~\ref{ch22:thm:sigmaeven}; two-root block\\
\texttt{MoonshineLRSigmaBlock3.lean} & individually & Thm~\ref{ch22:thm:sigmaeven}; octad block, \lean{sigmaBlockOFull_nat}\\
\texttt{MoonshineLRSigmaB.lean} & individually & Thm~\ref{ch22:thm:sigmarational}(a); reduction to odd columns\\
\texttt{MoonshineLRSigmaB2.lean} & individually & Thm~\ref{ch22:thm:sigmarational}(b),(c); seed, \lean{sigmaBOdd_liftable}\\
\texttt{MoonshineLRTauB.lean} & individually & Thm~\ref{ch22:thm:tauB}; conjugation, centralizer\\
\texttt{MoonshineLRMonsterFinal.lean} & individually & Thms~\ref{ch22:thm:tauB}, \ref{ch22:thm:cnmonster}\\
\texttt{MoonshineLRCommutant.lean} & collectively & \S\ref{ch22:sec:commutant}, Steps 1--3\\
\texttt{MoonshineLRCommutant2.lean} & collectively & \S\ref{ch22:sec:commutant}, Steps 4--7\\
\texttt{MoonshineLRCommutant3.lean} & collectively (quoted) & \S\ref{ch22:sec:commutant}, Steps 8--9; \lean{commutL_xi_blocks}, $A$, $B$\\
\texttt{MoonshineLRCommutant4.lean} & individually & Thm~\ref{ch22:thm:commutant}, Steps 9--10; certificates\\
\texttt{MoonshineLRCommutantFinal.lean} & individually & Thms~\ref{ch22:thm:commutant}, \ref{ch22:thm:converse}\\
\texttt{MoonshineLRMonsterHeadline.lean} & individually & \S\ref{ch22:sec:headline}; the headline\\
\texttt{MoonshineCNMonsterFinalAudit.lean} & individually + audit & Thms~\ref{ch22:thm:gmat}, \ref{ch22:thm:cnmonster}; \S\ref{ch22:sec:audits}\\
\texttt{MoonshineCNSymmetryFinalAudit.lean} & individually + audit & Thm~\ref{ch22:thm:converse}; \S\ref{ch22:sec:audits}\\
\texttt{MoonshineCNRationalTraces.lean} & individually & Thms~\ref{ch22:thm:galois}, \ref{ch22:thm:coeff}, \ref{ch22:thm:rationaltraces}\\
\texttt{MoonshineCNRationalRamanujan.lean} & individually & Thm~\ref{ch22:thm:ramanujan}\\
\texttt{MoonshineCNRationalSchur.lean} & individually & Thms~\ref{ch22:thm:schurspan}, \ref{ch22:thm:descent}\\
\texttt{MoonshineCNRationalGriess.lean} & individually & Def~\ref{ch22:def:P0}--Thm~\ref{ch22:thm:rationalcomm}\\
\texttt{MoonshineCNRationalAudit.lean} & audit & \S\ref{ch22:sec:audits}\\
\texttt{MoonshineLRLiftAudit.lean} & audit & \S\ref{ch22:sec:audits}\\
\bottomrule
\end{longtable}
\end{center}

\paragraph{Supporting declarations not displayed.}  The files contain many lemmas not quoted above (in total several hundred);
each is used in one of the proofs described.  In particular: the Fock-space covariance lemmas of
\leanpath{VOA/Moonshine/MoonshineLRLiftTwisted.lean} (annihilator chain rule, creation and annihilation exponentials, dressing
\lean{DD}); the field-product transport of \leanpath{VOA/Moonshine/MoonshineLRLiftTwisted2.lean} (\lean{liftTwConjField_prod},
\lean{liftTwConjField_star}); the rank-transport lemmas (\lean{liftTwCastEquiv_YM}, \lean{extraPhi_cast_YM}) that move between
the abstract lattice rank and the concrete $24$; the spin-coordinate data of \leanpath{VOA/Moonshine/MoonshineLRLiftWBlock.lean}
(\lean{liftWEquiv}, a complex coordinate isomorphism $(\C^{4096})^{24}\cong(T^+)_2$; the module docstring says, correctly, that
these statements do not assert rationality); and the sign bookkeeping lemmas of the triality files.

%======================================================================
\section{Audit files}\label{ch22:sec:audits}
%======================================================================

\begin{itemize}
\item \leanpath{VOA/Moonshine/MoonshineLRLiftAudit.lean}: at the release, $997$ \texttt{\#print axioms}
  commands (counted in the file) and $42$ imports: all \texttt{MoonshineLRLift*}, \texttt{MoonshineLRSigma*},
  \texttt{MoonshineLRTauB}, \texttt{MoonshineLRMonster\{Visible,Final,Headline\}}, \texttt{MoonshineLRCommutant*},
  \texttt{MoonshineLRGriessTransfer}, and the \texttt{Aut*.lean} files of
  Chapter~\ref{ch:20} (whose $124$ theorems were added).  Build record: green, $9{,}457$
  jobs for the headline state with $926$ reports; at the release, green, $997$ reports; all
  \lean{[propext, Classical.choice, Quot.sound]}, $0$ \lean{sorryAx}.
\item \leanpath{VOA/Moonshine/MoonshineCNRationalAudit.lean}: $94$ \texttt{\#print axioms} commands (plus one \texttt{\#print}
  of the definition \lean{GmatCommutantScalar}) over
  \texttt{MoonshineCNRational\{Schur,Ramanujan,Traces,Griess\}}; green ($9{,}449$ jobs), standard axioms, $0$ \lean{sorryAx}.  Its
  docstring notes that \texttt{\#print axioms} does not report \lean{hcomm} or \lean{hGmat}.
\item \leanpath{VOA/Moonshine/MoonshineCNMonsterFinalAudit.lean}: $11$ reports (\lean{conwayNorton_monster},
  \lean{conwayNorton_gmat}, \lean{gmat_eq_map_closure_insert_tauB}, \lean{isVnatLiftable_of_mem_gmat}, \lean{conwayNorton},
  \lean{M8.isVnatLiftable_tauB}, \lean{M8.conwayNorton_monster_of_Gmat}, \lean{tauB_mem_AutO}, \lean{Generated.Gmat_le_AutO},
  \lean{griessIso_mul}, \lean{griessIso_form}); green ($9{,}441$ jobs), standard axioms.
\item \leanpath{VOA/Moonshine/MoonshineCNSymmetryFinalAudit.lean}: $8$ reports (\lean{autVnat_rational},
  \lean{nonempty_autVnat_mulEquiv_liftable}, \lean{autVnat_mulEquiv_autO}, \lean{M8.gmatCommutantScalar},
  \lean{M8.autVnat_mulEquiv_liftable}, \lean{rationalTraces}, \lean{conwayNorton}, \lean{conwayNorton_gmat}); green
  ($9{,}455$ jobs), standard axioms.
\end{itemize}
These figures are those of the build record of the release (Chapter~\ref{ch:15}, Section~\ref{ch15:sec:record}).  The two
\texttt{*FinalAudit} files also contain theorems, presented above.

%======================================================================
\section{Referee's checklist}\label{ch22:sec:checklist}
%======================================================================

\subsection{This chapter}

\begin{center}\small
\begin{longtable}{p{0.31\textwidth}p{0.33\textwidth}p{0.17\textwidth}p{0.09\textwidth}}
\toprule
declaration & intended meaning & hypotheses & check\\
\midrule
\endhead
\chtwentytwolean{M8.vnatAutOfPair} & compatible form-preserving pair is an automorphism of $\Vnat$ & \lean{hφ}, \lean{hT} (data) & Thm~\ref{ch22:thm:natural}\\
\chtwentytwolean{M8.leechPlusLatMorphAut} & $\chtwentytwoCo_0$ on $V_\Leech^+$ & isometry, onto & Def~\ref{ch22:def:latplus}\\
\chtwentytwolean{M8.liftTwChargeIntertwiner_exists} & charge intertwiner $\rho$ exists & isometry & Thm~\ref{ch22:thm:rho}\\
\chtwentytwolean{M8.liftTwPlusEquiv_mode}, \chtwentytwolean{M8.liftTwExists_formScalar} & $\varphi_\iota$ covariant; normalisable & isometry, onto & Thm~\ref{ch22:thm:twcov}\\
\chtwentytwolean{M8.liftTwVnatAut} & lift of $\iota\in\chtwentytwoCo_0$ to $\Aut(\Vnat)$ (chosen) & isometry, onto & Def~\ref{ch22:def:liftTw}\\
\chtwentytwolean{M8.extraVnatAut} & lift of $\mu\in\Leech$ (extraspecial) & none & Def~\ref{ch22:def:extra}\\
\chtwentytwolean{M8.liftSignNat_untwisted} & untwisted block: $gAg^{\mathsf T}$, signed permutation & isometry, onto & Thm~\ref{ch22:thm:untwistedblock}\\
\chtwentytwolean{M8.liftLiftableIffSpinRational}, \chtwentytwolean{M8.liftMonoOfUnphasedIntertwiner} & rational criterion for liftability & rational orthogonal intertwiner & Thm~\ref{ch22:thm:criterion}\\
\chtwentytwolean{M8.liftSpinGolay_liftable}, \chtwentytwolean{M8.liftSpinMathieu_liftable} & $2^{12}{:}M_{24}$ lifts & codeword / code-preserving & Thm~\ref{ch22:thm:monomial}\\
\chtwentytwolean{M8.liftXi_liftable} & $\xi$ lifts & none & Thm~\ref{ch22:thm:xi}\\
\chtwentytwolean{M8.liftVisRAllVisible} & every element of $\chtwentytwoGvis$ lifts & none & Thm~\ref{ch22:thm:visible}\\
\chtwentytwolean{M8.conwayNorton_autO_of_tauB} & CN for $\chtwentytwoAutO$ from $\tau$ and generation & \lean{hτ}, \lean{hGmat} & Thm~\ref{ch22:thm:visible}\\
\chtwentytwolean{M8.sigmaBlockConjugate} & $\sigma z\sigma=$ frame sign & none & Thm~\ref{ch22:thm:sector}\\
\chtwentytwolean{M8.sigmaBlockOFull_nat} & $\sigma_0$ rational & even support & Thm~\ref{ch22:thm:sigmaeven}\\
\chtwentytwolean{M8.sigmaBOdd_rational}, \chtwentytwolean{M8.sigmaBOdd_liftable} & $\sigma$ rational; $\sigma=\mathrm{monsterLift}(\sigma_B)$ & none & Thm~\ref{ch22:thm:sigmarational}\\
\chtwentytwolean{M8.tauBLiftSigma_conj_xB} & $\sigma_Bx_B\sigma_B^{-1}\in\{y_B,z_B\}$ & \lean{hτ} (a theorem) & Thm~\ref{ch22:thm:tauB}\\
\chtwentytwolean{M8.isVnatLiftable_tauB} & $\tau_B$ lifts & none & Thm~\ref{ch22:thm:tauB}\\
\chtwentytwolean{conwayNorton_gmat} & CN for every $g\in\chtwentytwoGmat$, with lift & none beyond membership $g\in\chtwentytwoGmat$ & Thm~\ref{ch22:thm:gmat}\\
\chtwentytwolean{conwayNorton_monster}, \chtwentytwolean{M8.conwayNorton_monster_of_Gmat} & CN for all of $\chtwentytwoAutO$ & \lean{hGmat} & Thm~\ref{ch22:thm:cnmonster}\\
\chtwentytwolean{mckayThompsonFun_pow_eq_of_coprime_holds} & $T_{g^u}=T_g$ & none & Thm~\ref{ch22:thm:galois}\\
\chtwentytwolean{Rationality.exists_rat_trace_of_trace_pow_coprime} & Galois-invariant finite-order trace is rational & as stated & Thm~\ref{ch22:thm:ramanujan}\\
\chtwentytwolean{rationalTraces} & all traces rational & none & Thm~\ref{ch22:thm:rationaltraces}\\
\chtwentytwolean{Rationality.schur_span_formula}, \chtwentytwolean{Rationality.eq_map_ratCast_of_rational_traces} & Schur span; descent & as stated & Thms~\ref{ch22:thm:schurspan}, \ref{ch22:thm:descent}\\
\chtwentytwolean{MonsterAction.finite_gmat} & $\chtwentytwoGmat$ finite & none & Thm~\ref{ch22:thm:omega}\\
\chtwentytwolean{MonsterAction.GmatCommutantScalar} & commutant of $\chtwentytwoGmat$ is $\Q1+\Q P_0$ & (a \lean{Prop}) & Def~\ref{ch22:def:commutant}\\
\chtwentytwolean{MonsterAction.autVnat_rational_of_commutant} & $\Aut(\Vnat)$ rational given commutant & \lean{hcomm} & Thm~\ref{ch22:thm:rationalcomm}\\
\chtwentytwolean{M8.gmatCommutantScalar} & the commutant lemma & none & Thm~\ref{ch22:thm:commutant}\\
\chtwentytwolean{autVnat_rational} & every automorphism of $\Vnat$ is rational on $\Vnat_2$ & none & Thm~\ref{ch22:thm:converse}\\
\chtwentytwolean{nonempty_autVnat_mulEquiv_liftable}, \chtwentytwolean{M8.autVnat_mulEquiv_liftable} & $\Aut(\Vnat)\cong$ liftable subgroup of $\chtwentytwoAutO$ & none & Thm~\ref{ch22:thm:converse}\\
\chtwentytwolean{autVnat_mulEquiv_autO} & $\Aut(\Vnat)\cong\chtwentytwoAutO$ & \lean{hGmat} & Thm~\ref{ch22:thm:converse}\\
\chtwentytwolean{GmatEqAutO} & $\chtwentytwoGmat=\chtwentytwoAutO$ (AxiomMath's statement after renaming, \lean{phi0} restated definitionally equal) & (a \lean{Prop}) & Def~\ref{ch22:def:GmatEqAutO}\\
\chtwentytwolean{monsterEquivAutVnat}, \chtwentytwolean{monsterEquivAutVnat_restrict} & $\chtwentytwoAutO\cong\Aut(\Vnat)$, Griess-compatible & \lean{GmatEqAutO} & Thm~\ref{ch22:thm:equiv}\\
\chtwentytwolean{monstrousMoonshine} & monstrous moonshine for AxiomMath's Monster & \textbf{\lean{GmatEqAutO}} & Thm~\ref{ch22:thm:monstrous}\\
\bottomrule
\end{longtable}
\end{center}
Names without a prefix are in \chtwentytwolean{VOA.Moonshine.ConwayNorton}; \lean{M8.} abbreviates \chtwentytwolean{VOA.Moonshine.ConwayNorton.M8.};
\lean{MonsterAction.} and \lean{Rationality.} are in \lean{VOA.Moonshine}.

\subsection{The part's final table}

The headline statements of Chapters~\ref{ch:16}--\ref{ch:22}, with their hypotheses.  ``None'' means: no explicit hypothesis, and
\texttt{\#print axioms} is \lean{[propext, Classical.choice, Quot.sound]}.

\begin{center}\small
\begin{longtable}{p{0.34\textwidth}p{0.38\textwidth}p{0.2\textwidth}}
\toprule
declaration & meaning & hypotheses; chapter\\
\midrule
\endhead
\chtwentytwolean{ConwayNorton.conwayNorton} & Conway--Norton for every automorphism of $\Vnat$ & none; Ch.~\ref{ch:16}, \ref{ch:20}\\
\chtwentytwolean{ConwayNorton.conwayNortonGammaOne_holds} & the $\Gamma_1$ form & none; Ch.~\ref{ch:19}, \ref{ch:22}\\
\chtwentytwolean{ConwayNorton.M8.finite_autVnat} & $\Aut(\Vnat)$ finite & none; Ch.~\ref{ch:20}\\
\chtwentytwolean{ConwayNorton.rationalTraces} & all McKay--Thompson coefficients rational & none; Ch.~\ref{ch:22}\\
\chtwentytwolean{GriessDictionary.griessIso} & $\C\otimes\chtwentytwoB\cong\Vnat_2$ as algebras with forms & none; Ch.~\ref{ch:21}\\
\chtwentytwolean{ConwayNorton.conwayNorton_gmat} & CN for every element of $\chtwentytwoGmat$, acting on $\Vnat$ & none beyond membership in $\chtwentytwoGmat$; Ch.~\ref{ch:22}\\
\chtwentytwolean{ConwayNorton.nonempty_autVnat_mulEquiv_liftable} & $\Aut(\Vnat)\cong$ liftable subgroup of $\chtwentytwoAutO$ & none; Ch.~\ref{ch:22}\\
\chtwentytwolean{ConwayNorton.M8.gmatCommutantScalar} & $\chtwentytwoGmat$ absolutely irreducible on $\mathbf{e}_1^\perp$ & none; Ch.~\ref{ch:22}\\
\chtwentytwolean{ConwayNorton.monsterEquivAutVnat} & $\chtwentytwoAutO\cong\Aut(\Vnat)$ & \lean{GmatEqAutO}; Ch.~\ref{ch:22}\\
\chtwentytwolean{ConwayNorton.monstrousMoonshine} & monstrous moonshine for AxiomMath's Monster & \textbf{\lean{GmatEqAutO}}; Ch.~\ref{ch:22}\\
\bottomrule
\end{longtable}
\end{center}
All names are in the namespace \lean{VOA.Moonshine}.  Apart from the membership premise of \lean{conwayNorton_gmat}, the only
hypothesis anywhere in the table is \lean{GmatEqAutO}: AxiomMath's statement \lean{gmat_eq_griess} after renaming, with
\lean{phi0} restated definitionally equal; proved in AxiomMath's library on Lean~4.35 (the literature milestones on its path that
were inspected are theorems there), axiom footprint not re-checked, not available on this toolchain.

%======================================================================
\section{Notes and references}\label{ch22:sec:notes}
%======================================================================

The construction of $\Vnat$ and the action of the Monster on it (the centralizer $C=2^{1+24}.\mathrm{Co}_1$ in the orbifold
description, the triality element, and the identification of $\Vnat_2$ with the Griess algebra) are in I.~Frenkel,
J.~Lepowsky and A.~Meurman, \emph{Vertex Operator Algebras and the Monster}, Pure and Applied Mathematics 134, Academic Press,
1988, Chapters~10--12.  The Monster's Griess algebra is R.~L.~Griess Jr., \emph{The friendly giant}, Invent.\ Math.\ 69 (1982)
1--102.  The simplified construction followed by AxiomMath, including the visible group, the element $\xi$ and the triality
element, is J.~H.~Conway, \emph{A simple construction for the Fischer--Griess monster group}, Invent.\ Math.\ 79 (1985)
513--540.  That the automorphism group of the Griess algebra is the Monster is J.~Tits, \emph{On R.~Griess' ``friendly giant''},
Invent.\ Math.\ 78 (1984) 491--499 (not used in this chapter).  S.~P.~Norton's contributions to moonshine and to the structure
of the Griess algebra include J.~H.~Conway and S.~P.~Norton, \emph{Monstrous moonshine}, Bull.\ London Math.\ Soc.\ 11 (1979)
308--339, and S.~P.~Norton, \emph{More on moonshine}, in \emph{Computational Group Theory} (M.~D.~Atkinson, ed.), Academic Press,
1984, 185--193.  The proof of the Conway--Norton conjecture is R.~E.~Borcherds, \emph{Monstrous moonshine and monstrous Lie
superalgebras}, Invent.\ Math.\ 109 (1992) 405--444.  A modern account of the Monster through its Griess algebra and its axes,
including the Majorana/Sakuma approach mentioned among the dead ends, is A.~A.~Ivanov, \emph{The Monster Group and Majorana
Involutions}, Cambridge Tracts in Mathematics 176, Cambridge University Press, 2009.  The MOG and the hexacode are in
R.~T.~Curtis, \emph{A new combinatorial approach to $M_{24}$}, Math.\ Proc.\ Cambridge Philos.\ Soc.\ 79 (1976) 25--42, and
J.~H.~Conway and N.~J.~A.~Sloane, \emph{Sphere Packings, Lattices and Groups}, 3rd ed., Springer, 1999, Chapters~10, 11 and~29.
The Baby Monster recognition through which AxiomMath's library proves $\chtwentytwoGmat=\chtwentytwoAutO$ goes back to J.~S.~Leon and
C.~C.~Sims, \emph{The existence and uniqueness of a simple group generated by $\{3,4\}$-transpositions}, Bull.\ Amer.\ Math.\
Soc.\ 83 (1977) 1039--1040, and its later refinements.  The physics of the orbifold and of the triality is in L.~Dixon,
P.~Ginsparg and J.~Harvey, \emph{Beauty and the beast: superconformal symmetry in a Monster module}, Comm.\ Math.\ Phys.\ 119
(1988) 221--241.  Schur's orthogonality relations and the Schur index are textbook material: J.-P.~Serre, \emph{Linear
Representations of Finite Groups}, Graduate Texts in Mathematics 42, Springer, 1977, \S\S2.2 and~12.2.

The Monster used here is that of the AxiomMath CFSG library: D.~Cureton, V.~Ilin, B.~Mehta, N.~Voss, R.~Cabral, C.~Defant,
E.~Gregory, A.~Hill, L.~Hong, Q.~Hu, D.~Jamner, A.~Jiang, T.~Jordan, A.~Kurghinyan, K.~Lau, J.~Laurent, D.~Liang, J.~Limperg,
S.~Mahns, B.~Nugent, C.~Olson, K.~Ono, E.~Paul, D.~Rajan, K.~Ram, C.~Ray, H.~C.~Siu, A.~A.~Swaminathan, J.~Thorner, A.~Tsaplin and
C.~Wang, \emph{Lean formalization of the classification of finite simple groups}, 2026, \texttt{https://github.com/AxiomMath/CFSG},
commit \texttt{69cb972f587a3ecc035e092253167ff2e303c962} (Lean v4.35.0-rc3), Apache License 2.0; the theorem
\chtwentytwolean{CFSG.Blueprint.NativeSporadic.MonsterModel.gmat_eq_griess} is in \texttt{CFSG/Sporadic/Bridge/MonsterModelsGriess.lean}.
The port of its definitions to Lean~4.30 is by F.~K.~Popov (with Claude).

The development was recorded in the authors' working notes, which are not part of the release; the port audit \leanpath{Notes/\allowbreak monster\_aut/\allowbreak CFSG\_PORT\_FIDELITY.md} is.  The Lean files were drafted largely with the help of automated coding assistants under the authors' direction and verified by targeted builds and axiom audits.

\endgroup
